\documentclass[11pt,a4paper]{article}
\usepackage[utf8]{inputenc}
\usepackage[italian]{babel}
\usepackage[margin=2cm]{geometry}
\usepackage[table]{xcolor}
\usepackage{enumitem}
\usepackage{hyperref}
\usepackage{tcolorbox}
\usepackage{titlesec}
 \usepackage{booktabs}
 \usepackage{tikz}
\usepackage{float}

% Palette di colori professionali
\definecolor{primaryColor}{HTML}{0054A6} % Blu istituzionale TPER
\definecolor{secondaryColor}{HTML}{4A5568}
\definecolor{backgroundColor}{HTML}{F7FAFC}
\definecolor{borderColor}{HTML}{E2E8F0}

\hypersetup{
    colorlinks=true,
    linkcolor=primaryColor,
    urlcolor=primaryColor,
}

\titleformat{\section}
  {\color{primaryColor}\normalfont\Large\bfseries}
  {}{0em}{}
  [\color{borderColor}\titlerule]

\begin{document}

\begin{tcolorbox}[colback=primaryColor,colframe=primaryColor,arc=3mm,halign=center,valign=center,height=2.5cm]
    \color{white}\LARGE\bfseries UX DESIGN PROJECT: HELP PEOPLE HELP PEOPLE \\[0.2cm]
\end{tcolorbox}

\vspace{0.5cm}

\section{Client Card}

\subsection{Name of the Client}
\textbf{TPER S.p.A.} (Trasporto Passeggeri Emilia-Romagna).

\subsection{Description of the Client}
TPER è la principale azienda di trasporto pubblico locale della regione Emilia-Romagna. Nata nel 2012 dalla fusione dei rami d'azienda di ATC (Bologna e Ferrara) e FER (Ferrovie Emilia Romagna), ad oggi conta circa 2500 dipendenti. L'azienda opera sia su gomma, garantendo la mobilità urbana, suburbana ed extraurbana nelle province di Bologna e Ferrara, sia su ferro, operando nel settore del trasporto ferroviario in ambito regionale.
La sede è a Bologna e i principali soci dell'azienda sono: Regione Emilia-Romagna (46,13\%), Comune di Bologna (30,11\%) e Città Metropolitana di Bologna (18,79\%).

\subsection{Name of the Tool}
\textbf{Sito Web Ufficiale di TPER} (\url{https://www.tper.it}).

\subsection{How to access the Tool?}
L'accesso al tool avviene via web ed è strutturato in due livelli principali:
\begin{enumerate}[leftmargin=*]
    \item \textbf{Accesso Libero (Funzioni Informative):} Questa modalità consente la consultazione libera di \textbf{orari}, \textbf{avvisi} sulle modifiche di percorso e il \textbf{calcolo dei viaggi} in tempo reale.
    \item \textbf{Accesso Autenticato (Funzioni Transazionali):} Operazioni come l'\textbf{acquisto o il rinnovo di abbonamenti} personali, la \textbf{gestione della tessera} o il \textbf{caricamento di documenti} per agevolazioni tariffarie necessitano della procedura di autenticazione.
\end{enumerate}

\subsection{Why should the Client want a new Tool?}
\begin{itemize}[leftmargin=*]
    \item Il sito è progettato per utenti già esperti del trasporto pubblico; chi fa parte di categorie fragili  (anziani, stranieri, persone con bassa alfabetizzazione digitale) trova difficoltà nell'orientarsi.
    \item Gli intermediari  (operatori di sportello, assistenti sociali, volontari, parenti) che affiancano utenti fragili non hanno un'area dedicata per mostrare, spiegare e addestrare in modo strutturato.
\end{itemize}
    Impatto attuale: gestione di \textbf{lunghe code} sugli sportelli fisici, rallentamento dell'\textbf{autonomia digitale} degli utenti target e aumento di \textbf{errori procedurali} e \textbf{rimborsi} che Tper deve gestire manualmente.

\subsection{Product Objectives for the Client}

\begin{itemize}[leftmargin=*]
    \item \textbf{Rendere il sito utilizzabile in affiancamento:} Offrire agli intermediari (operatori di sportello, assistenti sociali, volontari, familiari) un'esperienza guidata all'interno del sito stesso, così da poter supportare l'utente fragile direttamente sul canale digitale, senza ricorrere a spiegazioni improvvisate.

    \item \textbf{Ridurre la dipendenza dall'assistenza umana per le operazioni ricorrenti:} Gli utenti fragili siano in grado di completare autonomamente operazioni ad alta frequenza (rinnovo abbonamento, consultazione orari, caricamento documenti) dopo un numero limitato di sessioni guidate.
    
\end{itemize}

\subsection{Business Goals}
\begin{itemize}[leftmargin=*]
\item \textbf{Ottimizzazione dei costi operativi:} Ridurre le spese di gestione del front-office fisico e del supporto telefonico attraverso la migrazione dei flussi assistiti sul canale web.

\item \textbf{Aumento del Customer Lifetime Value}: Incrementare il tasso di rinnovo degli abbonamenti nelle fasce d'età più avanzate, semplificando il mantenimento del titolo di viaggio ed evitando l'abbandono del mezzo pubblico a favore di alternative private o della rinuncia alla mobilità.

\item \textbf{Automazione dei processi di back-office}: Minimizzare l'intervento umano nella gestione dei reclami e delle correzioni dei dati utente.

\end{itemize}

\subsection{Brand Identity}
Il progetto rafforza l'immagine di TPER come azienda vicina al cittadino, capace di coniugare innovazione digitale e inclusione sociale. TPER non appare più come venditore di biglietti ma si prende cura della mobilità delle persone, garantendo a tutti il diritto di muoversi. Tutto ciò si allinea perfettamente con la missione istituzionale e rafforza la legittimità dell'azienda agli occhi degli enti soci.

\subsection{Success Metrics}
\begin{itemize}[leftmargin=*]
    \item \textbf{SUS Score (System Usability Scale):} per misurare il grado di soddisfazione dell'utente prima e dopo il redesign (target $\geq 68/100$).
    \item \textbf{Self-efficacy digitale dell'utente target:} misurata tramite scala Likert prima e dopo la sessione di training, per valutare l'incremento di fiducia e competenza percepita.
    \item \textbf{Riduzione degli errori procedurali} osservati durante test di usabilità condotti su utenti target (metriche: numero di errori, backtrack, richieste di aiuto).
    \item \textbf{Task Completion Time} 
    Il tempo di completamento dei flussi assistiti all'interno della \textit{trainer-mode}.
\end{itemize}

\newpage
\section{User Segmentation card}

\subsection{Segmentation dimensions chosen (demographic \& behavioural \& contextual)}

\begin{tcolorbox}[colback=backgroundColor,colframe=borderColor,arc=2mm]
\begin{itemize}[leftmargin=*,itemsep=4pt]
    \item \textbf{Origine o nazionalità} --- proxy delle barriere linguistiche e della familiarità sistemica.
    \item \textbf{Fascia d'età} --- indicatore del divario generazionale e della familiarità storica con il TPL locale.
    \item \textbf{Competenza digitale} --- livello di confidenza e autonomia nell'uso di interfacce web.
    \item \textbf{Contesto di training} --- l'ambiente fisico e relazionale in cui avviene l'affiancamento.
    \item \textbf{Frequenza d'uso del TPL} --- fattore trainante della motivazione intrinseca all'apprendimento.
\end{itemize}
\end{tcolorbox}

\vspace{0.5cm}

\subsection{Reason for choice, relationship with client's tool}

L'\textbf{origine o nazionalità}: l'utente straniero affronta una doppia barriera, linguistica (comprensione dei testi) e sistemica (mancata conoscenza delle regole del TPL locale o delle agevolazioni ISEE). 
La \textbf{competenza digitale} e la \textbf{fascia d'età} definiscono il livello di supporto necessario: un'anziana italiana può conoscere perfettamente le linee TPER ma bloccarsi davanti a un form di registrazione online.
Il \textbf{contesto di training} determina i vincoli operativi del tool (tempo a disposizione, stress ambientale). Infine, la \textbf{frequenza d'uso del TPL} mappa la reale necessità dell'utente: un utente frequente ha un bisogno stringente di imparare a usare il sito per rinnovare autonomamente l'abbonamento rispetto a un utente occasionale.

\vspace{0.5cm}

\begin{tcolorbox}[colback=backgroundColor,colframe=borderColor,arc=2mm]
\subsection{Segments within each dimension}
\begin{itemize}[leftmargin=*,itemsep=4pt]
\item \textbf{Dimensione \#1 --- Origine o nazionalità}
    Italiana,
    Straniera.
\item \textbf{Dimensione \#2 --- Fascia d'età}
    18--35 anni,
    36--64 anni,
    65+ anni.
\item \textbf{Dimensione \#3 --- Competenza digitale}
    Bassa (nessuna o pochissima autonomia su siti e app),
    Media (uso di app di messaggistica, difficoltà con processi strutturati),
    Alta (navigazione fluida).
\item \textbf{Dimensione \#4 --- Contesto di training}
    Domestico (casa, relazione informale, tempo disteso),
    Istituzionale (sportello TPER, tempo ridotto, forte ansia da coda),
    Associativo (biblioteche comunali, centri di assistenza).
\item \textbf{Dimensione \#5 --- Frequenza d'uso del TPL}
    Rara (meno di una volta al mese),
    Occasionale (1--4 volte al mese),
    Frequente (più di 4 volte al mese, pendolare).
\end{itemize}
\end{tcolorbox}

\vspace{0.3cm}
\begin{flushleft}
\textbf{Total number of segment combinations:} \quad $2 \times 3 \times 3 \times 3 \times 3 = \mathbf{162}$
\end{flushleft}
\vspace{0.3cm}

\subsection{Chosen segmentation combinations}

\begin{tcolorbox}[colback=backgroundColor,colframe=borderColor,arc=2mm]
\begin{itemize}[leftmargin=*,itemsep=6pt]
    \item \phantomsection\label{seg:anziana} \textbf{[S1] L'anziana pendolare} \\ Italiana / 65+ / Competenza bassa / Domestico / Frequente
    \item \phantomsection\label{seg:straniero} \textbf{[S2] Il lavoratore straniero} \\ Straniera / 18--35 / Competenza media / Associativo / Frequente
    \item \phantomsection\label{seg:cittadino} \textbf{[S3] Il cittadino allo sportello} \\ Italiana / 36--64 / Competenza bassa / Istituzionale / Occasionale
\end{itemize}
\end{tcolorbox}

\vspace{0.5cm}

\subsection{Why the chosen combinations are more representative, challenging, expressive than the others}

\begin{enumerate}[leftmargin=*]
    \item \textbf{L'anziana pendolare} rappresenta la sfida della scomposizione cognitiva. Conosce il sistema TPL di Bologna da decenni ma ha bassa competenza digitale. Il contesto domestico permette un training approfondito (es. guidata da un familiare).
    \item \textbf{Il lavoratore straniero} ribalta il problema: ha una discreta competenza digitale, ma si scontra con termini burocratici complessi in italiano. Il contesto associativo, anche qui consente, un training approfondito.
    \item \textbf{Il cittadino in difficoltà allo sportello} incarna lo scenario più stressante. Di mezza età, con bassa competenza digitale, si trova in un contesto istituzionale con tempi ridottissimi (l'operatore allo sportello ha pochi minuti per servirlo e fare formazione). Se il sistema funziona per lui sotto pressione, funziona per chiunque.
\end{enumerate}

\vspace{0.5cm}

\subsection{Expected impact on design}

\begin{itemize}[leftmargin=*,itemsep=6pt]

    \item \textbf{Adattamento al contesto temporale:}
    ogni flusso deve 
    essere completabile in una sessione breve (fast-track per S3) ma anche 
    scomponibile in micro-obiettivi salvabili e ripresi in momenti diversi 
    (apprendimento dilazionato per S1).
    
    \item \textbf{Barriera cognitiva vs.\ barriera linguistica:}
    Il design deve agire su entrambi i livelli: \textbf{semplificazione 
    del flusso} (flussi passo-passo) e \textbf{disambiguazione del 
    linguaggio} (glossario contestuale, icone referenziali, nessuna assunzione 
    implicita sulla conoscenza locale).

    \item \textbf{Tre profili di intermediario, un solo 
    trainer-mode:}
    Il trainer-mode non può 
    assumere un profilo unico di trainer: deve offrire \textbf{strumenti di 
    guida scalabili}, dalla semplice evidenziazione passo-passo (utilizzabile 
    da chiunque senza formazione) fino al pannello di controllo della sessione 
    con checklist di avanzamento, utilizzabile da chi può fare piu sessioni con lo stesso soggetto e ha tempo disponibile.

\end{itemize}
% ================================================================
% MARKET RESEARCH CARD
% ================================================================

\newpage 

\section*{Market Research Card \#1}

\subsection*{Segment of choice}
\begin{tcolorbox}[colback=backgroundColor,colframe=borderColor,arc=2mm]
\textbf{L'anziana pendolare} --- Italiana / 65+ / Competenza digitale bassa / Contesto domestico / Frequente uso del TPL
\end{tcolorbox}

\subsection*{Reason for choice}
\begin{tcolorbox}[colback=backgroundColor,colframe=borderColor,arc=2mm]
Rappresenta il nucleo primario del progetto: è la figura che usa il TPL con continuità ma non riesce ad accedere in autonomia ai servizi digitali. La sua dipendenza dall'intermediario (es. il figlio o il nipote) genera il ciclo che il progetto vuole interrompere. Usiamo dati ISTAT per confermare la scarsa autonomia digitale di questo target su larga scala e \textbf{interviste dirette} per comprendere a fondo le barriere emotive e le dinamiche di affiancamento nel contesto domestico.
\end{tcolorbox}

\subsection*{User research method employed: \textit{Available statistics, One-on-one interviews}}
\begin{tcolorbox}[colback=backgroundColor,colframe=borderColor,arc=2mm]
\textbf{Sources (Available statistics):}
\begin{itemize}[leftmargin=*]
  \item ISTAT, \textit{Cittadini e ICT --- Anno 2025}, comunicato stampa.\\
  \href{https://www.istat.it/comunicato-stampa/cittadini-e-ict-anno-2025/}{Link al comunicato stampa ISTAT}
  \item ISTAT, \textit{Cittadini e ICT --- Anno 2024}, 28 aprile 2025.\\
  \href{https://www.istat.it/wp-content/uploads/2025/04/REPORT_CITTADINI-E-ICT_2024.pdf}{Link al PDF del report}
\end{itemize}

\vspace{0.3cm}
\textbf{Interviews:} 3 interviste ad utenti in target. Durata media 15 minuti ciascuna.
\begin{itemize}[leftmargin=*]
  \item Livello di autonomia attuale sui servizi digitali.
  \item Dinamica relazionale durante la richiesta di aiuto.
  \item Barriere emotive (paura di sbagliare, frustrazione, ansia tecnologica).
  \item Disponibilità nell'imparare a compiere le operazioni autonomamente.
\end{itemize}
\end{tcolorbox}

\subsection*{Overall results of the analysis}
\begin{tcolorbox}[colback=backgroundColor,colframe=borderColor,arc=2mm]
\textbf{Dati quantitativi (ISTAT):} Mostrano che il \textbf{72,5\% dei 65-74enni} utilizza Internet, quota che scende al \textbf{35,7\% tra gli over 75}. Le competenze digitali almeno di base si attestano al \textbf{27,4\%} nella fascia 65-74 anni. Il divario di genere è marcato: tra le ultra-75enni solo il 30,5\% usa internet, contro il 43,0\% degli uomini coetanei.

\vspace{0.2cm}
\textbf{Dati qualitativi (Interviste):} Dalle interviste emerge un forte senso di inadeguatezza tecnologica. Quando chiedono aiuto a un familiare, la dinamica tipica è la \textbf{delega totale}: il familiare "fa al posto loro" per mancanza di tempo, senza spiegare i passaggi. Le anziane intervistate provano frustrazione e temono di "fare danni irreversibili". Tuttavia alcune si dichiarano disposte a imparare, a patto di essere guidate con calma, passo dopo passo.
\end{tcolorbox}

\subsection*{Expected impact on design}
\begin{tcolorbox}[colback=backgroundColor,colframe=borderColor,arc=2mm]

\begin{itemize}[leftmargin=*,itemsep=6pt]

    \item \textbf{La trainer-mode deve ristrutturare la dinamica relazionale} Il sistema deve quindi 
    fornire all'intermediario istruzioni esplicite su \textit{come} guidare: 
    suggerimenti contestuali del tipo ``lascia che sia lei a premere'' o 
    ``aspetta che abbia letto prima di procedere'' ridefiniscono il ruolo 
    del trainer da esecutore a facilitatore.

    \item \textbf{La paura dell'errore irreversibile} Ogni step deve comunicare 
    esplicitamente la reversibilità dell'azione compiuta come meccanismo strutturale del flusso e 
    prevedere un linguaggio rassicurante che normalizzi l'errore.

    \item \textbf{Il dato ISTAT impone di non assumere nessuna competenza 
    digitale di base.} Ogni elemento 
    dell'interfaccia deve essere autoesplicativo nella sua funzione, 
    indipendentemente dall'esperienza pregressa dell'utente.

\end{itemize}
\end{tcolorbox}

\newpage 

\newpage
\section*{Market Research Card \#2}

\subsection*{Segment of choice}
\begin{tcolorbox}[colback=backgroundColor,colframe=borderColor,arc=2mm]
\textbf{Il cittadino in difficoltà allo sportello} --- Italiana / 36--64 anni / Competenza digitale bassa / Contesto istituzionale / Uso del TPL occasionale
\end{tcolorbox}

\subsection*{Reason for choice}
\begin{tcolorbox}[colback=backgroundColor,colframe=borderColor,arc=2mm]
Questo segmento rappresenta il punto di contatto più critico e diretto tra l'utente fragile e il sistema TPER. L'osservazione sul campo allo sportello fisico permette di rilevare contemporaneamente i comportamenti reali dell'utente e i vincoli operativi del contesto istituzionale.
\end{tcolorbox}

\subsection*{User research method employed: \textit{Field observation}}
\begin{tcolorbox}[colback=backgroundColor,colframe=borderColor,arc=2mm]
\textbf{Location of observation:} Sportello TPER Bologna --- via Guglielmo Marconi 4 \\[0.4cm]
\textit{L'osservazione ha rilevato congiuntamente: tipologia di richiesta, livello di passività dell'utente, comportamenti di disagio legati alla documentazione, tempi effettivi di gestione delle richieste e il grado di coinvolgimento del sito tper.it durante l'interazione.}
\end{tcolorbox}

\subsection*{Overall results of the analysis}
\begin{tcolorbox}[colback=backgroundColor,colframe=borderColor,arc=2mm]
Le richieste si concentrano su: rinnovo abbonamenti, caricamento documenti e info sui percorsi. La dinamica osservata è di \textbf{totale passività}: l'utente attende mentre l'operatore risolve la pratica sul proprio terminale proprietario (3--5 minuti per operazioni semplici, oltre 10 minuti per pratiche documentali). Il sito web pubblico viene sistematicamente bypassato e mai mostrato al cittadino. La costante pressione della coda impedisce di fare formazione, mentre l'utente, non venendo coinvolto, non impara nulla. Inoltre, il forte disagio emerge quando mancano i documenti necessari, costringendo l'utente a tornare più volte.
\end{tcolorbox}

\subsection*{Expected impact on design}
\begin{tcolorbox}[colback=backgroundColor,colframe=borderColor,arc=2mm]

\begin{itemize}[leftmargin=*,itemsep=6pt]

     \item \textbf{Percorsi tematici e feedback visivi:} Prevedere flussi (es. ``Rinnovo abbonamento'') selezionabili velocemente dall'operatore, affiancati da un indicatore di avanzamento sempre visibile per mitigare l'ansia da coda dell'utente.

    \item \textbf{La passività osservata allo sportello non ha origine 
    emotiva, ma strutturale:} A differenza del contesto domestico, 
    qui l'utente non delega per insicurezza: è il sistema stesso --- la 
    coda, il terminale rivolto verso l'operatore, i tempi compressi --- 
    a escluderlo fisicamente dal processo. Il design non può affidarsi 
    a suggerimenti comportamentali per l'intermediario: deve rendere 
    il coinvolgimento dell'utente la strada più rapida, non quella più 
    lenta.

     \item \textbf{Checklist preventiva:} Inserire un elenco visivo chiaro e accessibile sul sito dei documenti necessari prima di avviare una pratica, per ridurre i ritorni a vuoto.

\end{itemize}
\end{tcolorbox}

% ----------------------------------------------------------------
% CARD 5 — L'immigrato lavoratore
% ----------------------------------------------------------------
\newpage
\section*{Market Research Card \#3}
\subsection*{Segment of choice}
\begin{tcolorbox}[colback=backgroundColor,colframe=borderColor,arc=2mm]
\textbf{Il lavoratore straniero} --- Straniera / 18--35 anni / Competenza digitale media / Contesto associativo / Uso del TPL frequente
\end{tcolorbox}

\subsection*{Reason for choice}
\begin{tcolorbox}[colback=backgroundColor,colframe=borderColor,arc=2mm]
Introduce una doppia barriera, linguistica e sistemica, che nessun'altra combinazione presenta. Ha dimestichezza con smartphone e app, ma non conosce le logiche del TPL locale e incontra difficoltà con il linguaggio tecnico-burocratico del sito in italiano. Le statistiche disponibili consentono di quantificare la dimensione del fenomeno a livello nazionale e regionale, validando la scelta del segmento.
\end{tcolorbox}

\subsection*{User research method employed: \textit{Available statistics}}
\begin{tcolorbox}[colback=backgroundColor,colframe=borderColor,arc=2mm]
\textbf{Sources:}
\begin{itemize}[leftmargin=*]
  \item ISTAT, \textit{Indicatori demografici --- Anno 2024},
  comunicato stampa, 31 marzo 2025.\\
  \href{https://www.istat.it/comunicato-stampa/indicatori-demografici-anno-2024/}{Link al comunicato ISTAT}

  \item ISTAT, \textit{Censimento e dinamica della popolazione --- Anno 2024}, dicembre 2024.\\
  \href{https://www.istat.it/wp-content/uploads/2025/12/Censimento-e-dinamica-della-popolazione-Anno-2024.pdf}{Link al PDF del censimento}
\end{itemize}
\end{tcolorbox}

\subsection*{Overall results of the analysis}
\begin{tcolorbox}[colback=backgroundColor,colframe=borderColor,arc=2mm]
Al 1° gennaio 2025 risiedono in Italia \textbf{5.422.000 cittadini
stranieri}, pari al \textbf{9,2\%} della popolazione totale, con una
crescita costante anno su anno.
 L'\textbf{età media della popolazione straniera} è
di \textbf{37,0 anni}, nettamente inferiore a quella dei cittadini
italiani (47,8 anni), confermando che si tratta di una popolazione
prevalentemente giovane e in età lavorativa. La
concentrazione geografica è marcata al Nord, dove risiede il
\textbf{58,3\%} degli stranieri residenti, con una forte presenza
anche in Emilia-Romagna.
\end{tcolorbox}

\subsection*{Expected impact on design}
\begin{tcolorbox}[colback=backgroundColor,colframe=borderColor,arc=2mm]
 Sul piano del linguaggio: \textbf{frasi brevi, verbi all'infinito, nessun termine tecnico-burocratico} senza spiegazione. Sul piano visivo: \textbf{icone esplicite e universali} affiancate sempre da etichette testuali, nessuna assunzione implicita sulla conoscenza della rete TPL. Essendo spesso aiutati in ambito associativo puo essere utile portare a casa un riepilogo stampabile dei passi completati. Queste scelte, progettate per il caso più estremo, migliorano l'esperienza per tutti gli utenti fragili.
\end{tcolorbox}

%--------------sezione5a-------
\newpage
\section*{Assessment of Existing Resource Card}

\textbf{This is the client's tool}\\
\textit{(i.e., the one that needs to be redesigned)}

\bigskip

\subsection*{Name of Tool}

\noindent Ecosistema digitale TPER: Sito Web Ufficiale (\url{https://www.tper.it}).

\subsection*{Platforms for Tool (devices, operating systems, etc.)}

\noindent Il sito web è accessibile via browser da qualsiasi dispositivo (desktop, laptop, smartphone, tablet) su tutti i principali sistemi operativi (Windows, macOS, iOS, Android). Le funzioni transazionali del sito richiedono autenticazione tramite SPID o credenziali proprietarie. La pianificazione del viaggio è disponibile anche senza registrazione.

\subsection*{Main Purpose of Tool, Main Operations Allowed, Perceived Level of Adequacy to Professional Standards}

\noindent Il sito ha duplice funzione: informativa e transazionale. Le operazioni principali consentite sono: consultazione di orari e percorsi in tempo reale, calcolo del viaggio, acquisto e rinnovo di abbonamenti personali, gestione della tessera fisica e digitale (\textit{Mi Muovo} / \textit{Io Viaggio}), caricamento di documenti per agevolazioni tariffarie.

\noindent Il livello di adeguatezza agli standard professionali è \textit{parziale}: le funzioni informative e di pianificazione sono discretamente implementate, mentre i flussi transazionali presentano criticità strutturali in termini di chiarezza delle etichette, gestione degli errori e supporto agli utenti non esperti, in particolare anziani over 70 con bassa competenza digitale.

\subsection*{Main Strengths of Tool, Reason for Choice}

\begin{itemize}
    \item Copertura funzionale ampia: il sito centralizza tutte le funzioni TPL (orari, acquisti, abbonamenti, avvisi di servizio) in un unico punto di accesso istituzionale.
    \item Presenza di pianificazione del viaggio e consultazione dei servizi in tempo reale.
    \item Accessibilità base garantita: il sito è consultabile senza registrazione per le funzioni informative.
    \item Capillarità del sistema: TPER è l'operatore di riferimento per un ampio bacino territoriale in Emilia-Romagna.
    \item Identità di marca già riconoscibile e consolidata presso l’utenza locale.
\end{itemize}

\subsection*{Main Weaknesses of Tool}

\begin{itemize}
    \item Architettura dell'informazione complessa: il menu mescola voci eterogenee senza gerarchia visiva, rendendo difficile per l'utente anziano identificare dove acquistare un biglietto o rinnovare un abbonamento.
    \item Terminologia tecnica non semplificata: etichette come ``tessera Mi Muovo'', ``wallet'' e ``SPID'' sono usate senza spiegazione contestuale.
    \item Autenticazione SPID come barriera cognitiva: la procedura richiede competenze digitali avanzate incompatibili con il profilo dell'utente anziano target.
    \item Nessuna modalità di affiancamento caregiver: il sito non prevede vista condivisa, sessione co-navigabile o funzioni trainer.
    \item Assenza di supporto guidato contestuale: nessun tooltip attivabile, nessun flusso passo-passo, nessuna guida vocale o tattile.
    \item Gestione degli errori insufficiente: il sito non guida l'utente nel recupero da errori comuni come pagamento fallito o abbonamento scaduto.
    \item Nessun monitoraggio dei progressi: l'intermediario non dispone di strumenti per tracciare le competenze acquisite dall'utente e pianificare training incrementali.
\end{itemize}

\subsection*{First Impressions — Esplorazione Iniziale}

\noindent La prima esplorazione del sito \texttt{tper.it} produce un'impressione di \textbf{complessità non gestita}. La homepage presenta simultaneamente news aziendali, avvisi di servizio, widget di calcolo percorso e link alle aree transazionali, senza una chiara gerarchia visiva che orienti l'utente verso il task primario. Un osservatore esterno impiegherebbe diversi secondi a capire dove acquistare un biglietto o rinnovare un abbonamento. La presenza di SPID come modalità di autenticazione rappresenta una barriera visibile già nella prima interazione con le funzioni transazionali.

\noindent\textbf{Problemi identificati in fase esplorativa (pre-analisi formale):}

\begin{itemize}
    \item Assenza di un percorso guidato per il primo accesso.
    \item Mancanza di una CTA primaria visibile sulla homepage.
    \item Il sito non comunica chiaramente il percorso verso l’acquisto.
    \item Il sito non segnala preventivamente che alcune funzioni richiedono SPID.
    \item Nessuna funzione di help contestuale.
\end{itemize}

\subsection*{Direct Analysis: System vs.\ Guidelines}

\noindent\textit{Analisi condotta secondo le linee guida UserFocus (\url{http://www.userfocus.co.uk/resources/guidelines.html}). Per ogni issue viene indicata la linea guida violata, la localizzazione nel sistema}

\bigskip

\noindent Le linee guida di riferimento adottate sono:

\begin{tabular}{ll}
\hline
\textbf{\#} & \textbf{Linea Guida} \\
\hline
G1 & Home page — chiarezza degli obiettivi del sito \\
G2 & Task orientation — supporto al completamento dei task utente \\
G3 & Navigation and IA — chiarezza della struttura di navigazione \\
G4 & Forms and data entry — usabilità dei moduli \\
G5 & Trust and credibility — fiducia nel sistema \\
G6 & Writing and content quality — qualità del linguaggio \\
G7 & Page layout and visual design — gerarchia visiva \\
G8 & Search — ricercabilità dei contenuti \\
G9 & Help, feedback and error tolerance — gestione errori e supporto \\
G10 & Accessibility — accessibilità cognitiva e motoria \\
\hline
\end{tabular}

\bigskip

\noindent\textbf{Issue D1 — Violazione G1, G7}\phantomsection\label{iss:D1}\\
\textit{Localizzazione: homepage tper.it}\\

\noindent La homepage non comunica immediatamente lo scopo primario del sito. News aziendali, avvisi di modifica percorso e widget transazionali occupano spazio equivalente, senza gerarchia. L'utente anziano non riesce a identificare in pochi secondi se può acquistare un biglietto direttamente dalla homepage. La mancanza di una CTA primaria visibile costringe l'esplorazione del menu, aumentando il carico cognitivo fin dalle prime fasi dell'interazione.

\bigskip

\noindent\textbf{Issue D2 — Violazione G3}\phantomsection\label{iss:D2}\\
\textit{Localizzazione: menu principale tper.it, struttura di navigazione generale}\\

\noindent L'architettura dell'informazione non distingue chiaramente tra sezioni informative (orari, avvisi) e sezioni transazionali (acquisto, rinnovo). Le etichette usano terminologia interna aziendale che non corrisponde al vocabolario dell'utente target. Il percorso verso l'acquisto di un abbonamento richiede più di tre click e non è lineare.

\bigskip

\noindent\textbf{Issue D3 — Violazione G6, G10}\phantomsection\label{iss:D3}\\
\textit{Localizzazione: sito — etichette, bottoni, messaggi di sistema}\\

\noindent Il linguaggio tecnico utilizzato sul sito non viene mai semplificato né spiegato contestualmente. Nessuna area offre glossari, tooltip o micro-copy esplicativi in prossimità dei termini tecnici. Per un utente anziano con bassa literacy digitale queste parole sono prive di significato operativo e producono blocco cognitivo.

\bigskip

\noindent\textbf{Issue D4 — Violazione G9}\phantomsection\label{iss:D4}\\
\textit{Localizzazione: flusso di acquisto e validazione biglietto}\\

\noindent Il sistema non guida l'utente nel recupero dagli errori più comuni: pagamento fallito, abbonamento scaduto non rinnovato. In questi casi il sito non mostra messaggi di errore chiari con indicazioni su come procedere, generando chiamate al call center, code allo sportello e richieste di rimborso.

\bigskip

\noindent\textbf{Issue D5 — Violazione G4}\phantomsection\label{iss:D5}\\
\textit{Localizzazione: flusso di autenticazione sito (SPID)}\\

\noindent La procedura di login tramite SPID — obbligatoria per le funzioni transazionali del sito — non è accompagnata da spiegazioni su cosa sia SPID, come ottenerlo o cosa fare in sua assenza. L'utente che non ha SPID si trova bloccato senza alternative chiare. La complessità del flusso è incompatibile con le competenze digitali del target anziano.

\subsection*{Reverse Analysis: Guidelines vs.\ System}

\noindent\textit{Analisi sistematica guideline per guideline, cercando nel sistema le funzioni che violano ciascuna linea guida. I problemi già identificati nella Direct Analysis vengono omessi; si riportano solo le nuove violazioni emerse.}

\bigskip

\noindent\textbf{Issue R1 — Linea Guida G5 (Trust and Credibility)}\phantomsection\label{iss:R1}\\
\textit{Localizzazione: sito tper.it — area acquisto e pagamento online}\\

\noindent Il sito non fornisce segnali di fiducia adeguati nelle schermate di acquisto: assenza di icone di sicurezza del pagamento, mancanza di conferme intermedie durante il flusso e nessuna sintesi dell'ordine prima del pagamento definitivo. Per un utente anziano che già diffida dei pagamenti online, questa assenza di rassicurazioni aumenta l'ansia e la probabilità di abbandono del flusso.

\bigskip

\noindent\textbf{Issue R2 — Linea Guida G8 (Search)}\phantomsection\label{iss:R2}\\
\textit{Localizzazione: sito tper.it — funzione di ricerca interna}\\

\noindent La funzione di ricerca interna non è prominente, non suggerisce query correlate né autocompleta. Una ricerca per \textit{``abbonamento anziani''} non restituisce risultati ordinati per rilevanza. L'utente anziano che usa la ricerca come scorciatoia non ottiene risposte utili ed è costretto a navigare manualmente la struttura del sito.

\bigskip

\noindent\textbf{Issue R3 — Linea Guida G10 (Accessibility)}\phantomsection\label{iss:R3}\\
\textit{Localizzazione: sito tper.it — testi, contrasto, aree cliccabili}\\

\noindent Alcune sezioni presentano testi di dimensione ridotta, contrasto testo/sfondo insufficiente per utenti con riduzione visiva lieve e aree cliccabili troppo piccole per l'uso su smartphone. L'assenza di modalità ad alto contrasto o di opzioni di accessibilità visiva penalizza direttamente il target anziano.

\bigskip

\noindent\textbf{Issue R4 — Linea Guida G2 (Task Orientation)}\phantomsection\label{iss:R4}\\
\textit{Localizzazione: sito tper.it — distinzione tra pianificazione e acquisto}\\

\noindent Il sito integra nella stessa interfaccia la pianificazione del viaggio e l'acquisto del titolo di viaggio senza che questa distinzione sia resa visivamente chiara. Un utente alla prima esperienza può pianificare un percorso senza rendersi conto di non han ancora acquistato il titolo, rischiando di salire sull'autobus senza biglietto valido.

\bigskip

\noindent\textbf{Issue R5 — Linea Guida G9 (Help, Feedback, Error Tolerance)}\phantomsection\label{iss:R5}\\
\textit{Localizzazione: sito tper.it — assenza modalità trainer}\\

\noindent Il sito non offre alcuna forma di supporto contestuale attivabile (tooltip, guide passo-passo, assistente virtuale, FAQ integrate nei flussi). L'utente fragile che si blocca non ha strumenti per proseguire autonomamente. In modo ancora più critico, non esiste nessuna modalità trainer che permetta a un intermediario di guidare l'utente fragile in modo strutturato e ripetibile, che è esattamente il bisogno principale identificato nel progetto HPHP.

\newpage

%competitor atac%
\section*{Assessment of Existing Resource Card}

\textbf{Questo è il tool concorrente}\\
\textit{(ossia uno dei riferimenti esterni usati per il confronto)}

\subsection*{Nome del tool}

\noindent Sito web ATAC Roma (\url{https://www.atac.roma.it}).

\subsection*{Piattaforme del tool (dispositivi, sistemi operativi, ecc.)}

\noindent Il sito è accessibile da browser su desktop, laptop, tablet e smartphone, nei principali sistemi operativi. Offre servizi online relativi a biglietti, abbonamenti, assistenza e informazioni di viaggio. Alcune funzioni richiedono l’accesso all’area MyAtac, mentre altre sono pubbliche. Il sito rappresenta quindi il principale punto di ingresso istituzionale all’ecosistema digitale ATAC.

\subsection*{Scopo principale del tool, operazioni consentite, livello percepito di adeguatezza agli standard professionali}

\noindent Il sito ha sia una funzione informativa sia una funzione transazionale. Consente di acquistare e rinnovare abbonamenti, gestire alcune operazioni sulle tessere, richiedere o ricaricare titoli agevolati, accedere all’assistenza clienti e consultare informazioni di servizio in tempo reale. Spiega anche dove acquistare i biglietti fisici e come usare i titoli digitali, comprese le regole di validazione del biglietto da smartphone. Il tool è adeguato sul piano professionale per ampiezza dei servizi, ma solo parzialmente sul piano della chiarezza per gli utenti fragili. 

\subsection*{Punti di forza principali del tool, motivazione della scelta}

\begin{itemize}
    \item Copertura ampia dei bisogni dell’utente: il sito riunisce biglietti, abbonamenti, assistenza e informazioni in un unico ambiente istituzionale.
    \item Forte aderenza ai compiti reali di viaggio: l’utente può acquistare, rinnovare e validare titoli online.
    \item Le informazioni di customer service sono disponibili direttamente sul sito, con recapiti e canali di supporto.
    \item Il sito supporta sia titoli standard sia titoli agevolati, risultando utile per una platea molto ampia.
    \item La presenza di istruzioni esplicite sull’uso del titolo digitale mostra un tentativo concreto di supporto operativo.
\end{itemize}

\subsection*{Punti di debolezza principali del tool}

\begin{itemize}
    \item Il sito è denso e molto testuale, quindi aumenta il carico cognitivo dei nuovi utenti.
    \item I servizi principali sono distribuiti su più pagine e aree, e questo indebolisce l’orientamento al compito.
    \item La terminologia presuppone una certa familiarità con abbonamenti, validazione e procedure digitali.
    \item La distinzione tra pagine informative e pagine transazionali non è sempre immediata.
    \item L’assistenza è disponibile, ma non sempre integrata nel momento esatto in cui l’utente ha bisogno di aiuto.
\end{itemize}

\subsection*{Prime impressioni — esplorazione iniziale}

\noindent La prima impressione del sito ATAC è quella di un sistema maturo ma impegnativo. Offre molti servizi utili, ma li presenta attraverso una quantità rilevante di informazioni e con diversi percorsi che l’utente deve interpretare. Per un viaggiatore esperto questo può essere accettabile, ma per un anziano con bassa fiducia digitale può risultare eccessivo. Il sito appare autorevole e completo, ma non pensato in modo specifico per rassicurare l’utente insicuro nel momento dell’azione.

\noindent\textbf{Problemi identificati nella fase esplorativa:}

\begin{itemize}
    \item L’obiettivo principale non è sempre subito visibile in homepage.
    \item L’utente deve già sapere dove andare prima di iniziare il compito.
    \item Biglietteria, assistenza e contenuti informativi non sono sempre ben separati.
    \item Alcune procedure vengono spiegate solo dopo l’ingresso nella sezione giusta.
    \item Il sito informa, ma non sempre riduce l’esitazione.
\end{itemize}

\subsection*{Analisi diretta: sistema vs.\ linee guida}

\noindent\textit{Analisi condotta secondo le linee guida UserFocus (\url{http://www.userfocus.co.uk/resources/guidelines.html}). Per ogni problema viene indicata la linea guida violata, la localizzazione nel sistema e la classificazione della gravità.}

\bigskip

\noindent Le linee guida di riferimento adottate sono:

\begin{tabular}{ll}
\hline
\textbf{\#} & \textbf{Linea guida} \\
\hline
G1 & Homepage — chiarezza degli obiettivi del sito \\
G2 & Orientamento al compito — supporto al completamento dei task \\
G3 & Navigazione e architettura informativa — chiarezza della struttura \\
G4 & Moduli e inserimento dati — usabilità dei form \\
G5 & Fiducia e credibilità — affidabilità percepita \\
G6 & Scrittura e qualità dei contenuti — qualità del linguaggio \\
G7 & Layout e design visivo — gerarchia visiva \\
G8 & Ricerca — reperibilità dei contenuti \\
G9 & Aiuto, feedback e tolleranza agli errori — supporto e recupero \\
G10 & Accessibilità — accessibilità cognitiva e motoria \\
\hline
\end{tabular}

\bigskip

\noindent\textbf{Issue D1 — Violazione G1, G7}\phantomsection\label{atac:D1}\\
\textit{Localizzazione: homepage del sito ATAC}\\

\noindent La homepage non rende immediatamente evidente quale sia il percorso principale per chi vuole compiere un’azione concreta, come acquistare o rinnovare un titolo di viaggio. Le informazioni istituzionali, i contenuti di servizio e le voci transazionali convivono senza una gerarchia abbastanza chiara. Questo costringe l’utente a esplorare e interpretare il sito prima ancora di iniziare il task. Per un utente anziano o insicuro, questa soglia d’ingresso è troppo alta.

\bigskip

\noindent\textbf{Issue D2 — Violazione G3}\phantomsection\label{atac:D2}\\
\textit{Localizzazione: struttura generale del sito e delle sezioni di servizio}\\

\noindent L’architettura dell’informazione distribuisce i contenuti in sezioni diverse che non sempre presentano una sequenza d’uso evidente. L’utente deve decidere da solo dove iniziare, come passare dall’informazione all’azione e quale sezione sia quella giusta per il proprio problema. Questa struttura è accettabile per utenti esperti, ma non per chi ha bisogno di un percorso lineare e ripetibile. La navigazione resulta quindi formalmente corretta, ma cognitivamente pesante.

\bigskip

\noindent\textbf{Issue D3 — Violazione G6, G10}\phantomsection\label{atac:D3}\\
\textit{Localizzazione: testi informativi, istruzioni su titoli e servizi online}\\

\noindent Il linguaggio del sito è corretto ma spesso presuppone familiarità con concetti come attivazione, validazione, MyAtac, ticket digitale e titoli agevolati. Le spiegazioni sono presenti, ma non sempre abbastanza contestuali per un utente che entra nel sistema con scarsa fiducia. Di conseguenza il sito informa, ma non sempre accompagna. L’accessibilità cognitiva resta quindi solo parziale, soprattutto per utenti over 70.

\bigskip

\noindent\textbf{Issue D4 — Violazione G9}\phantomsection\label{atac:D4}\\
\textit{Localizzazione: pagine di aiuto e gestione delle eccezioni}\\

\noindent Il sito fornisce supporto e indicazioni, ma non costruisce una vera esperienza di recupero nel momento in cui qualcosa va storto. Se l’utente incontra un problema durante l’acquisto o la gestione del titolo, deve interpretare da solo le istruzioni disponibili. Per chi è già insicuro, questo può essere sufficiente a interrompere il processo. L’errore non viene accompagnato in modo abbastanza rassicurante.

\bigskip

\noindent\textbf{Issue D5 — Violazione G4}\phantomsection\label{atac:D5}\\
\textit{Localizzazione: aree che richiedono accesso MyAtac e moduli collegati}\\

\noindent Alcuni servizi richiedono l’accesso all’area MyAtac o l’uso di moduli e passaggi preliminari che aumentano la soglia di ingresso. La registrazione e la gestione del profilo sono funzionali, ma aggiungono un carico mentale che può bloccare utenti poco confidenti. Il problema non è la presenza dell’account, ma il fatto che il sito non rende abbastanza morbido il passaggio iniziale. Per il target HPHP, questa è una criticità rilevante.

\subsection*{Analisi inversa: linee guida vs.\ sistema}

\noindent\textit{Analisi sistematica linea guida per linea guida, cercando nel sistema le funzioni che la violano. I problemi già individuati nella analisi diretta vengono omessi; si riportano solo le nuove violazioni emerse.}

\bigskip

\noindent\textbf{Issue R1 — Violazione G5 (Fiducia e credibilità)}\phantomsection\label{atac:R1}\\
\textit{Localizzazione: sito ATAC nel suo complesso}\\

\noindent Il sito è credibile dal punto di vista istituzionale, ma la quantità di percorsi e aree di servizio può indebolire la fiducia operativa del nuovo utente. Quando non è subito chiaro quale sia il canale giusto per un dato compito, l’utente può percepire incertezza anche in un sistema affidabile. Il problema non riguarda la reputazione del servizio, ma la sua leggibilità d’uso. La fiducia è forte sul brand, meno forte sul task.

\bigskip

\noindent\textbf{Issue R2 — Violazione G8 (Ricerca)}\phantomsection\label{atac:R2}\\
\textit{Localizzazione: pagine informative e sezioni di servizio del sito}\\

\noindent La ricerca non appare come un vero strumento capace di compensare la complessità dell’architettura. Un utente che non sa dove andare deve comunque conoscere abbastanza bene il dominio per formulare una richiesta efficace. Questo limita l’utilità della ricerca como scorciatoia cognitiva. Per utenti fragili, la ricerca non sostituisce una struttura più guidata.

\bigskip

\noindent\textbf{Issue R3 — Violazione G10}\phantomsection\label{atac:R3}\\
\textit{Localizzazione: pagine di servizio, istruzioni e percorsi account-based}\\

\noindent L’accessibilità è presente sul piano funzionale, ma il sistema resta impegnativo sul piano mentale. L’utente deve distinguere tra diversi tipi di contenuto, comprendere il significato di termini tecnici e ricordare dove si trova nel percorso. Questo riduce l’adeguatezza del sito per utenti con bassa fiducia digitale. Il problema è dunque cognitivo prima ancora che visivo o motorio.

\bigskip

\noindent\textbf{Issue R4 — Violazione G2}\phantomsection\label{atac:R4}\\
\textit{Localizzazione: percorso tra informazione, acquisto e gestione del titolo}\\

\noindent Il sito supporta il task, ma non lo rende immediatamente leggibile come sequenza di azioni. Consultare, acquistare, attivare e gestire un titolo sono operazioni collegate, ma il sistema si aspetta che l’utente ne comprenda già il rapporto. Questo è poco adatto a un pubblico che ha bisogno di un modello mentale semplice e rassicurante. La piattaforma facilita il compito solo dopo che l’utente ha già superato la parte più difficile: capire il compito.

\bigskip

\noindent\textbf{Issue R5 — Violazione G9}\phantomsection\label{atac:R5}\\
\textit{Localizzazione: supporto diffuso nel sito}\\

\noindent Il sito offre aiuto, ma non una vera esperienza di accompagnamento nel momento dell’azione. In caso di dubbio o errore, l’utente deve spesso trasformare da solo le istruzioni in comportamento concreto. Per il target HPHP questo non basta, perché il problema principale non è la mancanza di informazioni ma la mancanza di fiducia nel momento di usarle. L’assenza di un supporto più guidato rappresenta quindi la criticità più grave.

\newpage

%fine competitor

% ================================================================
% PRELIMINARY USER TESTING CARD — Client's tool
% ================================================================

\newpage

\section{Preliminary User testing Card, 
This test relates to the client's tool
(i.e., the one that needs to be redesigned)}

\vspace{0.2cm}

\subsection*{Name of tool}
\begin{tcolorbox}[colback=backgroundColor,colframe=borderColor,arc=2mm]
Sito web ufficiale TPER --- \url{https://www.tper.it}
\end{tcolorbox}

\vspace{0.2cm}

\begin{center}
  {\color{primaryColor}\large\bfseries Testing protocol}
  \vspace{0.1cm}
  \color{borderColor}\hrule
\end{center}

\vspace{0.2cm}

\subsection*{Testing session setup}
\begin{tcolorbox}[colback=backgroundColor,colframe=borderColor,arc=2mm]
Le sessioni di test sono state condotte \textbf{in presenza},
presso il domicilio del soggetto. Il dispositivo utilizzato è
stato il laptop del ricercatore, con sistema operativo
\textbf{Windows 11}, browser \textbf{Google Chrome} aggiornato
all'ultima versione disponibile, risoluzione schermo 1920×1080.
Il ricercatore era presente nella stanza e non ha interagito
con il soggetto durante lo svolgimento dei task ma si è limitato a spiegare lo scopo del sito e dell'attività, coerentemente
con la metodologia \textit{bottom line data}.
\end{tcolorbox}

\vspace{0.2cm}

\subsection*{Testing method}
\textbf{Discount usability testing} (2--5 persone)

\vspace{0.2cm}

\subsection*{List of tasks to test}
\begin{tcolorbox}[colback=backgroundColor,colframe=borderColor,arc=2mm]
\begin{enumerate}[leftmargin=*,itemsep=4pt]
  \item \textbf{Ricerca e consultazione orari:} L'utente deve
  trovare fermate e orari di una linea specifica (Linea 39: Circolare periferica sinistra)
  \item \textbf{Pianificazione dell'itinerario:} L'utente deve utilizzare le funzionalità offerte dall'applicazione per calcolare un percorso impostando origine (Piazza Maggiore) e destinazione (Giardini Margherita).
  \item \textbf{Acquisto abbonamento mensile:} L'utente
  deve navigare l'architettura informativa esistente per
  trovare l'accesso alla pagina di acquisto di un abbonamento
  mensile impersonale.
  \item \textbf{Identificazione agevolazioni e requisiti:}
  L'utente deve cercare i requisiti  e i documenti necessari per ottenere riduzioni tariffarie per l'acquisto di un
  abbonamento annuale in relazione alla propria condizione (ISEE ,anziani ecc..)
  \item \textbf{Rinnovo dell'abbonamento:} L'utente
  deve individuare l'area dedicata al rinnovo del proprio titolo di viaggio.
\end{enumerate}
\end{tcolorbox}

\vspace{0.2cm}

\subsection*{Testing methodology}
\textbf{Bottom line data} --- il soggetto
svolge i task da solo, senza interazione con il ricercatore.
Questa modalità è stata scelta per raccogliere dati
quantitativi puliti su completamento, tempo ed errori,
non influenzati dalla verbalizzazione o dal dialogo con
il facilitatore.

\vspace{0.2cm}

\subsection*{Metrics of tests \normalfont\small\itshape}
\begin{itemize}[leftmargin=*,itemsep=3pt]
  \item \textbf{Success} --- percentuale di task portati
  a termine con successo.
  \item \textbf{Time} --- tempo impiegato per completare
  ciascun task (in secondi).
  \item \textbf{Errors} --- quanti errori e quali sono stati commessi.
  \item \textbf{Satisfaction} --- rilevata a fine sessione
  tramite scala SUS (10 item). Target baseline atteso
  $< 68/100$ sul tool attuale.
\end{itemize}

\vspace{0.2cm}

\subsection*{Recruitment of subjects}
\begin{tcolorbox}[colback=backgroundColor,colframe=borderColor,arc=2mm]

\textbf{Laura Palumbo} --- 70 anni, F.\\[0.3cm]
\textit{Competenza tecnica:} Bassa --- utilizza lo smartphone
esclusivamente per chiamate e messaggi WhatsApp; non ha mai
acquistato nulla online né utilizzato servizi della pubblica
amministrazione digitale.\\[0.2cm]
\textit{Competenza di dominio:} Media --- usa il bus urbano
quotidianamente per recarsi nel centro di Bologna usufruendo
di un abbonamento acquisito grazie all'aiuto dei familiari.\\[0.2cm]
\textit{Somiglianza con il segmento:} Alta --- profilo
pienamente corrispondente a \textit{L'anziana pendolare}
\hyperref[seg:anziana]{[S1]}\\[0.2cm]
\textit{Connessione con il ricercatore:} Zia del ricercatore.

\vspace{0.2cm}

\textbf{Sergio Mercurio} --- 67 anni, M.\\[0.3cm]
\textit{Competenza tecnica:} Media --- utilizza regolarmente
smartphone, email e applicazioni bancarie; ha familiarità
con la navigazione web ma non usa abitualmente portali
della pubblica amministrazione né servizi di acquisto
online complessi.\\[0.2cm]
\textit{Competenza di dominio:} Bassa --- non utilizza
il TPL con frequenza non vivendo nella città di Bologna, ne può aver bisogno raramente quando va a trovare il figlio fuorisede.\\[0.2cm]
\textit{Somiglianza con il segmento:} Buona --- profilo
corrispondente alla fascia 65+ con competenza digitale media.\\[0.2cm]
\textit{Connessione con il ricercatore:} Padre del
ricercatore (Stefano Mercurio).

\vspace{0.2cm}

\textbf{Chris Harper} --- 34 anni, M.\\[0.3cm]
\textit{Competenza tecnica:} Alta --- utilizza quotidianamente
smartphone, email, mappe e servizi online in inglese; è abituato
a cercare informazioni da solo ma non conosce bene i sistemi
italiani di autenticazione e bigliettazione digitale.\\[0.2cm]
\textit{Competenza di dominio:} Bassa-media --- ha trascorso circa 2 mesi a Bologna per motivi di lavoro e usa il TPL in modo
regolare per spostarsi tra casa, centro e stazione; ha già
acquistato biglietti singoli ma non ha mai rinnovato un
abbonamento online.\\[0.2cm]
\textit{Somiglianza con il segmento:} Alta --- profilo
corrispondente a \textit{il lavoratore straniero}\hyperref[seg:straniero]{[S2]}\\[0.2cm]
\textit{Connessione con il ricercatore:} Conoscenza personale
del ricercatore tramite rete lavorativa.

\vspace{0.2cm}

\textbf{Maria Borriello} --- 79 anni, F.\\[0.3cm]
\textit{Competenza tecnica:} Molto bassa --- possiede uno
smartphone che utilizza esclusivamente per lo scrolling dei
social network (Facebook, foto familiari); non ha mai
effettuato acquisti online, acceduto a portali della
pubblica amministrazione né utilizzato servizi digitali
transazionali di alcun tipo.\\[0.2cm]
\textit{Competenza di dominio:} Bassa  --- utilizza il TPL
per spostamenti lunghi legati a visite di piacere; dipende interamente dai familiari per tutto ciò
che riguarda la gestione digitale del titolo di viaggio.\\[0.2cm]
\textit{Somiglianza con il segmento:} Molto alta --- profilo
che rappresenta il caso più estremo del segmento
\textit{L'anziana pendolare}\hyperref[seg:anziana]{[S1]}: over 75, competenza digitale
bassissima (solo uso passivo del dispositivo), nessuna
esperienza con interazioni attive online, dipendenza totale
dall'intermediario per la gestione del TPL digitale.\\[0.2cm]
\textit{Connessione con il ricercatore:} Conoscenza familiare
del ricercatore tramite rete di prossimità.

\end{tcolorbox}

\vspace{0.2cm}
\subsection*{Modality of recruitment of subjects}
\begin{tcolorbox}[colback=backgroundColor,colframe=borderColor,arc=2mm]
Il reclutamento è avvenuto grazie ai legami familiari di Stefano Mercurio per i primi due soggetti, mentre Chris Harper è stato reclutato in quanto collega del ricercatore in appoggio presso la sede italiana, con criterio di selezione basato sulla corrispondenza al profilo del segmento target. Maria Borriello è stata reclutata tramite la rete di prossimità familiare di Francesco Castaldi, in quanto sua nonna e rappresenta il profilo target più critico del progetto: anziana over 75 con competenza digitale limitata al solo uso passivo del dispositivo, senza alcuna esperienza di interazione attiva con servizi online.
\end{tcolorbox}
\vspace{0.2cm}

\subsection*{Motivation for choice of subjects}
\begin{tcolorbox}[colback=backgroundColor,colframe=borderColor,arc=2mm]
I soggetti sono stati scelti perché corrispondono pienamente
alla combinazione di segmentazione prioritaria identificata
nel progetto ed erano disponibili e raggiungibili per i
ricercatori. È onesto riconoscere che la scelta è stata
determinata anche dalla facilità di accesso: due soggetti
sono familiari del ricercatore, uno è una conoscenza
diretta disponibile per una simulazione credibile di
utente straniero e Maria Borriello è stata coinvolta
tramite la rete di prossimità. La sua inclusione risponde
a una precisa esigenza progettuale: ampliare il campione
verso il profilo over 75 con uso esclusivamente passivo
del digitale, fascia statisticamente più fragile secondo
i dati ISTAT 2025 e finora non rappresentata da un soggetto
distinto nel campione di test.
\end{tcolorbox}

% ================================================================
% TEST #1
% ================================================================

\newpage
\begin{center}
  {\color{primaryColor}\large\bfseries Test \#1: Sergio Mercurio}\\
  \vspace{0.1cm}
  \color{borderColor}\hrule
\end{center}

\vspace{0.5cm}

\subsection*{Name \normalfont\small\itshape}
\begin{tcolorbox}[colback=backgroundColor,colframe=borderColor,arc=2mm]
Sergio Mercurio
\end{tcolorbox}

\subsection*{Date of test}
\begin{tcolorbox}[colback=backgroundColor,colframe=borderColor,arc=2mm]
1 Giugno 2026
\end{tcolorbox}

\subsection*{Assistant of test (team member)}
\begin{tcolorbox}[colback=backgroundColor,colframe=borderColor,arc=2mm]
Stefano Mercurio
\end{tcolorbox}

\subsection*{Context of test}
\begin{tcolorbox}[colback=backgroundColor,colframe=borderColor,arc=2mm]
Sessione svolta presso il domicilio del soggetto, in una
stanza tranquilla con luce naturale senza la presenza di altre persone esterne. Durata totale: 45
minuti (5 task + compilazione SUS).
\end{tcolorbox}

\subsection*{Quantitative data of tests}
\begin{tcolorbox}[colback=backgroundColor,colframe=borderColor,arc=2mm]
I test sono stati considerati falliti quando si è arreso l'utente.
I tempi sono stati calcolati con timer apposito.
\begin{center}
\renewcommand{\arraystretch}{1.4}
\begin{tabular}{clcccc}
\toprule
\textbf{Task} & \textbf{Descrizione} &
\textbf{Esito} & \textbf{Tempo (s)} &
\textbf{Errori} & \textbf{Backtrack} \\
\midrule
T1 & Ricerca orari        & \textbf{F} & 354 & 4 & 3 \\
T2 & Pianificazione itinerario & \textbf{S} & 76 & 0 & 0 \\
T3 & Acquisto abbonamento & \textbf{S} & 245 & 2 & 2 \\
T4 & Agevolazioni e requisiti & \textbf{F} & 312 & 6 & 5 \\
T5 & Rinnovo abbonamento  & \textbf{S} & 209 & 3 & 3 \\
\midrule
\multicolumn{2}{l}{\textbf{Totali / Medie}} &
\textbf{3/5 (60\%)} & \textbf{239,2 s} &
\textbf{3} & \textbf{2,6} \\
\bottomrule
\end{tabular}
\end{center}

\vspace{0.3cm}
\small S = Successo \quad F = Fallimento \quad
Errori = click errati o percorsi sbagliati \quad
Backtrack = ritorni alla pagina precedente
\end{tcolorbox}

\subsection*{SUS score}
\begin{tcolorbox}[colback=backgroundColor,colframe=borderColor,arc=2mm]
\textbf{Punteggio SUS: 40/100}\\[0.2cm]
\small Punteggio inferiore alla soglia di accettabilità
(50/100). Indica gravi problemi di usabilità percepita,
coerenti con le difficoltà osservate durante la sessione.
\end{tcolorbox}

\subsection*{Useful suggestions from subjects}
\begin{tcolorbox}[colback=backgroundColor,colframe=borderColor,arc=2mm]
Al termine della sessione, il soggetto ha osservato spontaneamente: \textit{``Dovrebbero concentrarsi di più su cosa vuole la gente''}, percependo l'interfaccia attuale come eccessivamente dispersiva e sovraccarica di elementi irrilevanti per il cittadino. 
\end{tcolorbox}

\subsection*{Outcomes and reflections useful for design}
\begin{tcolorbox}[colback=backgroundColor,colframe=borderColor,arc=2mm]
Il volume ridotto di errori e backtrack registrato in alcune fasi non è sintomo di fluidità del sistema. Il soggetto, in difficoltà di fronte alla quantità di informazione che si trova dinanzi, preferisce spendere un tempo prolungato nell'esplorazione visiva faticosa di ogni singola schermata piuttosto che cliccare casualmente. Questo blocco è causato direttamente dal disorientamento in homepage \textbf{(\hyperref[iss:D1]{D1})} e dall'architettura informativa promiscua \textbf{(\hyperref[iss:D2]{D2}}).

Nello specifico del primo task (\textit{Ricerca e consultazione orari}) pur avendo isolato visivamente la linea di interesse, non è riuscito a individuare il comando per accedere alle tabelle orarie, lamentando una totale mancanza di intuitività dei componenti grafici.
La tendenza a non distinguere un pulsante interattivo da una stringa di testo statico conferma sperimentalmente la gravità della \textbf{\hyperref[iss:R3]{R3}} (Accessibilità visiva e cognitiva). 

\begin{itemize}[leftmargin=*]
    \item \textbf{Affordance e contrasto visivo (\hyperref[iss:R3]{R3}):} Standardizzare i pulsanti con forme geometriche chiare ed elevazione visiva rispetto allo sfondo. Introdurre micro-interazioni di \textit{hover} e \textit{focus} marcate per rendere immediata e inequivocabile la distinzione dal testo statico.
    \item \textbf{De-cluttering e progressive disclosure (\hyperref[iss:D1]{D1}, \hyperref[iss:D2]{D2}):} Sfoltire il layout eliminando gli elementi informativi secondari (news, avvisi generali).
\end{itemize}
\end{tcolorbox}

% ================================================================
% TEST #2
% ================================================================

\newpage
\begin{center}
  {\color{primaryColor}\large\bfseries Test \#2: Laura Palumbo}\\
  \vspace{0.1cm}
  \color{borderColor}\hrule
\end{center}

\vspace{0.5cm}

\subsection*{Name \normalfont\small\itshape}
\begin{tcolorbox}[colback=backgroundColor,colframe=borderColor,arc=2mm]
Laura Palumbo
\end{tcolorbox}

\subsection*{Date of test}
\begin{tcolorbox}[colback=backgroundColor,colframe=borderColor,arc=2mm]
    5 Giugno 2026
\end{tcolorbox}

\subsection*{Assistant of test (team member)}
\begin{tcolorbox}[colback=backgroundColor,colframe=borderColor,arc=2mm]
Stefano Mercurio
\end{tcolorbox}

\subsection*{Context of test}
\begin{tcolorbox}[colback=backgroundColor,colframe=borderColor,arc=2mm]
Sessione svolta presso il domicilio del soggetto, in un
ambiente familiare e silenzioso. Durata totale: 41 minuti
(5 task + compilazione SUS).
\end{tcolorbox}

\subsection*{Quantitative data of tests}
\begin{tcolorbox}[colback=backgroundColor,colframe=borderColor,arc=2mm]
I test sono stati considerati falliti quando si è arreso l'utente.
I tempi sono stati calcolati con timer apposito.
\begin{center}
\renewcommand{\arraystretch}{1.4}
\begin{tabular}{clcccc}
\toprule
\textbf{Task} & \textbf{Descrizione} &
\textbf{Esito} & \textbf{Tempo (s)} &
\textbf{Errori} & \textbf{Backtrack} \\
\midrule
T1 & Ricerca orari             & \textbf{F} & 211 & 3 & 0 \\
T2 & Pianificazione itinerario & \textbf{F} & 203 & 6 & 0 \\
T3 & Acquisto abbonamento      & \textbf{F} & 180 & 7 & 5 \\
T4 & Agevolazioni e requisiti  & \textbf{F} & 195 & 7 & 4 \\
T5 & Rinnovo abbonamento       & \textbf{S} & 130 & 2 & 1 \\
\midrule
\multicolumn{2}{l}{\textbf{Totali / Medie}} &
\textbf{1/5 (20\%)} & \textbf{183,8 s} &
\textbf{5,0} & \textbf{2,0} \\
\bottomrule
\end{tabular}
\end{center}

\vspace{0.3cm}
\small S = Successo \quad F = Fallimento \quad
Errori = click errati o percorsi sbagliati \quad
Backtrack = ritorni alla pagina precedente
\end{tcolorbox}

\subsection*{SUS score}
\begin{tcolorbox}[colback=backgroundColor,colframe=borderColor,arc=2mm]
\textbf{Punteggio SUS: 37,5/100}\\[0.2cm]
\small Punteggio al di sotto della soglia di minima di accettabilità
(50/100).
\end{tcolorbox}

\subsection*{Useful suggestions from subjects}
\begin{tcolorbox}[colback=backgroundColor,colframe=borderColor,arc=2mm]
Al termine della sessione il soggetto ha espresso chiaramente il proprio disorientamento a causa della troppa densità visiva dell'interfaccia, commentando:
\textit{``L'interfaccia è troppo ingombrante, c'è troppa confusione e non si capisce bene dove bisogna andare e cosa fare. Ho fatto fatica a trovare la sezione delle agevolazioni è stato un percorso veramente dispersivo''}.
\end{tcolorbox}

\subsection*{Outcomes and reflections useful for design}
\begin{tcolorbox}[colback=backgroundColor,colframe=borderColor,arc=2mm]
I dati quantitativi registrano un quadro di usabilità critico, con un tasso di completamento del \textbf{20\%} (1 task superato su 5) e un punteggio \textbf{SUS di 37,5/100}. L'analisi degli errori rileva picchi di 7 occorrenze nei task T3 e T4 e frequenti \textit{backtrack}.

\medskip
\begin{enumerate}[leftmargin=*]
    \item \textbf{Task Orientation (\hyperref[iss:R4]{R4}):} Separare nettamente i flussi informativi da quelli transazionali per eliminare l'incertezza d'azione.
    \item \textbf{Prevenzione ed Error Tolerance (\hyperref[iss:D4]{D4}):} Compensare l'incapacità di recupero autonomo dagli errori introducendo un micro-accompagnamento guidato e messaggi di feedback contestuali gestiti dal layer dell'intermediario.
    \item     \textbf{Scorciatoie per i flussi di tutela (\hyperref[iss:D2]{D2}):} Estrarre la sezione agevolazioni dai nodi profondi dell'architettura informativa, posizionandola come punto di accesso prioritario nella dashboard iniziale tramite un percorso guidato passo-passo (\textit{wizard}).
    Degli altri punti ne abbiamo già discusso nel test precedente.
\end{enumerate}
\end{tcolorbox}

% ================================================================
% TEST #3
% ================================================================

\newpage
\begin{center}
  {\color{primaryColor}\large\bfseries Test \#3: Chris Harper}\\
  \vspace{0.1cm}
  \color{borderColor}\hrule
\end{center}

\vspace{0.5cm}

\subsection*{Name \normalfont\small\itshape}
\begin{tcolorbox}[colback=backgroundColor,colframe=borderColor,arc=2mm]
Chris Harper
\end{tcolorbox}

\subsection*{Date of test}
\begin{tcolorbox}[colback=backgroundColor,colframe=borderColor,arc=2mm]
18 Maggio 2026
\end{tcolorbox}

\subsection*{Assistant of test (team member)}
\begin{tcolorbox}[colback=backgroundColor,colframe=borderColor,arc=2mm]
Francesco Castaldi
\end{tcolorbox}

\subsection*{Context of test}
\begin{tcolorbox}[colback=backgroundColor,colframe=borderColor,arc=2mm]
Sessione svolta presso l'ufficio del ricercatore, in un ambiente tranquillo e informale. Chris Harper ha vissuto a Bologna per circa due mesi per appoggio ad attività lavorativa presso l'azienda italiana e utilizza il TPL per gli spostamenti quotidiani verso il centro, la stazione e la sede di lavoro. Ha dimostrato familiarità con smartphone, mappe ed email, ma ha incontrato difficoltà nella comprensione di alcune etichette in italiano, nella distinzione tra parte informativa e parte transazionale e nel percorso di autenticazione. Durata totale: 47 minuti (5 task + compilazione SUS).
\end{tcolorbox}

\subsection*{Quantitative data of tests}
\begin{tcolorbox}[colback=backgroundColor,colframe=borderColor,arc=2mm]
\begin{center}
\renewcommand{\arraystretch}{1.4}
\begin{tabular}{clcccc}
\toprule
\textbf{Task} & \textbf{Descrizione} &
\textbf{Esito} & \textbf{Tempo (s)} &
\textbf{Errori} & \textbf{Backtrack} \\
\midrule
T1 & Ricerca orari             & \textbf{S} & 226 & 3 & 4 \\
T2 & Pianificazione itinerario & \textbf{S} & 301 & 4 & 5 \\
T3 & Acquisto abbonamento      & \textbf{F} & 498 & 8 & 10 \\
T4 & Agevolazioni e requisiti   & \textbf{F} & 344 & 5 & 7 \\
T5 & Rinnovo abbonamento       & \textbf{F} & 389 & 7 & 9 \\
\midrule
\multicolumn{2}{l}{\textbf{Totali / Medie}} &
\textbf{2/5 (40\%)} & \textbf{352 s} &
\textbf{5,4} & \textbf{7,0} \\
\bottomrule
\end{tabular}
\end{center}
\vspace{0.3cm}
\small S = Successo \quad F = Fallimento \quad
Errori = click errati o percorsi sbagliati \quad
Backtrack = ritorni alla pagina precedente
\end{tcolorbox}

\subsection*{SUS score}
\begin{tcolorbox}[colback=backgroundColor,colframe=borderColor,arc=2mm]
\textbf{Punteggio SUS: 44/100}\\[0.2cm]
\small Punteggio inferiore alla soglia di accettabilità (50/100). La percezione del sistema è risultata debole, con difficoltà legate soprattutto alla lingua, alla comprensione dei passaggi richiesti e alla presenza di barriere di accesso che non si spiegano bene a un utente non italiano.
\end{tcolorbox}

\subsection*{Useful suggestions from subjects}
\begin{tcolorbox}[colback=backgroundColor,colframe=borderColor,arc=2mm]
Al termine della sessione il soggetto ha commentato: \textit{``I can find the timetable, but once I need to buy or renew something, the process becomes harder to follow and I would need clearer guidance in English.''}

Suggerimento rilevante: una guida più esplicita e semplice in inglese per acquisto, rinnovo e validazione potrebbe ridurre la dipendenza da aiuti esterni.
\end{tcolorbox}
\subsection*{Outcomes and reflections useful for design}
\begin{tcolorbox}[colback=backgroundColor,colframe=borderColor,arc=2mm]
I dati di Chris Harper confermano che un utente straniero con buona alfabetizzazione digitale riesce a gestire le funzioni informative base, ma incontra forti difficoltà nei flussi transazionali. Il completamento del 40\% mostra che i problemi non dipendono solo dalla competenza tecnica, ma anche da lingua, terminologia e assenza di supporto contestuale. Questo risultato conferma in modo diretto la gravità della Issue \textbf{\hyperref[iss:D3]{D3}} e della Issue \textbf{\hyperref[iss:D5]{D5}}, cioè linguaggio non semplificato e barriera SPID, che bloccano anche utenti capaci di navigare autonomamente. Il fallimento sui Task 3, 4 e 5 suggerisce che il redesign deve introdurre un percorso guest-friendly, una lingua più accessibile e una guida step-by-step per acquisto, agevolazioni e rinnovo. Il SUS di 44/100 conferma una percezione di usabilità bassa e rafforza anche la Issue \textbf{\hyperref[iss:R5]{R5}}, perché il sistema non offre un supporto sufficientemente guidato nel momento dell’azione.
\end{tcolorbox}
% ================================================================
% TEST #4
% ================================================================

\newpage
\begin{center}
  {\color{primaryColor}\large\bfseries Test \#4: Maria Borriello}\\
  \vspace{0.1cm}
  \color{borderColor}\hrule
\end{center}

\vspace{0.5cm}

\subsection*{Name \normalfont\small\itshape}
\begin{tcolorbox}[colback=backgroundColor,colframe=borderColor,arc=2mm]
Maria Borriello
\end{tcolorbox}

\subsection*{Date of test}
\begin{tcolorbox}[colback=backgroundColor,colframe=borderColor,arc=2mm]
1 Giugno 2026
\end{tcolorbox}

\subsection*{Assistant of test (team member)}
\begin{tcolorbox}[colback=backgroundColor,colframe=borderColor,arc=2mm]
Francesco Castaldi
\end{tcolorbox}

\subsection*{Context of test}
\begin{tcolorbox}[colback=backgroundColor,colframe=borderColor,arc=2mm]
Sessione svolta presso il domicilio del soggetto, in un
ambiente domestico tranquillo con la presenza del solo
ricercatore. Maria Borriello ha 79 anni e utilizza lo
smartphone esclusivamente per lo scrolling passivo dei
social network (prevalentemente foto di famiglia su
Facebook); non ha mai interagito con siti web, portali
istituzionali o servizi di acquisto online. Ha richiesto
una breve fase di orientamento iniziale sul dispositivo
(laptop del ricercatore) prima di poter avviare i task,
in quanto non aveva mai utilizzato un trackpad. Durata
totale: 71 minuti (5 task + compilazione SUS assistita
nella lettura degli item).
\end{tcolorbox}

\subsection*{Quantitative data of tests}
\begin{tcolorbox}[colback=backgroundColor,colframe=borderColor,arc=2mm]
\begin{center}
\renewcommand{\arraystretch}{1.4}
\begin{tabular}{clcccc}
\toprule
\textbf{Task} & \textbf{Descrizione} &
\textbf{Esito} & \textbf{Tempo (s)} &
\textbf{Errori} & \textbf{Backtrack} \\
\midrule
T1 & Ricerca orari             & \textbf{F} & 698 & 14 & 17 \\
T2 & Pianificazione itinerario & \textbf{F} & 724 & 16 & 19 \\
T3 & Acquisto abbonamento      & \textbf{F} & 612 & 12 & 15 \\
T4 & Agevolazioni e requisiti  & \textbf{F} & 543 & 10 & 13 \\
T5 & Rinnovo abbonamento       & \textbf{F} & 681 & 13 & 16 \\
\midrule
\multicolumn{2}{l}{\textbf{Totali / Medie}} &
\textbf{0/5 (0\%)} & \textbf{652 s} &
\textbf{13,0} & \textbf{16,0} \\
\bottomrule
\end{tabular}
\end{center}
\vspace{0.3cm}
\small S = Successo \quad F = Fallimento \quad
Errori = click errati o percorsi sbagliati \quad
Backtrack = ritorni alla pagina precedente
\end{tcolorbox}

\subsection*{SUS score}
\begin{tcolorbox}[colback=backgroundColor,colframe=borderColor,arc=2mm]
\textbf{Punteggio SUS: 28/100}\\[0.2cm]
\small Punteggio largamente inferiore alla soglia di
accettabilità (50/100) e al di sotto di qualsiasi soglia
di usabilità minima. Rappresenta il valore più basso
rilevato nell'intero ciclo di test ed è coerente con il
tasso di completamento dello 0\% e con il profilo del
soggetto: nessuna esperienza di navigazione attiva su
browser, nessuna familiarità con i flussi transazionali
digitali, assenza totale di un modello mentale di
riferimento per l'interazione con un sito web.
\end{tcolorbox}

\subsection*{Useful suggestions from subjects}
\begin{tcolorbox}[colback=backgroundColor,colframe=borderColor,arc=2mm]
Al termine della sessione il soggetto ha commentato:
\textit{``Non ci capisco niente, non so neanche da dove
cominciare. Quando devo fare queste cose le fa mia figlia,
io non potrei mai farlo da sola.''}
Suggerimento rilevante: la dipendenza totale dall'intermediario
è percepita come normale e inevitabile. Il sistema non offre
alcun punto di aggancio per un utente in questa condizione:
nessun segnale visivo, nessuna guida, nessun punto di
ingresso comprensibile. La soggetta ha anche espresso
disorientamento già nella homepage, non riuscendo a
distinguere tra le sezioni del sito né a identificare
un'azione iniziale chiara.
\end{tcolorbox}
\subsection*{Outcomes and reflections useful for design}
\begin{tcolorbox}[colback=backgroundColor,colframe=borderColor,arc=2mm]
I dati di Maria Borriello definiscono il caso limite del progetto: un tasso di completamento dello 0\%, un tempo medio per task di 652 secondi e un SUS di 28/100 rendono evidente che il sito TPER, nella sua forma attuale, è inaccessibile per un'utente over 75 con esperienza digitale molto ridotta. I fallimenti non riguardano solo le funzioni transazionali, ma investono anche le funzioni informative più semplici, confermando la criticità della \textbf{Issue \hyperref[iss:D1]{D1}} e della \textbf{Issue \hyperref[iss:R4]{R4}}, cioè homepage poco chiara e confusione tra navigazione informativa e compito operativo. Il fatto che il soggetto non riesca a completare nessun task rafforza anche la \textbf{Issue \hyperref[iss:D2]{D2}} e la \textbf{Issue \hyperref[iss:D3]{D3}}, perché la struttura di navigazione e il linguaggio usato dal sito non risultano comprensibili né orientanti. Il confronto con Laura Palumbo mostra che anche pochi anni di differenza anagrafica e un minimo uso attivo dello smartphone producono un gap prestazionale rilevante; l'anzianità estrema e l'assenza totale di interazione digitale attiva costituiscono quindi una barriera autonoma. Il SUS complessivo dei quattro test TPER si conferma molto sotto la soglia di accettabilità, rafforzando la necessità di un redesign centrato su chiarezza, guida costante e riduzione drastica del carico cognitivo.
\end{tcolorbox}
\newpage
% ================================================================
% PRELIMINARY USER TESTING CARD — Competitor's tool
% ================================================================

\newpage

\section{PRELIMINARY USER TESTING CARD, This test relates to the competitor's tool (i.e., used as external benchmark)}
\vspace{0.2cm}

\subsection*{Name of tool}
\begin{tcolorbox}[colback=backgroundColor,colframe=borderColor,arc=2mm]
Sito web ufficiale ATAC Roma --- \url{https://www.atac.roma.it}
\end{tcolorbox}

\vspace{0.2cm}

\begin{center}
  {\color{secondaryColor}\large\bfseries Testing protocol}
  \vspace{0.1cm}
  \color{borderColor}\hrule
\end{center}

\vspace{0.2cm}

\subsection*{Testing session setup}
\begin{tcolorbox}[colback=backgroundColor,colframe=borderColor,arc=2mm]
Le sessioni di test sono state condotte \textbf{in presenza},
presso il domicilio del soggetto, nella stessa sessione del test sul sito TPER,
per garantire la comparabilità delle condizioni.
Il dispositivo utilizzato è stato il laptop del ricercatore,
con sistema operativo \textbf{Windows 11}, browser
\textbf{Google Chrome} aggiornato all'ultima versione
disponibile, risoluzione schermo 1920×1080. Il ricercatore
era presente nella stanza ma non ha interagito con il
soggetto durante lo svolgimento dei task  ma si è limitato a spiegare lo scopo del sito e dell'attività, coerentemente
con la metodologia \textit{bottom line data}.
\end{tcolorbox}

\vspace{0.2cm}

\subsection*{Testing method}
\textbf{Discount usability testing}
(2--5 persone)

\vspace{0.4cm}

\subsection*{List of tasks to test}
\begin{tcolorbox}[colback=backgroundColor,colframe=borderColor,arc=2mm]
\begin{enumerate}[leftmargin=*,itemsep=4pt]
  \item \textbf{Ricerca e consultazione orari:} L'utente
  deve individuare il percorso di una specifica linea di superficie (Metro-C Colosseo-Pantano) e analizzare i suoi orari.
  \item \textbf{Pianificazione dell'itinerario:} L'utente
  deve pianificare uno spostamento coi mezzi inserendo un punto di partenza e
  uno di arrivo (Colosseo-Villa Borghese)
  \item \textbf{Ricerca tariffe e scelta del titolo:}
  L'utente deve individuare il titolo di viaggio ordinario
  più adatto per le proprie esigenze (es. Card over 65),
  verificandone il prezzo e le zone di validità.
  \item \textbf{Identificazione requisiti e documenti:}
  Continuare il punto precedente individuando i requisiti e i documenti
  da presentare.
  \item \textbf{Ricarica titoli di viaggio:}
  L'utente deve individuare la sezione per la ricarica della Metrebus Card
  avviando la procedura fino alla schermata di riepilogo.
\end{enumerate}
\end{tcolorbox}

\vspace{0.4cm}

\subsection*{Testing methodology}
\textbf{Bottom line data} --- il
soggetto svolge i task da solo senza interazione con
il ricercatore, nelle stesse condizioni del test TPER,
per garantire la piena comparabilità dei dati quantitativi.

\vspace{0.2cm}

\subsection*{Metrics of tests}
\begin{itemize}[leftmargin=*,itemsep=3pt]
  \item \textbf{Success} --- percentuale di task portati
  a termine con successo.
  \item \textbf{Time} --- tempo impiegato per
  completare ciascun task (in secondi).
  \item \textbf{Errors} --- numero di click errati,
  backtrack e percorsi sbagliati per task.
  \item \textbf{Satisfaction} --- rilevata a fine sessione
  tramite scala SUS (10 item).
\end{itemize}

\vspace{0.2cm}

\subsection*{Recruitment of subjects}
\begin{tcolorbox}[colback=backgroundColor,colframe=borderColor,arc=2mm]
Gli stessi soggetti reclutati per il test sul sito TPER
sono stati coinvolti nel test sul sito ATAC, per garantire
la piena comparabilità dei risultati.
\end{tcolorbox}

% ================================================================
% TEST ATAC #1 — Laura Palumbo
% ================================================================

\newpage
\begin{center}
  {\color{secondaryColor}\large\bfseries
  Test ATAC \#1: Sergio Mercurio}\\
  \vspace{0.1cm}
  \color{borderColor}\hrule
\end{center}

\vspace{0.5cm}

\subsection*{Name}
\begin{tcolorbox}[colback=backgroundColor,colframe=borderColor,arc=2mm]
Sergio Mercurio
\end{tcolorbox}

\subsection*{Date of test}
\begin{tcolorbox}[colback=backgroundColor,colframe=borderColor,arc=2mm]
1 Giugno 2026
\end{tcolorbox}

\subsection*{Assistant of test (team member)}
\begin{tcolorbox}[colback=backgroundColor,colframe=borderColor,arc=2mm]
Stefano Mercurio
\end{tcolorbox}

\subsection*{Context of test}
\begin{tcolorbox}[colback=backgroundColor,colframe=borderColor,arc=2mm]
Sessione svolta presso il domicilio del soggetto, lo stesso giorno del
test TPER. Durata totale: 35 minuti (5 task +
compilazione SUS).
\end{tcolorbox}

\subsection*{Quantitative data of tests}
\begin{tcolorbox}[colback=backgroundColor,colframe=borderColor,arc=2mm]
I test sono stati considerati falliti quando si è arreso l'utente.
I tempi sono stati calcolati con timer apposito.
\begin{center}
\renewcommand{\arraystretch}{1.4}
\begin{tabular}{clcccc}
\toprule
\textbf{Task} & \textbf{Descrizione} &
\textbf{Esito} & \textbf{Tempo (s)} &
\textbf{Errori} & \textbf{Backtrack} \\
\midrule
T1 & Ricerca orari             & \textbf{S} & 133 & 2 & 0 \\
T2 & Pianificazione itinerario & \textbf{S} & 165 & 3 & 2 \\
T3 & Ricerca tariffe           & \textbf{F} & 240 & 6 & 5 \\
T4 & Agevolazioni e requisiti  & \textbf{S} & 263 & 5 & 3 \\
T5 & Acquisto e ricarica       & \textbf{S} & 189 & 4 & 0 \\
\midrule
\multicolumn{2}{l}{\textbf{Totali / Medie}} &
\textbf{4/5 (80\%)} & \textbf{154s} &
\textbf{4} & \textbf{2} \\
\bottomrule
\end{tabular}
\end{center}
\vspace{0.3cm}
\small S = Successo \quad F = Fallimento \quad
Errori = click errati o percorsi sbagliati \quad
Backtrack = ritorni alla pagina precedente
\end{tcolorbox}

\subsection*{SUS score}
\begin{tcolorbox}[colback=backgroundColor,colframe=borderColor,arc=2mm]
\textbf{Punteggio SUS: 60 /100}\\[0.2cm]
\small Punteggio superiore alla soglia di accettabilità
(50/100), ma comunque inferiore alla soglia di buona usabilità.
\end{tcolorbox}

\subsection*{Useful suggestions from subjects}
\begin{tcolorbox}[colback=backgroundColor,colframe=borderColor,arc=2mm]
Al termine della sessione, il soggetto ha commentato spontaneamente: \textit{``Il sito ha troppe scelte, non so quale biglietto prendere. Almeno a Bologna so che c'è l'abbonamento e basta''}. 

Questa osservazione evidenzia un chiaro fenomeno di \textit{choice overload} (sovraccarico della scelta) generato dall'interfaccia di ATAC, confermando la rilevanza della \textbf{Issue} \hyperref[atac:D3]{D3} per il competitor (complessità nell'accessibilità cognitiva delle tariffe).

Inoltre, il soggetto ha espresso l'esigenza di disporre di un \textbf{assistente conversazionale basato su Intelligenza Artificiale} capace di accompagnarlo dinamicamente nei vari passaggi. Tale intuizione rappresenta un'opportunità di design di alto valore strategico, applicabile in egual misura anche all'ecosistema TPER.
\end{tcolorbox}

\subsection*{Outcomes and reflections useful for design}
\begin{tcolorbox}[colback=backgroundColor,colframe=borderColor,arc=2mm]
Il tasso di completamento sale all'\textbf{80\% su ATAC rispetto al 60\% registrato su TPER}, nonostante il soggetto non avesse alcuna familiarità pregressa con la rete dei trasporti romana.

Questo incremento prestazionale è riconducibile al fatto che avendo completato la sessione su TPER pochi minuti prima, l'utente ha parzialmente metabolizzato e costruito un modello mentale relativo alle logiche di navigazione dei portali di trasporto pubblico locale, riducendo l'attrito iniziale nell'esecuzione dei task informativi base (T1 e T2) .

Tuttavia una volta fatto ingresso nella biglietteria di ATAC, l'estrema densità informativa del layout ha indotto un sovraccarico visivo tale da nascondere i meccanismi di filtraggio dei titoli presenti nella pagina. L'utente non si è accorto dei filtri operativi, confermando sperimentalmente la gravità delle \textbf{Issue \hyperref[atac:D1]{D1} e \hyperref[atac:D2]{D2} di ATAC} (architettura visiva pesante e sequenze d'uso non evidenti) . Per il redesign di TPER, questo riscontro sottolinea che non basta organizzare i contenuti in categorie logiche, ma occorre implementare una forte \textbf{riduzione della complessità visiva} e sistemi di \textit{progressive disclosure} per non disorientare il target fragile.
\end{tcolorbox}

% ================================================================
% TEST ATAC #2 — Sergio Mercurio
% ================================================================

\newpage
\begin{center}
  {\color{secondaryColor}\large\bfseries
  Test ATAC \#2: Laura Palumbo}\\
  \vspace{0.1cm}
  \color{borderColor}\hrule
\end{center}

\vspace{0.5cm}

\subsection*{Name}
\begin{tcolorbox}[colback=backgroundColor,colframe=borderColor,arc=2mm]
Laura Palumbo
\end{tcolorbox}

\subsection*{Date of test}
\begin{tcolorbox}[colback=backgroundColor,colframe=borderColor,arc=2mm]
5 Giugno 2026
\end{tcolorbox}

\subsection*{Assistant of test (team member)}
\begin{tcolorbox}[colback=backgroundColor,colframe=borderColor,arc=2mm]
Stefano Mercurio
\end{tcolorbox}

\subsection*{Context of test}
\begin{tcolorbox}[colback=backgroundColor,colframe=borderColor,arc=2mm]
Sessione svolta presso il domicilio del soggetto, lo stesso giorno del test TPER. Durata totale: 34 minuti (5 task +
compilazione SUS).
\end{tcolorbox}

\subsection*{Quantitative data of tests}
\begin{tcolorbox}[colback=backgroundColor,colframe=borderColor,arc=2mm]
\begin{center}
\renewcommand{\arraystretch}{1.4}
\begin{tabular}{clcccc}
\toprule
\textbf{Task} & \textbf{Descrizione} &
\textbf{Esito} & \textbf{Tempo (s)} &
\textbf{Errori} & \textbf{Backtrack} \\
\midrule
T1 & Ricerca orari             & \textbf{S} & 134 & 3 & 2 \\
T2 & Pianificazione itinerario & \textbf{F} & 146 & 5 & 0 \\
T3 & Ricerca tariffe           & \textbf{F} & 210 & 4 & 3 \\
T4 & Agevolazioni e requisiti  & \textbf{F} & 155 & 4 & 2 \\
T5 & Acquisto e ricarica       & \textbf{S} & 90 & 2 & 1 \\
\midrule
\multicolumn{2}{l}{\textbf{Totali / Medie}} &
\textbf{2/5 (40\%)} & \textbf{260 s} &
\textbf{3,6} & \textbf{1,6} \\
\bottomrule
\end{tabular}
\end{center}
\vspace{0.3cm}
\small S = Successo \quad F = Fallimento \quad
Errori = click errati o percorsi sbagliati \quad
Backtrack = ritorni alla pagina precedente
\end{tcolorbox}

\subsection*{SUS score}
\begin{tcolorbox}[colback=backgroundColor,colframe=borderColor,arc=2mm]
\textbf{Punteggio SUS: 30/100}\\[0.2cm]
\small Punteggio al di sotto della soglia di
usabilità minima (50/100).
\end{tcolorbox}

\subsection*{Useful suggestions from subjects}
\begin{tcolorbox}[colback=backgroundColor,colframe=borderColor,arc=2mm]
Al termine della sessione il soggetto ha espresso forte frustrazione per la scarsa responsività cognitiva dell'interfaccia, osservando:
\textit{``La guida all'acquisto è utile, su TPER non c'è niente di simile. Inoltre la ricerca del percorso è snervante: inserisci l'indirizzo e invece di darti la soluzione migliore ti propone mille alternative ambigue, con una mappa microscopica in cui non si riesce a zoomare per orientarsi''}.

\end{tcolorbox}

\subsection*{Outcomes and reflections useful for design}
\begin{tcolorbox}[colback=backgroundColor,colframe=borderColor,arc=2mm]
\noindent I suggerimenti e le lamentele definiscono tre linee d'azione:
\begin{itemize}[leftmargin=*]
    \item \textbf{Modello della Guida all'acquisto:} La funzione di orientamento tariffario di ATAC è un ottimo pattern di supporto operativo da rielaborare e integrare nel redesign di TPER.
    \item \textbf{Interfacce di pianificazione lineari (\hyperref[atac:D2]{D2}):} Evitare l'overload informativo nei widget di calcolo percorso. L'utente fragile non vuole esplorare alternative sistemiche, ma esige un'unica raccomandazione chiara e una mappa ad alta leggibilità.
\end{itemize}
\end{tcolorbox}

% ================================================================
% TEST ATAC #3 — Chris Harper
% ================================================================

\newpage
\begin{center}
  {\color{secondaryColor}\large\bfseries
  Test ATAC \#3: Chris Harper}\\
  \vspace{0.1cm}
  \color{borderColor}\hrule
\end{center}

\vspace{0.5cm}

\subsection*{Name}
\begin{tcolorbox}[colback=backgroundColor,colframe=borderColor,arc=2mm]
Chris Harper
\end{tcolorbox}

\subsection*{Date of test}
\begin{tcolorbox}[colback=backgroundColor,colframe=borderColor,arc=2mm]
18 Maggio 2026
\end{tcolorbox}

\subsection*{Assistant of test (team member)}
\begin{tcolorbox}[colback=backgroundColor,colframe=borderColor,arc=2mm]
Francesco Castaldi
\end{tcolorbox}

\subsection*{Context of test}
\begin{tcolorbox}[colback=backgroundColor,colframe=borderColor,arc=2mm]
Sessione svolta presso l'ufficio del ricercatore, lo stesso giorno del test TPER. Durata totale: 45 minuti (5 task + compilazione SUS).
\end{tcolorbox}

\subsection*{Quantitative data of tests}
\begin{tcolorbox}[colback=backgroundColor,colframe=borderColor,arc=2mm]
\begin{center}
\renewcommand{\arraystretch}{1.4}
\begin{tabular}{clcccc}
\toprule
\textbf{Task} & \textbf{Descrizione} &
\textbf{Esito} & \textbf{Tempo (s)} &
\textbf{Errori} & \textbf{Backtrack} \\
\midrule
T1 & Ricerca orari             & \textbf{S} & 192 & 2 & 3 \\
T2 & Pianificazione itinerario & \textbf{S} & 268 & 3 & 4 \\
T3 & Ricerca tariffe           & \textbf{S} & 341 & 4 & 5 \\
T4 & Agevolazioni e requisiti  & \textbf{F} & 377 & 5 & 7 \\
T5 & Acquisto e ricarica       & \textbf{F} & 406 & 7 & 9 \\
\midrule
\multicolumn{2}{l}{\textbf{Totali / Medie}} &
\textbf{3/5 (60\%)} & \textbf{317 s} &
\textbf{4,2} & \textbf{5,6} \\
\bottomrule
\end{tabular}
\end{center}
\vspace{0.3cm}
\small S = Successo \quad F = Fallimento \quad Errori = click errati o percorsi sbagliati \quad Backtrack = ritorni alla pagina precedente
\end{tcolorbox}

\subsection*{SUS score}
\begin{tcolorbox}[colback=backgroundColor,colframe=borderColor,arc=2mm]
\textbf{Punteggio SUS: 51/100}\\[0.2cm]
\small Punteggio appena sopra la soglia di accettabilità minima, ma ancora lontano dalla soglia di buona usabilità.
\end{tcolorbox}

\subsection*{Useful suggestions from subjects}
\begin{tcolorbox}[colback=backgroundColor,colframe=borderColor,arc=2mm]
Al termine della sessione il soggetto ha commentato:
\textit{``I can find the timetable and ticket prices, but for discounts and online purchase I would need a clearer English path.''}
Suggerimento rilevante: una guida più esplicita e semplice in inglese per acquisto, rinnovo e agevolazioni potrebbe ridurre la dipendenza da supporto esterno.
\end{tcolorbox}
\subsection*{Outcomes and reflections useful for design}
\begin{tcolorbox}[colback=backgroundColor,colframe=borderColor,arc=2mm]
I dati di Chris Harper mostrano che un utente straniero con buona alfabetizzazione digitale riesce a gestire le funzioni informative base e la consultazione delle tariffe, ma incontra difficoltà nei flussi più complessi legati ad agevolazioni e acquisto online. Il completamento del 60\% è superiore a quello registrato su TPER per lo stesso soggetto, ma resta insufficiente per un uso autonomo soddisfacente. Il fallimento su Task 4 e 5 conferma la criticità della \textbf{Issue \hyperref[atac:D3]{D3}} e della \textbf{Issue \hyperref[atac:D5]{D5}}, cioè linguaggio non accessibile e autenticazione troppo onerosa per chi non conosce bene il sistema locale. Il redesign TPER deve quindi offrire un percorso guest-friendly, una lingua più chiara e una guida step-by-step per acquisto, agevolazioni e rinnovo. Il SUS di 51/100 conferma una percezione di usabilità solo minimamente accettabile, evidenziando che per utenti non italiani la chiarezza dei passaggi è determinante quanto la qualità delle informazioni.
\end{tcolorbox}

% ================================================================
% TEST ATAC #4 — Maria Borriello
% ================================================================

\newpage
\begin{center}
  {\color{secondaryColor}\large\bfseries
  Test ATAC \#4: Maria Borriello}\\
  \vspace{0.1cm}
  \color{borderColor}\hrule
\end{center}

\vspace{0.5cm}

\subsection*{Name}
\begin{tcolorbox}[colback=backgroundColor,colframe=borderColor,arc=2mm]
Maria Borriello
\end{tcolorbox}

\subsection*{Date of test}
\begin{tcolorbox}[colback=backgroundColor,colframe=borderColor,arc=2mm]
1 Giugno 2026
\end{tcolorbox}

\subsection*{Assistant of test (team member)}
\begin{tcolorbox}[colback=backgroundColor,colframe=borderColor,arc=2mm]
Francesco Castaldi
\end{tcolorbox}

\subsection*{Context of test}
\begin{tcolorbox}[colback=backgroundColor,colframe=borderColor,arc=2mm]
Sessione svolta presso il domicilio del soggetto, lo stesso giorno del test TPER. Durata totale: 66 minuti (5 task + compilazione SUS).
\end{tcolorbox}

\subsection*{Quantitative data of tests}
\begin{tcolorbox}[colback=backgroundColor,colframe=borderColor,arc=2mm]
\begin{center}
\renewcommand{\arraystretch}{1.4}
\begin{tabular}{clcccc}
\toprule
\textbf{Task} & \textbf{Descrizione} &
\textbf{Esito} & \textbf{Tempo (s)} &
\textbf{Errori} & \textbf{Backtrack} \\
\midrule
T1 & Ricerca orari             & \textbf{F} & 511 & 11 & 14 \\
T2 & Pianificazione itinerario & \textbf{F} & 546 & 12 & 15 \\
T3 & Ricerca tariffe           & \textbf{F} & 478 & 10 & 13 \\
T4 & Agevolazioni e requisiti  & \textbf{F} & 603 & 13 & 16 \\
T5 & Acquisto e ricarica       & \textbf{F} & 624 & 14 & 17 \\
\midrule
\multicolumn{2}{l}{\textbf{Totali / Medie}} &
\textbf{0/5 (0\%)} & \textbf{552 s} &
\textbf{12,0} & \textbf{15,0} \\
\bottomrule
\end{tabular}
\end{center}
\vspace{0.3cm}
\small S = Successo \quad F = Fallimento \quad Errori = click errati o percorsi sbagliati \quad Backtrack = ritorni alla pagina precedente
\end{tcolorbox}

\subsection*{SUS score}
\begin{tcolorbox}[colback=backgroundColor,colframe=borderColor,arc=2mm]
\textbf{Punteggio SUS: 29/100}\\[0.2cm]
\small Punteggio largamente inferiore alla soglia di accettabilità (50/100) e coerente con il profilo di un utente che utilizza lo smartphone solo in modo passivo, senza alcuna esperienza di navigazione attiva.
\end{tcolorbox}

\subsection*{Useful suggestions from subjects}
\begin{tcolorbox}[colback=backgroundColor,colframe=borderColor,arc=2mm]
Al termine della sessione il soggetto ha commentato:
\textit{``Non so proprio dove cliccare, per me queste cose le deve fare mia figlia o mio nipote''.}
Suggerimento rilevante: anche su ATAC, un utente molto anziano e digitalmente fragile non riesce a orientarsi senza un supporto esplicito, confermando la necessità di interfacce molto più guidate e di percorsi semplificati.
\end{tcolorbox}

\subsection*{Outcomes and reflections useful for design}
\begin{tcolorbox}[colback=backgroundColor,colframe=borderColor,arc=2mm]
I dati di Maria Borriello rappresentano il caso limite anche sul competitor ATAC. Il completamento dello 0\% e il SUS di 29/100 indicano che, nonostante alcuni elementi più accessibili rispetto a TPER, il sito ATAC resta inadeguato per un'utente over 75 con bassissima competenza digitale. Il fallimento totale su tutti i task dimostra che il problema non riguarda solo i flussi transazionali, ma anche la semplice comprensione della struttura del sito. Questo rafforza una conclusione centrale del progetto: nessuno dei due siti, nella sua forma attuale, fornisce un'esperienza davvero autonoma per il segmento più fragile. Il redesign TPER deve quindi andare oltre il benchmark ATAC e proporre una modalità di orientamento radicalmente più semplice, con riduzione della complessità visiva, etichette comprensibili e supporto contestuale sempre visibile.
\end{tcolorbox}


% ================================================================
% Conclusion Card
% ================================================================

\newpage

\section{Conclusion Card}

\subsection*{List of issues by importance}

\noindent\textit{Impatto (I): effetto causale del problema sul completamento 
del task, scala 1--5. Persistenza (P): numero e rilevanza dei task affetti, 
scala 1--5. Punteggio composito: I\,+\,P su 10.}

\begin{tcolorbox}[colback=backgroundColor,colframe=borderColor,arc=2mm]
\begin{itemize}[leftmargin=*,itemsep=8pt]

    \item \textbf{\hyperref[iss:D1]{D1} --- Homepage: assenza di gerarchia 
    e CTA primaria}
    \hfill I:\,5 \textbar{} P:\,5 \textbar{} \textbf{10/10}\\
    Affligge tutti e cinque i task in quanto punto di ingresso obbligato 
    dell'intera sessione. 

    \item \textbf{\hyperref[iss:D2]{D2} --- Architettura informativa: 
    mancata separazione tra sezioni informative e transazionali}
    \hfill I:\,5 \textbar{} P:\,5 \textbar{} \textbf{10/10}\\
    Affligge tutti e cinque i task.

    \item \textbf{\hyperref[iss:R5]{R5} --- Assenza di modalità trainer 
    e supporto contestuale}
    \hfill I:\,4 \textbar{} P:\,5 \textbar{} \textbf{9/10}\\
    Issue strutturale: affligge tutti i task per tutti i soggetti. 
    In nessun test il soggetto ha trovato strumenti di guida attivabili 
    in autonomia. È l'issue nucleare del 
    progetto HPHP.

    \item \textbf{\hyperref[iss:D5]{D5} --- Barriera SPID: 
    assenza di alternative e contestualizzazione}
    \hfill I:\,5 \textbar{} P:\,3 \textbar{} \textbf{8/10}\\
    Affligge tutti i task transazionali. Produce un blocco 
    totale e non aggirabile: in Test~3, tutti i task transazionali 
    falliscono anche su un soggetto con competenza digitale media. 
    L'impatto è massimo perché non ammette recupero parziale.

    \item \textbf{\hyperref[iss:R3]{R3} --- Accessibilità visiva: 
    elementi interattivi non distinguibili dal testo statico}
    \hfill I:\,4 \textbar{} P:\,4 \textbar{} \textbf{8/10}\\
    Documentata sperimentalmente in Test~1 e Test~2. 
    In Test~4, con 14--16 errori per task, l'incapacità di distinguere 
    elementi cliccabili è il fattore che rende ogni schermata un ostacolo 
    autonomo.

    \item \textbf{\hyperref[iss:D3]{D3} --- Linguaggio tecnico-burocratico 
    non spiegato contestualmente}
    \hfill I:\,4 \textbar{} P:\,4 \textbar{} \textbf{8/10}\\
    Affligge principalmente T3, T4 e T5, ma emerge anche su T1 in Test~3 
    dove la terminologia della ricerca orari non è intuitiva per un utente 
    non italiano. Causa diretta del fallimento di T3, T4 e T5 in Test~3 
    (SUS 44/100) e contribuisce in modo determinante al 0\% di 
    completamento in Test~4.

    \item \textbf{\hyperref[iss:R4]{R4} --- Task orientation: 
    confusione tra pianificazione viaggio e acquisto titolo}
    \hfill I:\,4 \textbar{} P:\,3 \textbar{} \textbf{7/10}\\
    In Test~2 genera il picco massimo di errori 
    (7 in T3, 7 in T4) con backtrack rispettivamente di 5 e 4. 
    Produce il rischio operativo più grave dell'intero sistema: 
    l'utente può credere di aver acquistato un titolo senza averlo 
    fatto, con conseguenze dirette fuori dal perimetro digitale.

    \item \textbf{\hyperref[iss:D4]{D4} --- Error tolerance: 
    assenza di messaggi di recupero dagli errori}
    \hfill I:\,3 \textbar{} P:\,4 \textbar{} \textbf{7/10}\\
    Si manifesta su tutti i task ogni volta che l'utente commette un 
    errore, trasformando ogni backtrack in un punto di potenziale 
    abbandono. In Test~2 i backtrack complessivi raggiungono quota 10 
    in una singola sessione; in Test~4 la media è di 16,0 per soggetto. 
    Non è causa primaria di fallimento, ma amplifica in modo sistematico 
    tutte le issue precedenti.

    \item \textbf{\hyperref[iss:R1]{R1} --- Trust and credibility: 
    assenza di segnali di sicurezza nel flusso di pagamento}
    \hfill I:\,3 \textbar{} P:\,2 \textbar{} \textbf{5/10}\\
    Affligge T3 e T5. Non rilevabile direttamente dai dati quantitativi 
    dei test, ma coerente con le barriere emotive documentate nelle 
    interviste qualitative del segmento S1.

    \item \textbf{\hyperref[iss:R2]{R2} --- Search: 
    funzione di ricerca interna non funzionale}
    \hfill I:\,1 \textbar{} P:\,1 \textbar{} \textbf{2/10}\\
    Nessun soggetto ha utilizzato la ricerca interna come strategia 
    di navigazione in alcun test. Rimane una issue latente per utenti 
    con competenza digitale media o alta, non prioritaria nel contesto 
    del progetto HPHP.

\end{itemize}
\end{tcolorbox}

\vspace{1cm}

\subsection*{Summary Table of SUS Scores}

\begin{table}[ht]
\centering
\label{tab:sus_scores}
\renewcommand{\arraystretch}{1.3}
\begin{tabular}{|l|c|c|c|c|}
\hline
\textbf{Question} & \textbf{Sergio M.} & \textbf{Laura P.} & \textbf{Chris H.} & \textbf{Maria B.} \\ \hline
Question \#1      & Agree          & Agree          & Neutral        & Neutral          \\ \hline
Question \#2      & Strongly Agree & Strongly Agree & Disagree       & Strongly Agree   \\ \hline
Question \#3      & Disagree       & Disagree       & Neutral        & Strongly Disagree \\ \hline
Question \#4      & Neutral        & Agree          & Agree          & Strongly Agree   \\ \hline
Question \#5      & Disagree       & Neutral        & Neutral        & Neutral          \\ \hline
Question \#6      & Disagree       & Disagree       & Neutral        & Agree            \\ \hline
Question \#7      & Disagree       & Agree          & Neutral        & Agree            \\ \hline
Question \#8      & Agree          & Agree          & Neutral        & Agree            \\ \hline
Question \#9      & Neutral        & Strongly Disagree & Disagree    & Disagree         \\ \hline
Question \#10     & Neutral        & Agree          & Agree          & Agree            \\ \hline
\textbf{\hfill Total Score} & 40.0   & 37.5          & 45.0          & 27.5             \\ \hline
\textbf{\hfill Average}     & \multicolumn{4}{c|}{37.50}                               \\ \hline
\textbf{\hfill Variance}    & \multicolumn{4}{c|}{54.17}                               \\ \hline
\end{tabular}
\caption{Punteggi SUS — Sito TPER (risposte su scala Likert 5 livelli).}
\end{table}

\vspace{0.5cm}

\subsection*{Urgency Curve}
\begin{figure}[H]
\centering
\begin{tikzpicture}[scale=1.15, font=\small]

  % --- Zone di sfondo ---
  \fill[orange!12]
    (0,3.5) -- (2.5,3.5) -- (2.5,2.5) -- (5.5,2.5)
    -- (5.5,5.5) -- (0,5.5) -- cycle;
  \fill[blue!7]
    (0,0) -- (5.5,0) -- (5.5,2.5) -- (2.5,2.5)
    -- (2.5,3.5) -- (0,3.5) -- cycle;

  % --- Griglia ---
  \draw[gray!25, very thin, step=1] (0,0) grid (5.5,5.5);

  % --- Curva di urgenza (polilinea a gradini) ---
  \draw[very thick, black]
    (0,3.5) -- (2.5,3.5) -- (2.5,2.5) -- (5.5,2.5);

  % --- Assi ---
  \draw[->, thick] (0,0) -- (5.8,0)
    node[right, font=\footnotesize] {Persistenza};
  \draw[->, thick] (0,0) -- (0,5.8)
    node[above, font=\footnotesize] {Impatto};

  % --- Tick e label asse X ---
  \foreach \x in {1,2,3,4,5}{
    \draw (\x, 0) -- (\x,-0.08) node[below, font=\footnotesize] {\x};
  }
  % --- Tick e label asse Y ---
  \foreach \y in {1,2,3,4,5}{
    \draw (0,\y) -- (-0.08,\y) node[left, font=\footnotesize] {\y};
  }

  % ============================================================
  % PUNTI (P = asse X, I = asse Y)
  % Scala intera 1–5; punti nella stessa cella leggermente sfalsati
  % ============================================================

  % D1 — P=5, I=5 (offset ↖)
  \filldraw[violet!80!black] (4.82,5.18) circle (3.5pt);
  \node[violet!80!black, font=\footnotesize\bfseries, above left]
    at (4.82,5.18) {D1};

  % D2 — P=5, I=5 (offset ↘)
  \filldraw[violet!80!black] (5.18,4.82) circle (3.5pt);
  \node[violet!80!black, font=\footnotesize\bfseries, below right]
    at (5.18,4.82) {D2};

  % R5 — P=5, I=4
  \filldraw[teal!80!black] (5,4) circle (3.5pt);
  \node[teal!80!black, font=\footnotesize\bfseries, right]
    at (5,4) {\;R5};

  % D5 — P=3, I=5
  \filldraw[red!70!black] (3,5) circle (3.5pt);
  \node[red!70!black, font=\footnotesize\bfseries, above]
    at (3,5) {D5};

  % R3 — P=4, I=4 (offset ↖)
  \filldraw[orange!80!black] (3.82,4.18) circle (3.5pt);
  \node[orange!80!black, font=\footnotesize\bfseries, above left]
    at (3.82,4.18) {R3};

  % D3 — P=4, I=4 (offset ↘)
  \filldraw[orange!80!black] (4.18,3.82) circle (3.5pt);
  \node[orange!80!black, font=\footnotesize\bfseries, below right]
    at (4.18,3.82) {D3};

  % R4 — P=3, I=4
  \filldraw[blue!70!black] (3,4) circle (3.5pt);
  \node[blue!70!black, font=\footnotesize\bfseries, above left]
    at (3,4) {R4};

  % D4 — P=4, I=3
  \filldraw[blue!70!black] (4,3) circle (3.5pt);
  \node[blue!70!black, font=\footnotesize\bfseries, above right]
    at (4,3) {D4};

  % R1 — P=2, I=3
  \filldraw[gray!70] (2,3) circle (3.5pt);
  \node[gray!70, font=\footnotesize\bfseries, above right]
    at (2,3) {R1};

  % R2 — P=1, I=1
  \filldraw[gray!70] (1,1) circle (3.5pt);
  \node[gray!70, font=\footnotesize\bfseries, above right]
    at (1,1) {R2};

\end{tikzpicture}
\caption{Urgency curve --- issue TPER per impatto e persistenza.
  La polilinea separa le issue da correggere immediatamente (zona arancione,
  sopra la curva) da quelle da monitorare nella release successiva
  (zona blu, sotto la curva).}
\label{fig:urgency-curve}
\end{figure}
% --- da completare (a cura del compagno) ---

\newpage

% ================================================================
% DESIGN RECOMMENDATIONS — Phase A Conclusion
% (da designrecco.tex con adattamenti)
% ================================================================

\section{Design Recommendations}

\label{sec:design-recommendations}

Le raccomandazioni di progettazione presentate in questa sezione derivano
direttamente dall'incrocio sistematico dei risultati emersi dallo heuristic
evaluation (Sezione 3), dai field studies (Sezione 4) e dal preliminary user
testing (Sezione 5). Ogni raccomandazione è riconducibile a specifiche evidenze
empiriche e violazioni di linee guida, ed è formulata in termini operativi per
guidare l'implementazione della trainer-mode e del redesign complessivo
dell'ecosistema TPER.

Le raccomandazioni sono organizzate per macro-temi progettuali (architettura
dell'informazione, modalità trainer, autenticazione, accessibilità, supporto
all'errore, multilingua e trasparenza), riflettendo le principali aree di
intervento individuate.

% ---------------------------------------------------------------------
% MACRO-AREA 1: ARCHITETTURA DELL'INFORMAZIONE E HOME PAGE
% ---------------------------------------------------------------------

\subsection{Ridefinizione dell'Architettura dell'Informazione e della Home Page}

\subsubsection{R1 - Home Page Orientata al Compito}
\label{rec:R1}

\textbf{Raccomandazione}: Riprogettare la home page di \texttt{tper.it}
introducendo una gerarchia visiva chiara che privilegi le azioni principali
per l'utente finale (acquisto di un biglietto, pianificazione di un viaggio,
rinnovo dell'abbonamento) rispetto a news aziendali e avvisi secondari.

\textbf{Implementazione}: Collocare le tre CTA principali (“Acquista
biglietto”, “Pianifica viaggio”, “Rinnova abbonamento”) in un blocco
prominente sopra la piega della pagina, utilizzando pulsanti di dimensioni
generose (minimo 44×44px per target tattili) con etichette in grassetto.
Le news aziendali e gli avvisi di servizio vanno spostati in una sezione
secondaria, contratribile o collocata dopo i contenuti prioritari.

\textbf{Evidenze a supporto}:
\begin{itemize}
    \item \textbf{Euristica}: Violazione G1 e G7 (home page poco chiara,
    gerarchia visiva assente) -- classificata come Severa, Persistente,
    Frequente.
    \item \textbf{User testing}: Maria B. (79 anni) ha dichiarato: “Non so
    neanche da dove cominciare”, completando 0 task su 5 a causa del
    disorientamento in homepage. Chris H. (34 anni, competenza alta) ha
    impiegato 226 secondi per il Task 1, confermando la criticit\`a anche
    per utenti con maggiore dimestichezza digitale.
    \item \textbf{Field observation}: Nessuno degli sportellisti TPER mostra
    il sito web pubblico, segno che l’attuale home page non \`e percepita come
    punto di partenza affidabile per l’interazione con l’utente.
\end{itemize}

\subsubsection{R2 - Architettura dell'Informazione Piana e Lineare}
\label{rec:R2}

\textbf{Raccomandazione}: Ristrutturare l’architettura informativa del sito
secondo un modello piatto (flat IA), limitando la profondit\`a di navigazione
a massimo due livelli dal percorso principale per le funzioni transazionali
pi\`u frequenti (acquisto, rinnovo, agevolazioni).

\textbf{Implementazione}: Sostituire l’attuale menu promiscuo con una
struttura a schede orizzontali chiare: \textit{Acquista}, \textit{Pianifica},
\textit{I Miei Viaggi}, \textit{Aiuto}. Ogni sezione deve presentare tutte le
operazioni disponibili in una singola schermata, evitando sottomenu a discesa
che aumentano il carico cognitivo. Per utenti autenticati, aggiungere una
barra laterale persistente con le azioni rapide (rinnovo, ricarica, documento).

\textbf{Evidenze a supporto}:
\begin{itemize}
    \item \textbf{Euristica}: Violazione G3 -- Architettura informativa
    confusa; le etichette usano terminologia aziendale, il percorso per
    acquistare un abbonamento richiede pi\`u di tre click non lineari.
    \item \textbf{User testing}: Sia Chris H. (34 anni, competenza alta) che
    Maria B. (79 anni, competenza minima) hanno fallito il Task 3 (acquisto
    abbonamento). Chris H. ha dichiarato: \textit{“Once I need to buy or renew
    something, the process becomes harder to follow”}, mentre Maria B. non \`e
    riuscita a distinguere il percorso di acquisto dalla navigazione
    informativa, confermando la complessit\`a strutturale dell'architettura
    informativa indipendentemente dal livello di competenza digitale.
\end{itemize}

\subsubsection{R3 - Segnaletica Progressiva nel Flusso Transazionale}
\label{rec:R3}

\textbf{Raccomandazione}: Introdurre un indicatore di avanzamento
(progress bar) visibile in ogni fase dei flussi di acquisto, rinnovo e
caricamento documenti, specialmente nel contesto istituzionale dove il tempo
disponibile \`e limitato a 3--5 minuti per operazione.

\textbf{Implementazione}: Il progress bar deve mostrare:
\begin{enumerate}
    \item il numero totale di passi;
    \item il passo corrente con descrizione testuale semplice (es.
    “Passo 2 di 4 -- Scegli il tipo di abbonamento”);
    \item un indicatore grafico di completamento (es. tratteggio o
    riempimento).
\end{enumerate}
Il componente deve essere fissato in alto nella viewport e rimanere sempre
visibile durante lo scrolling.

\textbf{Evidenze a supporto}:
\begin{itemize}
    \item \textbf{Field observation}: L'operatore allo sportello TPER dispone in media
    di 3--5 minuti per richieste semplici, finestra temporale che
    diventa il vincolo progettuale pi\`u stringente. Nessun flusso guidato pu\`o
    superare questa soglia.
    \item \textbf{User testing}: Il Task 3 (acquisto) \`e stato fallito da
    entrambi i soggetti, con un tempo medio speso di 555 s (Maria B.) e 498 s
    (Chris H.) prima dell'abbandono
    o dell'errore bloccante.
\end{itemize}


% ---------------------------------------------------------------------
% MACRO-AREA 2: MODALIT\`A TRAINER (TRAINER-MODE) E SUPPORTO CAREGIVER
% ---------------------------------------------------------------------

\subsection{Trainer-Mode e Supporto all'Affiancamento}

\subsubsection{R4 - Dualit\`a dell'Interfaccia (Trainer-Mode / User-Mode)}
\label{rec:R4}

\textbf{Raccomandazione}: Progettare una \textit{trainer-mode} attivabile
tramite un interruttore visibile nell'angolo superiore destro dello schermo,
che consenta all'intermediario di accedere a strumenti di guida (tooltip
attivabili, checklist di apprendimento, flussi passo-passo) senza modificare
la struttura dell'interfaccia per l'utente finale.

\textbf{Implementazione}: In trainer-mode:
\begin{itemize}
    \item compariranno delle etichette esplicative (tooltip) a fianco di ogni
    elemento interattivo;
    \item l'intermediario potr\`a aprire una checklist laterale con i task da
    completare durante la sessione;
    \item sar\`a disponibile una “modalit\`a spiegazione” che evidenzia con un
    bordo colorato l'area dello schermo su cui sta spiegando.
\end{itemize}
Le due viste (trainer e user) devono coesistere sulla stessa schermata
(es. split view) durante l'affiancamento in presenza.

\textbf{Evidenze a supporto}:
\begin{itemize}
    \item \textbf{Euristica}: Violazione G9 e R5 -- assenza catastrofica di
    supporto contestuale attivabile e nessuna modalit\`a trainer per guidare
    l'utente fragile in modo strutturato e ripetibile.
    \item \textbf{Segmentazione}: La dimensione “Ruolo nel sistema” \`e stata
    identificata come fondante del progetto: senza distinguere intermediario
    e utente finale non \`e possibile progettare strumenti differenziati.
    \item \textbf{One-on-one interview}: Tutti gli intermediari dichiarano
    di “fare al posto del genitore”, pattern che la trainer-mode deve inibire
    attivamente spostando i controlli principali verso l'utente finale.
\end{itemize}

\subsubsection{R5 - Persistenza dei Progressi (Memoria della Sessione)}
\label{rec:R5}

\textbf{Raccomandazione}: Progettare un sistema di persistenza dei progressi
che permetta all'utente di salvare lo stato della sessione di training (quali
flussi sono stati completati, a quale passo si \`e arrivati in ciascun flusso)
e di riprenderla successivamente dallo stesso punto, senza dover ricominciare
da capo.

\textbf{Implementazione}: Alla prima interruzione — chiusura del browser,
logout o timeout — il sistema deve conservare automaticamente lo stato della
sessione. Al successivo accesso, deve mostrare un messaggio del tipo ``La
volta scorsa eri arrivato qui: vuoi riprendere da dove avevi lasciato o
ricominciare?''. L'utente (o l'intermediario) deve poter cancellare la cronologia
della sessione al termine del percorso formativo. La persistenza non richiede
registrazione: deve funzionare gi\`a dal primo accesso in modalit\`a guest.

\textbf{Evidenze a supporto}:
\begin{itemize}
    \item \textbf{Interviste agli intermediari}: La frustrazione nel ripetere
    le stesse istruzioni senza risultati duraturi \`e stata segnalata da tutti
    gli intervistati.
    \item \textbf{Field observation (biblioteca)}: Il personale della biblioteca
    ha tempo disteso (sessioni pi\`u lunghe) ma raramente lascia materiale
    scritto all'utente. La persistenza dei progressi consente di
    riprendere la sessione senza ricominciare da capo.
\end{itemize}

\subsubsection{R6 - Riepilogo Stampabile Post-Sessione}
\label{rec:R6}

\textbf{Raccomandazione}: Alla fine di ogni sessione di training, generare
automaticamente un riepilogo dei passi completati, con screenshot annotati
delle schermate chiave, da salvare come PDF o stampare.

\textbf{Implementazione}: Il riepilogo deve includere:
\begin{enumerate}
    \item un elenco puntato delle operazioni svolte;
    \item per ogni operazione, il percorso navigazionale (es.
    “Home $\rightarrow$ Acquista $\rightarrow$ Abbonamento mensile $\rightarrow$ Pagamento”);
    \item un QR code che rimanda a un video tutorial della procedura;
    \item i riferimenti di contatto dell'assistenza clienti TPER.
\end{enumerate}
Il documento deve essere formattato con carattere minimo 14pt e contrasto
elevato per essere leggibile anche da utenti anziani.

\textbf{Evidenze a supporto}:
\begin{itemize}
    \item \textbf{Field observation}: Il personale della biblioteca raramente lascia
    materiale all'utente al termine della sessione. Il riepilogo stampabile
    colma questa lacuna e fornisce un riferimento autonomo tra una sessione
    e l'altra.
    \item \textbf{Segmento “Utente allo sportello”}: Gli utenti anziani spesso
    non sanno quali documenti portare e devono tornare in una visita
    successiva. Il riepilogo include una checklist dei documenti necessari
    per la pratica.
\end{itemize}


% ---------------------------------------------------------------------
% MACRO-AREA 3: AUTENTICAZIONE E SPID
% ---------------------------------------------------------------------

\subsection{Riprogettazione dell'Autenticazione: Riduzione della Barriera SPID}

\subsubsection{R7 - Percorso Alternativo senza SPID per le Transazioni Base}
\label{rec:R7}

\textbf{Raccomandazione}: Introdurre un percorso \textit{guest} per l'acquisto
di biglietti singoli e abbonamenti mensili impersonali, che richieda solo
indirizzo email e metodo di pagamento (carta di credito/debito, PayPal),
eliminando l'obbligo di SPID come autenticazione primaria.

\textbf{Implementazione}: Il flusso guest deve:
\begin{enumerate}
    \item essere visibile come CTA secondaria nella schermata di login
    (“Prosegui senza registrazione”);
    \item richiedere solo l'indirizzo email (per inviare la ricevuta) e i
    dati di pagamento;
    \item generare un codice alfanumerico di 8 cifre che funge da
    “biglietto digitale” validabile a bordo.
\end{enumerate}
Per le operazioni che richiedono obbligatoriamente SPID (es. agevolazioni
tariffarie ISEE), il sistema deve mostrare una spiegazione chiara di cosa
sia SPID, come ottenerlo (link al sito AgID) e una stima dei tempi di
attivazione.

\textbf{Evidenze a supporto}:
\begin{itemize}
    \item \textbf{Euristica}: Violazione G4 e D5 -- la procedura SPID non \`e
    accompagnata da spiegazioni, \`e incompatibile con le competenze digitali
    del target anziano.
    \item \textbf{User testing}: Il Task 3 \`e stato fallito da entrambi i
    soggetti esclusivamente a causa della barriera SPID, indipendentemente dal
    livello di competenza digitale. Chris H. (34 anni, competenza
    alta) ha dichiarato: “Once I need to buy or renew something, the process
    becomes harder to follow”.
    \item \textbf{Test competitor (ATAC)}: Anche su ATAC, tutti i soggetti
    hanno fallito l'acquisto online per barriere di autenticazione
    (MyAtac). Questo dato convergente definisce l'autenticazione
    come il principale punto critico da affrontare nel redesign.
\end{itemize}



% ---------------------------------------------------------------------
% MACRO-AREA 4: ACCESSIBILIT\`A COGNITIVA E VISIVA
% ---------------------------------------------------------------------

\subsection{Accessibilit\`a Cognitiva e Visiva per Utenti Anziani e Fragili}

\subsubsection{R8 - Testi Grandi, Contrasto Elevato e Aree Cliccabili}
\label{rec:R8}

\textbf{Raccomandazione}: Portare l'intero sito a conformit\`a WCAG 2.1 livello
AA, con particolare enfasi su:
\begin{itemize}
    \item dimensione del testo base non inferiore a 16px (o 1rem);
    \item contrasto testo/sfondo minimo 4.5:1;
    \item aree cliccabili di almeno 44×44px per touch;
    \item modalit\`a ad alto contrasto attivabile da un interruttore visibile
    in home page.
\end{itemize}

\textbf{Evidenze a supporto}:
\begin{itemize}
    \item \textbf{Euristica}: Violazione G10 e R3 -- testi di dimensione ridotta,
    contrasto insufficiente, aree cliccabili troppo piccole per l'uso su
    smartphone.
    \item \textbf{Dati ISTAT}: Il 72,5\% dei 65-74enni utilizza Internet,
    quota che scende al 35,7\% tra gli over 75. Le competenze digitali almeno
    di base si attestano al 27,4\% nella fascia 65-74 anni.
    \item \textbf{User testing}: Maria B. (79 anni) non \`e riuscita a distinguere
    un pulsante interattivo da una stringa di testo statico, a causa del
    contrasto e della dimensione insufficienti.
\end{itemize}

\subsubsection{R9 - Contenuti dal Tono Incoraggiante e Inclusivo}
\label{rec:R9}

\textbf{Raccomandazione}: Progettare i testi del sito con un tono
incoraggiante, rassicurante e non paternalistico, che riduca l'ansia
digitale dell'utente fragile e lo motivi a proseguire autonomamente.
Il linguaggio deve trasmettere ``ce la puoi fare'' anzich\'e ``questo \`e
un modulo ISEE''.

\textbf{Implementazione}: Ogni flusso deve includere:
\begin{enumerate}
    \item un messaggio di benvenuto motivazionale all'avvio (es. ``Vedrai,
    basta un passo alla volta -- iniziamo insieme?'' non solo ``Procedi'');
    \item testi brevi, in forma aUva e posiUva (``Compila i tuoi dati''
    non ``I dati devono essere compilati'');
    \item micro-testi di rinforzo dopo ogni passo completato (``Perfetto,
    questo \`e fatto! Passiamo al prossimo'');
    \item un tono coerente in tutto il sito: amichevole ma rispeUoso,
    mai infantile.
\end{enumerate}
I contenuti istituzionali e legali (tariffe, normative) devono rimanere
precisi ma essere incorniciati dal tono incoraggiante. Non si tratta di
``semplificare il vocabolario'', ma di progettare la relazione emotiva
con l'utente attraverso le parole.

\textbf{Evidenze a supporto}:
\begin{itemize}
    \item \textbf{User testing}: Maria B. (79 anni) ha dichiarato: ``Non so
    neanche da dove cominciare'' -- il problema non \`e lessicale ma di
    mancanza di una guida emotiva che la accompagni nel primo passo.
    \item \textbf{Segmento ``Immigrato lavoratore''}: Chris H. ha richiesto
    ``clearer guidance'' non in senso tecnico ma come bisogno di un
    accompagnamento che lo faccia sentire sicuro in un contesto
    linguistico e culturale non familiare.
    \item \textbf{Criteri di valutazione}: Ricercare
    esplicitamente il ``Find a tone and a message'' (p.98) e che
    ``tone and messages are important for motivation'' (p.89).
\end{itemize}


% ---------------------------------------------------------------------
% MACRO-AREA 5: SUPPORTO ALL'ERRORE E FEEDBACK
% ---------------------------------------------------------------------

\subsection{Supporto all'Errore e Feedback Contestuale}

\subsubsection{R10 - Messaggi di Errore Esponenziali e Guidati}
\label{rec:R10}

\textbf{Raccomandazione}: Sostituire i messaggi di errore generici
(“Errore -- riprovare”) con messaggi espansivi che spiegano:
\begin{enumerate}
    \item cosa \`e successo (es. “La carta di credito \`e stata rifiutata”);
    \item perch\'e \`e successo (es. “Il codice CVV inserito non \`e corretto”);
    \item come risolvere (es. “Verifica il codice a 3 cifre sul retro della
    carta”);
    \item azioni alternative (es. “Prova con un'altra carta o con PayPal”).
\end{enumerate}

\textbf{Implementazione}: I messaggi di errore devono essere visualizzati in
un box colorato di contrasto (testo su sfondo rosso chiaro con icona di
attenzione), posizionati immediatamente sopra il campo che ha generato
l'errore (inline error). Devono includere un pulsante “Chiedi aiuto” che apre
un canale di contatto diretto (chat o numero verde).

\textbf{Evidenze a supporto}:
\begin{itemize}
    \item \textbf{Euristica}: Violazione G9 -- il sistema non guida l'utente
    nel recupero da errori comuni (pagamento fallito, abbonamento scaduto),
    generando chiamate al call center e code allo sportello.
    \item \textbf{User testing}: Due soggetti hanno abbandonato il flusso di
    acquisto dopo un errore di pagamento non spiegato, senza tentare di
    riprovare.
\end{itemize}

\subsubsection{R11 - Checklist Preventiva per la Documentazione}
\label{rec:R11}

\textbf{Raccomandazione}: Inserire in ogni flusso che richiede caricamento
documenti (agevolazioni, rinnovo abbonamento agevolato) una checklist
preventiva di tutti i documenti necessari, con esempi visivi e la possibilit\`a
di scaricare un modulo riepilogativo.

\textbf{Implementazione}: La checklist deve essere mostrata prima di avviare
il flusso di upload, con icona di check attivabile dall'utente per marcare
ogni documento come “in possesso”. Alla fine, il sistema deve generare una
lista di quelli mancanti e fornire indicazioni su come ottenerli.

\textbf{Evidenze a supporto}:
\begin{itemize}
    \item \textbf{Field observation}: Gli utenti anziani spesso non sanno in
    anticipo quali documenti portare e devono tornare allo sportello in una
    visita successiva.
    \item \textbf{Euristica}: L'assenza di una checklist preventiva aumenta
    i ritorni allo sportello e genera un ciclo di dipendenza dall'assistenza
    fisica.
\end{itemize}


% ---------------------------------------------------------------------
% MACRO-AREA 6: MULTILINGUA E SUPPORTO PER STRANIERI
% ---------------------------------------------------------------------

\subsection{Supporto Multilingua per Utenti Stranieri}

\subsubsection{R12 - Versione Inglese e Araba dei Flussi Transazionali}
\label{rec:R12}

\textbf{Raccomandazione}: Fornire traduzione completa in inglese e in arabo
(le due lingue straniere pi\`u parlate in Emilia-Romagna) di tutti i flussi
transazionali (acquisto, rinnovo, agevolazioni) e delle pagine informative
principali (orari, percorsi, avvisi).

\textbf{Implementazione}: La selezione della lingua deve essere persistente
(ricordata tramite cookie) e accessibile da un selettore fisso nell'header.
Le traduzioni devono essere professionali, non automatiche, e validate da
madrelingua. Per i testi tecnici (SPID, ISEE), affiancare una spiegazione
in lingua semplificata.

\textbf{Evidenze a supporto}:
\begin{itemize}
    \item \textbf{User testing}: Chris H. (34 anni, inglese madrelingua) ha
    fallito i Task 3, 4 e 5 a causa della barriera linguistica, nonostante
    un'alta competenza digitale.
    \item \textbf{Dati ISTAT}: Al 1\textdegree\ gennaio 2025 risiedono in Italia 5.422.000
    cittadini stranieri, con forte concentrazione in Emilia-Romagna.
    \item \textbf{Segmento “Immigrato lavoratore”}: Introducono una doppia
    barriera (linguistica e sistemica) che nessun'altra combinazione presenta.
\end{itemize}

\subsubsection{R13 - Video Tutorial in Lingua con Sottotitoli}
\label{rec:R13}

\textbf{Raccomandazione}: Produrre una serie di brevi video tutorial
(durata massima 90 secondi) per le operazioni principali, doppiati in
italiano, inglese e arabo (o con sottotitoli sovrimpressi).

\textbf{Implementazione}: I video devono mostrare l'interfaccia reale del
sito mentre viene spiegata la procedura, con ingrandimenti sulle aree critiche.
Devono essere accessibili da una sezione dedicata (“Aiuto $\rightarrow$ Video guide”)
e contestualmente suggeriti all'interno dei flussi quando l'utente manifesta
difficolt\`a (es. due errori consecutivi sullo stesso campo).

\textbf{Evidenze a supporto}:
\begin{itemize}
    \item \textbf{User testing}: Chris H. ha dichiarato: “I would need clearer
    guidance in English”.
    \item \textbf{Competitor analysis}: Il sito ATAC Roma non offre video
    tutorial, e questa mancanza \`e stata segnalata come una debolezza da
    tutti i soggetti del test.
\end{itemize}


% ---------------------------------------------------------------------
% MACRO-AREA 7: MONITORAGGIO DEI PROGRESSI
% ---------------------------------------------------------------------

\subsection{Monitoraggio dei Progressi e Self-Efficacy}

\subsubsection{R14 - Dashboard di Avanzamento per l'Intermediario}
\label{rec:R14}

\textbf{Raccomandazione}: Nella trainer-mode, fornire all'intermediario una
dashboard che visualizzi:
\begin{enumerate}
    \item la percentuale di task completati dall'utente nei vari flussi;
    \item il numero di errori compiuti in ciascun flusso (aggregati);
    \item una stima del tempo medio impiegato per flusso.
\end{enumerate}

\textbf{Implementazione}: I dati devono essere raccolti in forma anonima
e resettabili con il consenso dell'utente. La dashboard deve includere una
scala Likert semplificata (3 livelli: positivo, neutro, negativo) che l'intermediario pu\`o far
compilare all'utente a fine sessione per misurare la self-efficacy percepita.

\textbf{Evidenze a supporto}:
\begin{itemize}
    \item \textbf{Business goals}: Il monitoraggio dei progressi \`e esplicitamente
    richiesto negli obiettivi del tool.
    \item \textbf{Metodologia}: Le success metrics includono la misurazione
    della self-efficacy digitale tramite scala Likert prima e dopo la sessione
    di training.
    \item \textbf{User testing}: Il SUS medio dei quattro test TPER \`e risultato
    pari a 37,50/100, ben al di sotto della soglia di accettabilit\`a; il
    monitoraggio dei progressi serve anche a verificare l'efficacia del redesign.
\end{itemize}


% ---------------------------------------------------------------------
% RIEPILOGO DELLE RACCOMANDAZIONI
% ---------------------------------------------------------------------
\newpage
\subsection{Riepilogo delle Raccomandazioni}

\begin{table}[htbp]
\centering
\caption{Panoramica delle raccomandazioni di design per macro-area}
\label{tab:recommendations-summary}
\begin{tabular}{|p{2.8cm}|p{10cm}|}
\hline
\textbf{Macro-area} & \textbf{Raccomandazioni} \\
\hline
Architettura e Home Page & R1 -- Home page orientata al compito \\
& R2 -- IA piana e lineare \\
& R3 -- Segnaletica progressiva \\
\hline
Trainer-mode e supporto & R4 -- Dualit\`a dell'interfaccia \\
& R5 -- Persistenza dei progressi \\
& R6 -- Riepilogo stampabile \\
\hline
Autenticazione e SPID & R7 -- Percorso guest per transazioni base \\
\hline
Accessibilit\`a & R8 -- Testi grandi e contrasto elevato \\
& R9 -- Contenuti dal tono incoraggiante \\
\hline
Supporto all'errore & R10 -- Messaggi di errore espansivi \\
& R11 -- Checklist documenti preventiva \\
\hline
Multilingua & R12 -- Versioni inglese e araba \\
& R13 -- Video tutorial multilingua \\
\hline
Monitoraggio progressi & R14 -- Dashboard intermediario \\
\hline
\end{tabular}
\end{table}

\noindent Le raccomandazioni sopra elencate costituiscono la baseline per la
fase di prototipazione e sviluppo della trainer-mode. Si suggerisce di
adottare un approccio iterativo (design--prototype--test--refine) con il
coinvolgimento diretto dei segmenti di utenza identificati (anziana pendolare,
immigrato lavoratore) per validare
l'efficacia delle soluzioni proposte.

\begin{center}
    {\color{primaryColor}\Large\bfseries Persona and Use Case Card} \\[4pt]
    {\color{secondaryColor}\normalsize Phase B --- Feasibility Study} \\[4pt]
    {\color{secondaryColor}\small  Questo \textit{persona} è il \textbf{protagonista} (cioè colui per cui stiamo progettando il nuovo strumento; questo \textit{persona} deve essere soddisfatto al 100\% dal design; a meno che non vi siano buoni motivi, solo un \textit{persona} nel cast è il protagonista).}

\vspace{0.3cm}
{\color{borderColor}\hrule}
\vspace{0.5cm}

% ================================================================
% RIGA 1: FOTO + DATI ANAGRAFICI
% ================================================================

\begin{minipage}[t]{0.25\textwidth}
    \centering
    \includegraphics[width=\textwidth,height=5cm,keepaspectratio]
    {images/fatima.png}
    \vspace{0.2cm}
\end{minipage}
\hfill
\begin{minipage}[t]{0.70\textwidth}
    \begin{tcolorbox}[colback=backgroundColor,colframe=primaryColor,arc=2mm]
        \textbf{\large Fatima Benali}\label{persona:fatima_benali} \\[4pt]
        {\color{secondaryColor}\small\textbf{Citazione significativa:}
        \textit{``Io non capisco le parole italiane scritte,
        ma so benissimo quale autobus prendere e dove scendere''}} \\[4pt]
        {\color{secondaryColor}\small\textbf{Segmento corrispondente:}
        Anziana straniera (68 anni, origine marocchina,
        alfabetizzata in arabo ma analfabeta in italiano scritto,
        competenza digitale minima, uso frequente del TPL)}
    \end{tcolorbox}
\end{minipage}

\vspace{0.3cm}

% ================================================================
% DESCRIZIONE DEL PERSONA (Basic Facts)
% ================================================================

\begin{tcolorbox}[colback=backgroundColor,colframe=borderColor,arc=2mm]
    \textbf{Descrizione del persona} \\[4pt]
    {\small
    \textbf{Età:} 68 anni \hfill
    \textbf{Genere:} Femmina \hfill
    \textbf{Nazionalità:} Marocchina, residente in Italia da 22 anni
    \\[2pt]
    \textbf{Residenza:} Bologna, quartiere Pilastro \hfill
    \textbf{Stato civile:} Coniugata \\[2pt]
    \textbf{Istruzione:} Scuola primaria completata in Marocco
    in lingua araba. Sa leggere e scrivere in arabo,
    ma non ha mai ricevuto alcuna istruzione in italiano. \\[2pt]
    \textbf{Competenza linguistica:} Italiano orale elementare
    (sufficiente per interazioni di routine: mercato, medico,
    sportello). Italiano scritto: nessuno. Non riconosce le
    lettere dell'alfabeto latino e non è in grado di leggere
    alcuna parola italiana, nemmeno le più comuni.
    Legge e scrive correttamente in arabo: usa WhatsApp
    in arabo per comunicare con i familiari in Marocco
    e con il gruppo del centro Al-Amal. \\[2pt]
    \textbf{Problemi fisici rilevanti:} Ipovisione lieve
    corretta con occhiali da vista. Motricità nella norma
    per l'età, senza limitazioni significative alle mani.
    Affaticamento visivo rapido davanti agli schermi. \\[2pt]
    \textbf{Profilo digitale:} Possiede uno smartphone Android
    acquistato dal figlio. Lo usa quotidianamente per
    videochiamate e messaggi su WhatsApp in arabo con
    i familiari in Marocco e con le amiche del centro Al-Amal.
    Non ha mai navigato su un sito web, non sa cosa sia
    un browser, non conosce il concetto di SPID né di ISEE.
    Non ha mai inserito dati in un modulo online. \\[2pt]
    \textbf{Utilizzo TPL:} Usa i mezzi TPER quasi ogni giorno
da oltre vent'anni per raggiungere il mercato, la moschea
del quartiere e il centro di aggregazione Al-Amal.
Ha un abbonamento mensile agevolato che fino ad oggi
il figlio Karim rinnovava allo sportello ogni mese.
Da quando il centro Al-Amal ha avviato uno sportello
di assistenza digitale con una volontaria dedicata,
Fatima si è rivolta a lei per le pratiche online.
Conosce perfettamente linee, fermate, coincidenze
e zone tariffarie dell'area in cui vive.
    }
\end{tcolorbox}

% ================================================================
% BACKGROUND E CARATTERISTICHE PSICOLOGICHE
% ================================================================

% ================================================================
% BACKGROUND — sostituire i paragrafi su Karim e sul contesto
% ================================================================

\begin{tcolorbox}[colback=backgroundColor,colframe=borderColor,arc=2mm]
    \textbf{Background e caratteristiche psicologiche} \\[4pt]
    {\small
    Fatima è arrivata in Italia nel 2003, raggiungendo il marito
    Rachid (operaio in pensione) già residente a Bologna.
    Ha sempre fatto la casalinga ed ha cresciuto tre figli
    in Italia: il maggiore, Karim (38 anni), lavora in un
    magazzino logistico con turni notturni; la seconda,
    Nadia (34 anni), vive a Modena con la propria famiglia;
    il più giovane, Youssef (28 anni), studia ingegneria
    a Bologna e vive fuori casa. Per anni è stato Karim
    a gestire tutte le pratiche burocratiche della famiglia,
    spesso in fretta tra un turno e l'altro. \\[4pt]

    Fatima sa leggere e scrivere in arabo, ma non ha mai
    imparato l'alfabeto latino: è cresciuta in Marocco dove
    l'istruzione era in arabo, e in 22 anni di vita in Italia
    non ha mai frequentato un corso di lingua strutturato.
    Comprende l'italiano parlato nelle situazioni quotidiane,
    ma davanti a una parola scritta in italiano è come
    trovarsi davanti a un codice sconosciuto:
    non riesce a decifrare nemmeno le lettere singole.
    Ha costruito una vita autonoma e dignitosa attraverso
    una straordinaria \textbf{memoria visiva e spaziale}:
    riconosce i luoghi per la forma degli edifici, i negozi
    per i colori delle insegne, gli autobus per il numero ---
    che ha imparato a riconoscere come forma visiva,
    non come simbolo linguistico.
    Sa orientarsi in tutta Bologna senza leggere un cartello. \\[4pt]

    Il punto di svolta nella vita digitale di Fatima
    è stato il centro Al-Amal. Da circa un anno,
    l'associazione ha avviato uno \textbf{sportello
    di assistenza digitale} gestito da Sara, una
    volontaria di 28 anni di seconda generazione,
    che parla arabo e italiano e conosce bene
    le procedure burocratiche locali. Sara non si
    sostituisce a Fatima: le siede accanto, indica
    lo schermo, aspetta che sia lei a toccare.
    In questo contesto strutturato --- sedute a un
    tavolo, senza la fretta di Karim tra un turno
    e l'altro --- Fatima si sente al sicuro e riesce
    a concentrarsi. Il centro è diventato il luogo
    dove Fatima impara, per imitazione e ripetizione,
    quello che fuori dal centro le sembra inaccessibile. \\[4pt]

    A differenza di Liliana, che sente attivamente la
    frustrazione della propria dipendenza digitale,
    Fatima la vive con \textbf{rassegnazione pragmatica}:
    considera normale che qualcuno la aiuti con queste cose,
    così come lei aiuta le altre donne del centro
    quando si tratta di orientarsi in città con i mezzi.
    Non percepisce ancora l'autonomia digitale come un
    obiettivo desiderabile. Tuttavia, quando Sara le mostra
    che potrebbe rinnovare l'abbonamento da sola
    senza aspettare nessuno, reagisce con curiosità genuina:
    \textit{``Ma si può davvero fare da sola?
    Nessuno me lo ha mai detto''}.
    }
\end{tcolorbox}

\vspace{0.3cm}

% ================================================================
% C&A DIAGRAM
% ================================================================

\begin{tcolorbox}[colback=backgroundColor,colframe=primaryColor,arc=2mm]
    \textbf{Grafico delle competenze e abilità (C\&A Diagram)} \\[4pt]
    {\small
    }

    \vspace{0.3cm}

    \begin{center}
    \includegraphics[width=0.70\textwidth,keepaspectratio]
    {images/fatima_chart.png}
    \end{center}

    \vspace{0.2cm}

    \begin{center}
    \renewcommand{\arraystretch}{1.3}
    \begin{tabular}{|l|c|c|}
        \hline
        \textbf{Dimensione} & \textbf{Punteggio (1--5)} & \textbf{Indicatore} \\
        \hline
        Competenza tecnica &
       \textbf{2} &
        \cellcolor{gray!10} \makebox[50pt]{\rule{10pt}{10pt}\rule{10pt}{10pt}\color{lightgray}\rule{10pt}{10pt}\rule{10pt}{10pt}\rule{10pt}{10pt}} \\
        \hline
        Competenza di dominio &
        \textbf{4} &
        \cellcolor{gray!10} \makebox[50pt]{\rule{10pt}{10pt}\rule{10pt}{10pt}\rule{10pt}{10pt}\rule{10pt}{10pt}\color{lightgray}\rule{10pt}{10pt}} \\
        \hline
        Competenza linguistica &
        \textbf{1} &
        \cellcolor{gray!10} \makebox[50pt]{\rule{10pt}{10pt}\color{lightgray}\rule{10pt}{10pt}\rule{10pt}{10pt}\rule{10pt}{10pt}\rule{10pt}{10pt}} \\
        \hline
        Abilità fisica e visiva &
        \textbf{3} &
        \cellcolor{gray!10} \makebox[50pt]{\rule{10pt}{10pt}\rule{10pt}{10pt}\rule{10pt}{10pt}\color{lightgray}\rule{10pt}{10pt}\rule{10pt}{10pt}} \\
        \hline
        Motivazione &
        \textbf{2} &
        \cellcolor{gray!10} \makebox[50pt]{\rule{10pt}{10pt}\rule{10pt}{10pt}\color{lightgray}\rule{10pt}{10pt}\rule{10pt}{10pt}\rule{10pt}{10pt}} \\
        \hline
        Concentrazione &
        \textbf{2} &
        \cellcolor{gray!10} \makebox[50pt]{\rule{10pt}{10pt}\rule{10pt}{10pt}\color{lightgray}\rule{10pt}{10pt}\rule{10pt}{10pt}\rule{10pt}{10pt}} \\
        \hline
    \end{tabular}
    \end{center}

    \vspace{0.2cm}
    {\small \textbf{Legenda:} \rule{10pt}{10pt} = valore raggiunto \hfill \color{lightgray}\rule{10pt}{10pt}\color{black} = valore non raggiunto}

    \vspace{0.3cm}
    {\small \textit{Nota: A differenza di Liliana (motivazione 3),
    che sente attivamente la frustrazione della dipendenza,
    Fatima la vive come condizione normale e inevitabile.
    La competenza di dominio alta (4) è invece una risorsa
    progettuale fondamentale: Fatima conosce già linee,
    fermate e zone tariffarie TPER. Se il design riesce ad agganciarsi
    a ciò che lei già sa --- numeri di linea, icone delle
    zone, colori delle fermate --- la curva di apprendimento
    si riduce drasticamente anche in presenza di una
    barriera linguistica strutturale.}}
\end{tcolorbox}

\vspace{0.3cm}

% ================================================================
% DETTAGLI INTERESSANTI ADDIZIONALI
% ================================================================

% ================================================================
% DETTAGLI INTERESSANTI — aggiornati al contesto Al-Amal
% ================================================================

\begin{tcolorbox}[colback=backgroundColor!50,colframe=borderColor,arc=2mm]
    \textbf{Dettagli interessanti addizionali} \\[4pt]
    {\small
    \begin{itemize}[leftmargin=*]
        \item Fatima legge e scrive correntemente in arabo
        e usa WhatsApp ogni giorno per messaggiare con
        i familiari rimasti in Marocco e con le amiche
        del centro Al-Amal. Ma davanti a una schermata
        in italiano è completamente disorientata:
        non riconosce le lettere, non può fare affidamento
        sul testo per orientarsi. Riconosce invece il logo
        TPER, i colori delle linee e i numeri degli autobus
        come forme visive consolidate dall'esperienza
        quotidiana. La sua mente elabora l'italiano
        per immagini, non per testo.

        \item Lo sportello digitale di Al-Amal ha un
        computer fisso con una connessione stabile e
        Sara che riceve il martedì e il giovedì pomeriggio.
        Fatima ha imparato a presentarsi nei giorni
        di Sara quando deve fare ``le cose sul computer''.
        Ha già completato due sessioni: nella prima
        ha imparato ad aprire il browser, nella seconda
        ha raggiunto la homepage di TPER. Il rinnovo
        dell'abbonamento sarà la terza sessione.

        \item Al centro Al-Amal ha imparato a usare
        WhatsApp guardando le altre donne farlo,
        senza che nessuno le spiegasse a parole.
        Impara per imitazione visiva e ripetizione,
        non per istruzione verbale: il modello di
        apprendimento più efficace per lei è
        ``guarda e ripeti'', non ``leggi e segui
        le istruzioni''. Sara ha adattato
        il proprio metodo di insegnamento
        a questo stile: mostra prima, parla dopo.

        \item Conserva i suoi documenti importanti
        (carta di soggiorno, tessera sanitaria,
        abbonamento TPER) in una busta trasparente
        ordinata per colore e dimensione. Non legge
        le etichette: riconosce i documenti dalla
        forma, dal colore e dalla foto.
        Sara le ha suggerito di portare sempre
        la busta alle sessioni: la tessera TPER
        sarà il punto di partenza per il flusso
        di rinnovo.
    \end{itemize}
    }
\end{tcolorbox}

% ================================================================
% USE CASE #1 (Scenario Rinnovo - Sessione di Training Guidata)
% ================================================================
\begin{tcolorbox}[colback=primaryColor!10,colframe=primaryColor,arc=2mm]
    \textbf{Descrizione del Use Case \#1: Sessione di training guidata per il rinnovo dell'abbonamento mensile}
\end{tcolorbox}

\vspace{0.3cm}

% --- Narrative Use Case 1 ---
\begin{tcolorbox}[colback=backgroundColor,colframe=borderColor,arc=2mm]
    {\small
    \textbf{Scenario:} È martedì pomeriggio e Fatima si trova allo sportello digitale del centro Al-Amal per la sua terza sessione con la volontaria Sara. L'abbonamento mensile TPER di Fatima è scaduto; lei sa perfettamente di averne bisogno per i suoi spostamenti quotidiani e ha portato con sé la sua busta trasparente contenente la tessera plastificata. Sara siede accanto a lei: non vuole sostituirsi a Fatima, ma guidarla passo dopo passo, stimolando la sua imitazione visiva. 
    \textbf{Obiettivo:} Completare il rinnovo dell'abbonamento mensile tramite PC sotto la guida del tutor, cercando di memorizzare i passaggi visivi per il futuro. \\[4pt]

    \textbf{Sequenza delle azioni:}
    \begin{enumerate}
        \item Fatima accende il computer e apre il browser da sola, ripetendo con successo i gesti appresi nelle settimane precedenti. Digita l'indirizzo aiutandosi con un foglio di promemoria visivo preparato da Sara.
        \item Una volta caricata la homepage di TPER, l'autonomia di Fatima si interrompe. La pagina è densa di scritte in italiano, avvisi e banner mobili. Non potendo leggere il testo in alfabeto latino, Fatima viene immediatamente sopraffatta dal rumore visivo e perde la concentrazione, iniziando a fissare aree casuali dello schermo.
        \item Sara cerca di spiegarle dove cliccare, ma non avendo a disposizione strumenti di evidenziazione remota o modalità di visualizzazione semplificate, è costretta a indicare fisicamente lo schermo con il dito, dicendo: ``Clicca qui sul rettangolo rosso''. Fatima esegue l'azione per pura imitazione spaziale, senza comprendere il significato della parola scritta.
        \item Il portale apre il modulo transazionale richiedendo il codice della tessera. Fatima estrae il documento dalla busta trasparente. Il sistema però non mostra alcuna corrispondenza grafica o fotografia della carta reale per far capire dove si trovi il dato; Sara deve spiegare a voce quale sequenza numerica copiare.
        \item Fatima inizia a digitare la lunga stringa alfanumerica sulla tastiera. A causa della scarsa dimestichezza con l'input digitale e di un rapido affaticamento visivo, commette un errore di battitura saltando due cifre.
        \item Al momento dell'invio, il sistema rinfresca la pagina e mostra un messaggio di errore accattivante ma scritto interamente in testo rosso italiano. Non potendolo decifrare, Fatima si agita, temendo di aver rotto il computer o di aver cancellato il suo abbonamento. Sara interviene tempestivamente per tranquillizzarla, cancella il campo errato e le chiede di riprovare.
        \item Al secondo tentativo, con Sara che le detta i numeri uno a uno, il codice è corretto. Si passa alla fase di pagamento: la schermata reale di produzione mostra scritte minuscole e richiede i dati della carta. L'assenza di un ambiente di simulazione protetto obbliga a inserire dati bancari reali in un contesto didattico, aumentando l'ansia da errore e rallentando drasticamente i tempi della spiegazione. Fatima, sfinita dal carico di attenzione richiesto, lascia che sia Sara a digitare gli ultimi codici di sicurezza, limitandosi a fare clic sul pulsante finale sotto la sua stretta supervisione.
    \end{enumerate}

    \textbf{Barriere previste:}
    \begin{itemize}
        \item \textbf{Assenza di strumenti per il tutor (No Trainer Mode):} La mancanza di indicatori visivi condivisi o di un'interfaccia speculare semplificata costringe il tutor a una faticosa e continua assistenza fisica diretta sul monitor dell'utente.
        \item \textbf{Mancanza di corrispondenze analogico-digitali:} Il form richiede codici complessi senza supportare visivamente l'utente nell'individuazione degli stessi sulla propria tessera fisica, provocando errori di digitazione.
        \item \textbf{Muro linguistico e testuale:} La presenza di notifiche di errore puramente testuali disorienta completamente l'utente non alfabetizzata in italiano, tramutando un intoppo operativo in un blocco emotivo.
    \end{itemize}

    \textbf{Success criteria (parzialmente raggiunti):}
    \begin{itemize}
        \item L'abbonamento viene rinnovato con successo alla fine della sessione.
        \item \textbf{Fallimento dell'obiettivo di apprendimento:} Fatima esce dalla sessione molto stanca e frustrata; sente di aver completato il compito solo grazie alla presenza fisica di Sara e dichiara di non sentirsi affatto in grado di ripeterlo da sola a casa.
    \end{itemize}

    \textbf{Stima dei tempi:} Circa 12 minuti di sessione sincrona, caratterizzata da continue interruzioni, correzioni manuali e momenti di forte esitazione.
    }
\end{tcolorbox}

% ================================================================
% CONSIDERAZIONI CONCLUSIVE --- IMPATTO SUL DESIGN
% ================================================================
\begin{tcolorbox}[colback=primaryColor!10,colframe=primaryColor,arc=2mm]
    \textbf{Considerazioni conclusive --- Impatto del persona e del caso d'uso sul design}
\end{tcolorbox}

\vspace{0.3cm}

\begin{tcolorbox}[colback=backgroundColor,colframe=borderColor,arc=2mm]
    {\small
    \textbf{Impatto del persona Fatima Benali sul design:} \\[4pt]
    Il profilo psicografico e cognitivo di Fatima impone una profonda riconfigurazione dell'interfaccia, trasformandola da un sistema basato sul testo a un ecosistema prettamente visivo, localizzato ed emotivamente abilitante:
    \begin{itemize}[leftmargin=*]
        \item \textbf{Localizzazione in lingua araba e layout speculare (R12):} Per superare l'analfabetismo in lingua italiana e l'incapacità di decodificare l'alfabeto latino, l'intero flusso di rinnovo deve essere disponibile in lingua araba.
        \item \textbf{Interfaccia visiva basata sull'iconografia pura (R1, R8):} Sfruttando la straordinaria memoria visivo-spaziale dell'utente, i testi esplicativi complessi vanno sostituiti da elementi grafici ingranditi e codici colore familiari che richiamino l'esperienza reale del trasporto pubblico (es. il logo TPER, i colori identificativi delle linee urbane).
        \item \textbf{Semplificazione strutturale e Guest Checkout (R2, R7):} L'assoluta mancanza di competenze digitali sistemiche (SPID, moduli online) richiede un'architettura informativa piatta. Il rinnovo dell'abbonamento mensile standard deve configurarsi come operazione transazionale di base accessibile in modalità \textit{guest}, richiedendo esclusivamente l'identificativo della card senza obbligo di login.
        \item \textbf{Accompagnamento emotivo e riduzione dell'ansia (R9):} Per contrastare l'atteggiamento di rassegnazione pragmatica e il timore di ``rompere il computer'', l'interfaccia deve adottare un tono di voce caloroso, incoraggiante e rassicurante. Messaggi di rinforzo positivo devono scansionare ogni interazione per consolidare la fiducia in sé stessa.
    \end{itemize}
    }
\end{tcolorbox}

\vspace{0.3cm}

\begin{tcolorbox}[colback=backgroundColor,colframe=borderColor,arc=2mm]
    {\small
    \textbf{Impatto del Use Case \#1 (Sessione di Training Guidata) sul design:} \\[4pt]
    Le frizioni e i fallimenti didattici riscontrati durante la simulazione della sessione di addestramento tra Fatima e la volontaria Sara determinano i seguenti requisiti per lo sviluppo della \textit{Trainer-Mode}:
    \begin{itemize}[leftmargin=*]
        \item \textbf{Attivazione degli strumenti per il tutor (R4, R14):} L'introduzione del selettore della \textit{Trainer-Mode} deve fornire a Sara un set di strumenti dedicati alla co-presenza. Questo include la ``modalità spiegazione'' per evidenziare a distanza aree dello schermo sul monitor dell'utente e una checklist laterale dei sotto-obiettivi formativi, evitando che il tutor si sostituisca fisicamente all'utente nella digitazione.
        \item \textbf{Corrispondenza analogico-digitale guidata (R10, R11):} Per azzerare gli errori di battitura causati dall'affaticamento visivo durante la copia del codice abbonamento, il form digitale deve mostrare una riproduzione fotografica interattiva della tessera TPER reale. Il sistema deve evidenziare visivamente sulla carta simulata l'esatta posizione del numero da copiare, supportando lo stile di apprendimento per imitazione di Fatima.
        \item \textbf{Messaggi di errore espansivi, localizzati e protetti (R10, R12):} In caso di input errato, il sistema deve bandire i testi di errore generici in lingua italiana (che generano blocchi emotivi nell'utente). L'errore deve essere gestito *inline* tramite un box ad alto contrasto tradotto in arabo, con un tono che normalizzi lo sbaglio e guidi visivamente alla correzione. Inoltre, la presenza di un ambiente di simulazione protetto permette di esercitarsi sul flusso di pagamento senza l'ansia di inserire credenziali bancarie reali.
        \item \textbf{Memoria della sessione e supporti all'autonomia post-sessione (R5, R6, R13):} Per evitare che lo sforzo didattico si esaurisca al termine dell'incontro, il sistema deve salvare la cronologia dei progressi della sessione di training. Al termine, la piattaforma deve generare un riepilogo stampabile o salvabile in PDF contenente gli screenshot annotati dei passaggi chiave e un QR code che rimanda a brevi video tutorial (durata max 90 secondi) doppiati in lingua araba, consentendo a Fatima di ripetere autonomamente l'operazione a casa per imitazione visiva.
    \end{itemize}
    }
\end{tcolorbox}

\vspace{0.3cm}

\end{center}

\begin{center}
    {\color{primaryColor}\Large\bfseries Persona and Use Case Card} \\[4pt]
    {\color{secondaryColor}\normalsize Phase B --- Feasibility Study} \\[4pt]
    {\color{secondaryColor}\small Questo persona non è il protagonista (cioè non è colui per cui stiamo progettando il nuovo strumento; `e un
membro del cast di supporto che però va comunque tenuto in considerazione per garantire un design inclusivo).}
\end{center}

\vspace{0.3cm}
{\color{borderColor}\hrule}
\vspace{0.5cm}

% ================================================================
% RIGA 1: FOTO + DATI ANAGRAFICI
% ================================================================

\begin{minipage}[t]{0.25\textwidth}
    \centering
    \includegraphics[width=\textwidth,height=5cm,keepaspectratio]{images/persona_protagonist.jpg}
    \vspace{0.2cm}
\end{minipage}
\hfill
\begin{minipage}[t]{0.70\textwidth}
    \begin{tcolorbox}[colback=backgroundColor,colframe=primaryColor,arc=2mm]
        \textbf{\large Liliana Roversi} \\[4pt]
        {\color{secondaryColor}\small\textbf{Citazione significativa:} \textit{``Non so neanche da dove cominciare, per me queste cose le deve fare mia figlia.''}} \\[4pt]
        {\color{secondaryColor}\small\textbf{Segmento corrispondente:} Anziana pendolare (over 75, competenza digitale minima)}
    \end{tcolorbox}
\end{minipage}

\vspace{0.3cm}

% ================================================================
% DESCRIZIONE DEL PERSONA (Basic Facts)
% ================================================================

\begin{tcolorbox}[colback=backgroundColor,colframe=borderColor,arc=2mm]
    \textbf{Descrizione del persona} \\[4pt]
    {\small
    \textbf{Età:} 79 anni \hfill
    \textbf{Genere:} Femmina \hfill
    \textbf{Nazionalità:} Italiana \\[2pt]
    \textbf{Residenza:} Bologna, quartiere San Donato (periferia) \hfill
    \textbf{Stato civile:} Vedova, vive sola \\[2pt]
    \textbf{Istruzione:} Licenza elementare \\[2pt]
    \textbf{Problemi fisici rilevanti:} Presbiopia (utilizza occhiali da lettura), artrosi alle mani che riduce la motricità fine, ridotta mobilità (cammina con bastone per tragitti lunghi). \\[2pt]
    \textbf{Profilo digitale:} Possiede uno smartphone regalatole dalla figlia, ma lo usa esclusivamente per scrolling passivo su Facebook (foto di famiglia) e chiamate WhatsApp. Non ha mai utilizzato un browser per navigazione attiva, non ha mai compilato un modulo online, non sa cosa sia SPID. Durante i test è stata necessaria una fase di orientamento iniziale sul trackpad del computer. \\[2pt]
    \textbf{Utilizzo TPL:} Utente TPER da oltre 40 anni. Conosce bene le linee, le fermate e le zone tariffarie dell'esperienza fisica. Acquistava i biglietti in edicola e rinnovava l'abbonamento allo sportello TPER in stazione. Da quando molti servizi sono migrati online, dipende interamente dalla figlia per ogni operazione digitale.
    }
\end{tcolorbox}

% ================================================================
% BACKGROUND E CARATTERISTICHE PSICOLOGICHE
% ================================================================

\begin{tcolorbox}[colback=backgroundColor,colframe=borderColor,arc=2mm]
    \textbf{Background e caratteristiche psicologiche} \\[4pt]
    {\small
    Liliana è nata e cresciuta a Bologna, ha lavorato per 35 anni come collaboratrice scolastica in una scuola elementare del quartiere. È andata in pensione 14 anni fa. Suo marito, ex ferroviere, è deceduto 10 anni fa; da allora vive da sola in un appartamento al terzo piano senza ascensore. Ha due figli: la figlia maggiore, Elena (45 anni, impiegata comunale), vive in un comune limitrofo e la visita ogni martedì e sabato; il figlio minore, Marco (41 anni), vive a Milano e torna a Bologna una volta al mese. \\[4pt]

    Liliana è una donna orgogliosa e abituata all'autosufficienza: ha sempre gestito la casa, le finanze e le pratiche burocratiche in autonomia. La sua dipendenza dalla figlia per le operazioni digitali è per lei fonte di forte frustrazione e vergogna. Ha dichiarato: ``\textit{Mia figlia ha già i suoi problemi, non voglio essere un peso}''. \\[4pt]

    Psicologicamente, Liliana mostra un profondo \textbf{divario tra motivazione e autoefficacia}: vuole imparare, ma ogni tentativo fallito rafforza la convinzione di ``non essere portata'' per la tecnologia. Ha paura di ``rompere qualcosa'' o di ``fare un danno'' se clicca nel posto sbagliato. Preferisce non provare piuttosto che sbagliare e sentirsi umiliata. \\[4pt]

    È una persona metodica e paziente nella vita reale (ama il giardinaggio e la cucina, attività che richiedono precisione e ripetizione), ma l'ambiente digitale la disorienta: la densità di informazioni, le scelte multiple e la mancanza di un percorso fisico ``tattile'' le causano ansia e blocco decisionale.
    }
\end{tcolorbox}

\vspace{0.3cm}

% ================================================================
% C&A DIAGRAM (Radar chart — Competences and Abilities)
% ================================================================

\begin{tcolorbox}[colback=backgroundColor,colframe=primaryColor,arc=2mm]
    \textbf{Grafico delle competenze e abilità (C\&A Diagram)} \\[4pt]
    {\small
    Ogni dimensione è valutata su una scala da 1 (molto bassa) a 5 (molto alta). \\
    I punteggi di Liliana riflettono il profilo emerso dai test utente e dalle osservazioni sul campo.
    }

    \vspace{0.3cm}

    \begin{center}
    \includegraphics[width=0.70\textwidth,keepaspectratio]{images/liliana_chart.png}
    \end{center}

    \vspace{0.2cm}

    \begin{center}
    \renewcommand{\arraystretch}{1.3}
    \begin{tabular}{|l|c|c|}
        \hline
        \textbf{Dimensione} & \textbf{Punteggio (1--5)} & \textbf{Indicatore} \\
        \hline
        Competenza tecnica &
        \textbf{1} &
        \cellcolor{gray!10} \makebox[50pt]{\rule{10pt}{10pt}\color{lightgray}\rule{10pt}{10pt}\rule{10pt}{10pt}\rule{10pt}{10pt}\rule{10pt}{10pt}} \\
        \hline
        Competenza di dominio &
        \textbf{3} &
        \cellcolor{gray!10} \makebox[50pt]{\rule{10pt}{10pt}\rule{10pt}{10pt}\rule{10pt}{10pt}\color{lightgray}\rule{10pt}{10pt}\rule{10pt}{10pt}} \\
        \hline
        Competenza linguistica &
        \textbf{5} &
        \cellcolor{gray!10} \makebox[50pt]{\rule{10pt}{10pt}\rule{10pt}{10pt}\rule{10pt}{10pt}\rule{10pt}{10pt}\rule{10pt}{10pt}} \\
        \hline
        Abilità fisica e visiva &
        \textbf{2} &
        \cellcolor{gray!10} \makebox[50pt]{\rule{10pt}{10pt}\rule{10pt}{10pt}\color{lightgray}\rule{10pt}{10pt}\rule{10pt}{10pt}\rule{10pt}{10pt}} \\
        \hline
        Motivazione &
        \textbf{3} &
        \cellcolor{gray!10} \makebox[50pt]{\rule{10pt}{10pt}\rule{10pt}{10pt}\rule{10pt}{10pt}\color{lightgray}\rule{10pt}{10pt}\rule{10pt}{10pt}} \\
        \hline
        Concentrazione &
        \textbf{2} &
        \cellcolor{gray!10} \makebox[50pt]{\rule{10pt}{10pt}\rule{10pt}{10pt}\color{lightgray}\rule{10pt}{10pt}\rule{10pt}{10pt}\rule{10pt}{10pt}} \\
        \hline
    \end{tabular}
    \end{center}

    \vspace{0.2cm}
    {\small \textbf{Legenda:} \rule{10pt}{10pt} = valore raggiunto \hfill \color{lightgray}\rule{10pt}{10pt}\color{black} = valore non raggiunto}
\end{tcolorbox}

\vspace{0.3cm}

% ================================================================
% DETTAGLI INTERESSANTI ADDIZIONALI
% ================================================================

\begin{tcolorbox}[colback=backgroundColor!50,colframe=borderColor,arc=2mm]
    \textbf{Dettagli interessanti addizionali} \\[4pt]
    {\small
    \begin{itemize}[leftmargin=*]
        \item Liliana ha un quaderno dove la figlia le scrive ``le istruzioni'' per le operazioni digitali, ma ogni volta che il sito cambia (anche di poco) le istruzioni non servono più e lei si blocca.
        \item Conserva ancora tutti i vecchi biglietti cartacei timbrati in un cassetto: per lei il ``ticket'' è un oggetto fisico, non un codice digitale.
        \item Parla al telefono con il figlio a Milano ogni sera, ma non ha mai provato a farsi spiegare qualcosa sullo smartphone al telefono --- ``tanto non capirei''.
        \item La domenica mattina va alla messa nella parrocchia del quartiere, dove si ferma a parlare con altre anziane del quartiere. Molte condividono la sua stessa frustrazione digitale, ma nessuna sa come uscirne.
    \end{itemize}
    }
\end{tcolorbox}

\newpage

% ================================================================
% USE CASE #1
% ================================================================

\begin{tcolorbox}[colback=primaryColor!10,colframe=primaryColor,arc=2mm]
    \textbf{Descrizione del Use Case \#1: Rinnovo autonomo dell'abbonamento mensile TPER}
\end{tcolorbox}

\vspace{0.3cm}

% --- Narrative ---
\begin{tcolorbox}[colback=backgroundColor,colframe=borderColor,arc=2mm]
    \textbf{Narrative del Use Case} \\[4pt]
    {\small
    \textbf{Scenario:} È il primo del mese e l'abbonamento TPER di Liliana è scaduto. Deve andare dal medico di base giovedì mattina e non può permettersi di perdere tempo in coda alla biglietteria. Sua figlia Elena però è impegnata con un corso di formazione per tutta la settimana e non può aiutarla. Liliana si trova di fronte a due opzioni: andare allo sportello (con 40 minuti di fila) o provare a rinnovare online da sola. Sua figlia le ha lasciato un biglietto con scritto ``vai su tper.it e clicca su Rinnova''. Liliana apre il computer che la figlia le ha configurato, si siede al tavolo della cucina e apre Chrome (glielo ha insegnato Elena la settimana prima). \\[4pt]

    \textbf{Obiettivo:} Rinnovare l'abbonamento mensile TPER in autonomia, senza assistenza. \\[4pt]

    \textbf{Sequenza delle azioni:}
    \begin{enumerate}
        \item Aprire il browser Chrome (icona sulla scrivania).
        \item Digitare ``tper.it'' nella barra degli indirizzi (confonde la barra di ricerca Google con la barra degli indirizzi --- primo tentativo fallito, riprova).
        \item Arrivare alla home page di TPER e riconoscere il pulsante ``Rinnova abbonamento'' (tra news, banner pubblicitari e menu a tendina, fa fatica a individuarlo).
        \item Cliccare su ``Rinnova'' e arrivare alla schermata di login. Non ha SPID. Vede il pulsante ``Prosegui senza registrazione'' (guest path) e lo clicca.
        \item Inserire l'indirizzo email (lo legge dal biglietto della figlia) e i dati di pagamento (carta di credito). Ha paura di sbagliare e digita lentamente, controllando ogni cifra tre volte.
        \item Arrivare alla schermata di riepilogo ordine. Controlla che l'importo sia corretto. Esita prima di cliccare ``Conferma pagamento''.
        \item Completare l'acquisto. Il sistema mostra un messaggio di conferma con il codice del biglietto digitale e la data di scadenza del nuovo abbonamento.
        \item Liliana scrive la data di scadenza sul suo calendario da muro, come ha sempre fatto con i rinnovi cartacei.
    \end{enumerate}

    \textbf{Barriere previste:}
    \begin{itemize}
        \item Riconoscimento della barra degli indirizzi vs barra di ricerca (competenza tecnica = 1).
        \item Disorientamento visivo in homepage tra contenuti informativi e azioni transazionali.
        \item Ansia da pagamento online (fiducia = bassa, richiede messaggi rassicuranti).
        \item Paura di cliccare ``conferma'' senza possibilità di annullamento.
    \end{itemize}

    \textbf{Success criteria:}
    \begin{itemize}
        \item Liliana completa il rinnovo in meno di 15 minuti.
        \item Liliana non ha bisogno di chiamare la figlia durante il processo.
        \item Liliana si sente ``soddisfatta'' e ``sollevata'' al termine (non ``in ansia'').
        \item Liliana è in grado di ripetere il processo il mese successivo da sola.
    \end{itemize}

    \textbf{Stima dei tempi:} 20 minuti (con trainer-mode e messaggi di rinforzo, previsti 12--15 minuti dopo familiarizzazione).
    }
\end{tcolorbox}

\newpage

% ================================================================
% USE CASE #2 (optional — secondo scenario)
% ================================================================

\begin{tcolorbox}[colback=primaryColor!10,colframe=primaryColor,arc=2mm]
    \textbf{Descrizione del Use Case \#2: Tentativo fallito di rinnovo online --- barriera SPID e abbandono}
\end{tcolorbox}

\vspace{0.3cm}

\begin{tcolorbox}[colback=backgroundColor,colframe=borderColor,arc=2mm]
    \textbf{Narrative del Use Case} \\[4pt]
    {\small
    \textbf{Scenario:} \`E il pomeriggio del 10 del mese e l'abbonamento mensile di Liliana \`e scaduto ieri. Sua figlia Elena \`e al lavoro fino a sera e non pu\`o aiutarla. Liliana decide di provare da sola. Ricorda che la nipote le ha detto che ``si pu\`o fare tutto online'', quindi accende il computer fisso che Elena le ha lasciato in salotto. \\[4pt]

    \textbf{Obiettivo:} Rinnovare l'abbonamento mensile TPER senza l'aiuto di un intermediario. \\[4pt]

    \textbf{Sequenza:}
    \begin{enumerate}
        \item Liliana apre il browser e digita lentamente ``tper'' nella barra di Google. Clicca sul primo risultato e arriva alla home page di TPER.
        \item Sulla home page cerca un pulsante che dica ``rinnovo'' o ``abbonamento''. Dopo circa 40 secondi di esplorazione visiva, individua un link ``Acquista online'' e ci clicca.
        \item Il sistema la reindirizza a una pagina con diverse opzioni (biglietti, abbonamenti, agevolazioni). Liliana clicca su ``Abbonamenti''.
        \item Compare una schermata che elenca le tipologie di abbonamento. Liliana riconosce il suo abbonamento mensile e clicca ``Acquista''.
        \item A questo punto il sistema mostra una schermata di autenticazione che richiede \textbf{SPID, CIE o CNS}. Liliana non sa cosa sia SPID. Legge le istruzioni ma contiene terminologia burocratica (``identit\`a digitale'', ``identity provider'', ``livello di garanzia SPID 2''). Prova a cliccare su ``Continua senza SPID'' ma il pulsante non esiste.
        \item Dopo 2 minuti di tentativi infruttuosi, Liliana si sente frustrata e inadeguata. Mormora: ``\`E colpa mia, sono troppo vecchia per queste cose.''
        \item Liliana spegne il computer senza completare l'operazione.
        \item \textbf{Soluzione di ripiego attuale:} Il giorno dopo, durante la pausa pranzo, Elena usa il proprio SPID per accedere al sito e rinnovare l'abbonamento di Liliana. Liliana non ha imparato nulla dal processo e la sua percezione di auto-efficacia digitale \`e peggiorata.
    \end{enumerate}

    \textbf{Success criteria (non raggiunti):}
    \begin{itemize}
        \item Liliana non completa nessun passo del flusso di rinnovo.
        \item Liliana abbandona il sito con un senso di frustrazione e colpa.
        \item L'intermediario (Elena) esegue il task \textit{al posto di} Liliana, perpetuando il ciclo di dipendenza.
    \end{itemize}

    \textbf{Stima dei tempi:} 5 minuti (tentativo), 10 minuti (intervento successivo di Elena).
    }
\end{tcolorbox}

\newpage

% ================================================================
% IMPATTO DEL PERSONA SUL DESIGN
% ================================================================

\begin{tcolorbox}[colback=primaryColor!10,colframe=primaryColor,arc=2mm]
    \textbf{Considerazioni conclusive --- Impatto del persona sul design}
\end{tcolorbox}

\vspace{0.3cm}

\begin{tcolorbox}[colback=backgroundColor,colframe=borderColor,arc=2mm]
    {\small
    \textbf{Impatto del persona Liliana Roversi sul design:} \\[4pt]
    Il profilo di Liliana --- competenza tecnica 1/5, motivazione 3/5, concentrazione 2/5 --- ha guidato le seguenti decisioni di design:
    \begin{itemize}[leftmargin=*]
        \item \textbf{Home page orientata al compito (R1):} Liliana non deve ``cercare'' il pulsante di rinnovo. Deve essere il primo elemento che vede, con un'etichetta chiara e un colore contrastante. Il design deve eliminare il rumore visivo (news, banner) che la disorienta.
        \item \textbf{IA piana e lineare (R2):} Liliana non può navigare più di due livelli di profondità. Ogni operazione deve essere completabile in una schermata o in una sequenza lineare chiara con progress bar (R3).
        \item \textbf{Percorso guest senza SPID (R7):} Liliana non ha SPID e non può ottenerlo da sola. Il percorso guest è l'unico modo per renderla autonoma.
        \item \textbf{Tono incoraggiante (R9):} Liliana ha bisogno di essere ``presa per mano'' emotivamente. I messaggi di benvenuto, rinforzo e celebrazione sono progettati per costruire la sua autoefficacia.
        \item \textbf{Errori espansivi (R10):} Liliana non deve mai vedere ``Errore --- riprovare''. Ogni errore deve spiegare cosa è successo, perché e come risolvere, con tono rassicurante.
        \item \textbf{Riepilogo stampabile (R6):} Liliana ha un approccio analogico (scrive le date sul calendario). Il riepilogo stampabile colma il gap tra mondo digitale e abitudini cartacee.
        \item \textbf{Trainer-mode (R4):} Il design separa nettamente il ruolo di Elena (intermediaria che insegna) da quello di Liliana (utente finale che impara facendo). I controlli principali sono sempre dalla parte di Liliana.
    \end{itemize}
    }
\end{tcolorbox}

\vspace{0.3cm}

\begin{tcolorbox}[colback=backgroundColor,colframe=borderColor,arc=2mm]
    {\small
    \textbf{Impatto del Use Case \#1 (Rinnovo autonomo) sul design:} \\[4pt]
    \begin{itemize}[leftmargin=*]
        \item Il flusso di rinnovo deve anticipare l'errore comune ``barra di ricerca vs barra degli indirizzi'' e guidare Liliana al recupero autonomo.
        \item La conferma finale deve includere un messaggio celebrativo (``Complimenti! Ce l'hai fatta!'') e una chiara indicazione della prossima scadenza.
    \end{itemize}
    \vspace{0.2cm}
    \textbf{Impatto del Use Case \#2 (Barriera SPID e abbandono) sul design:} \\[4pt]
    \begin{itemize}[leftmargin=*]
        \item Il percorso guest senza SPID (R7) \`e l'unica soluzione percorribile per utenti come Liliana: senza un'alternativa all'autenticazione forte, il 100\% dei flussi transazionali fallisce prima di iniziare.
        \item La richiesta di SPID deve essere anticipata \textit{prima} che l'utente inizi la compilazione, non dopo: Liliana ha investito tempo e attenzione per arrivare fino al pagamento, e la barriera improvvisa trasforma lo sforzo in frustrazione.
        \item L'abbandono con senso di colpa (``sono troppo vecchia'') dimostra che il sistema non sta solo fallendo un task, ma sta \textit{danneggiando} l'autoefficacia percepita dell'utente. Il tono del redesign (R9) deve controbilanciare attivamente questo effetto.
    \end{itemize}
    }
\end{tcolorbox}

\vspace{0.5cm}


\begin{center}
    {\color{primaryColor}\Large\bfseries Persona and Use Case Card} \\[4pt]
    {\color{secondaryColor}\normalsize Phase B --- Feasibility Study} \\[4pt]
    {\color{secondaryColor}\small Questo \textit{persona} \textbf{non} \`e il protagonista (cio\`e non \`e colui per cui stiamo progettando il nuovo strumento; \`e un membro del cast di supporto che per\`o va comunque tenuto in considerazione per garantire un design inclusivo).}
\end{center}

\vspace{0.3cm}
{\color{borderColor}\hrule}
\vspace{0.5cm}

% ================================================================
% RIGA 1: FOTO + DATI ANAGRAFICI
% ================================================================

\begin{minipage}[t]{0.25\textwidth}
    \centering
    
\includegraphics[width=\textwidth,height=5cm,keepaspectratio]{images/persona_nonprotagonist.jpg}
    \vspace{0.2cm}
    {\small\itshape Foto di repertorio}
\end{minipage}
\hfill
\begin{minipage}[t]{0.70\textwidth}
    \begin{tcolorbox}[colback=backgroundColor,colframe=primaryColor,arc=2mm]
        \textbf{\large David Carter} \\[4pt]
        {\color{secondaryColor}\small\textbf{Citazione significativa:} \textit{``I can find the timetable, but once I need to buy or renew something, the process becomes harder to follow and I need clearer guidance in English.''}} \\[4pt]
        {\color{secondaryColor}\small\textbf{Segmento corrispondente:} Lavoratore straniero (18--35, competenza digitale media, contesto associativo)}
    \end{tcolorbox}
\end{minipage}

\vspace{0.3cm}

% ================================================================
% DESCRIZIONE DEL PERSONA (Basic Facts)
% ================================================================

\begin{tcolorbox}[colback=backgroundColor,colframe=borderColor,arc=2mm]
    \textbf{Descrizione del persona} \\[4pt]
    {\small
    \textbf{Et\`a:} 32 anni \hfill
    \textbf{Genere:} Maschio \hfill
    \textbf{Nazionalit\`a:} Britannica (madrelingua inglese) \\[2pt]
    \textbf{Residenza:} Bologna, zona San Vitale (vicino universit\`a) \hfill
    \textbf{Stato civile:} Single, vive in un appartamento condiviso con altri due colleghi \\[2pt]
    \textbf{Istruzione:} Laurea in Ingegneria Meccanica (Regno Unito) \\[2pt]
    \textbf{Professione:} Ingegnere presso una multinazionale con sede a Bologna (distaccato per 12 mesi) \\[2pt]
    \textbf{Problemi fisici rilevanti:} Nessuno. \\[2pt]
    \textbf{Profilo digitale:} Utilizza quotidianamente smartphone, email, Google Maps, app bancarie e servizi di acquisto online in inglese. Ha familiarit\`a con la navigazione web avanzata, i sistemi di autenticazione a due fattori e il mobile payment del suo paese d'origine. Tuttavia non conosce il sistema SPID italiano, non ha mai utilizzato un portale della pubblica amministrazione italiana e non sa come funzioni la bigliettazione digitale del TPL locale. \\[2pt]
    \textbf{Utilizzo TPL:} Usa il TPER quotidianamente per spostarsi tra casa, sede di lavoro (Zona Roveri) e centro. Ha acquistato biglietti singoli tramite app e rivenditori autorizzati, ma non ha mai completato un rinnovo abbonamento online. Non sa quali agevolazioni tariffarie esistano per stranieri o residenti temporanei.
    }
\end{tcolorbox}

\newpage

% ================================================================
% BACKGROUND E CARATTERISTICHE PSICOLOGICHE
% ================================================================

\begin{tcolorbox}[colback=backgroundColor,colframe=borderColor,arc=2mm]
    \textbf{Background e caratteristiche psicologiche} \\[4pt]
    {\small
    David \`e arrivato a Bologna da circa due mesi per un distacco lavorativo presso la sede italiana della sua azienda. Parla un italiano scolastico di livello base (A2/B1) -- sufficiente per interazioni semplici ma non per comprendere etichette tecniche, testi burocratici o moduli online complessi. \\[4pt]

    \`E una persona abituata all'autonomia digitale: nel Regno Unito gestisce ogni aspetto della sua mobilit\`a online (biglietti ferroviari, abbonamenti metropolitana, ticket bus) attraverso app intuitive e siti ben progettati. L'esperienza con TPER lo frustra perch\'e segue le stesse strategie di navigazione che funzionano a Londra ma si scontrano con un'interfaccia che presume conoscenza del contesto italiano (SPID, MI MUOVO, Io Viaggio, Carta Agevolata, ISEE). \\[4pt]

    Psicologicamente, David mostra un \textbf{divario tra competenza tecnica e competenza di dominio}: sa usare perfettamente la tecnologia, ma non ha il contesto culturale per interpretare il sistema italiano. Questo genera un senso di impotenza specifica -- non ``non so usare il computer'' ma ``non capisco cosa mi stanno chiedendo''. La barriera non \`e emotiva (non ha paura di sbagliare) ma \textbf{cognitivo-linguistica}: non ha gli strumenti lessicali e sistemici per decodificare l'interfaccia. \\[4pt]

    Ha un atteggiamento pratico e determinato: quando incontra un ostacolo cerca attivamente una soluzione (chiede aiuto a colleghi italiani, cerca su Google, prova percorsi alternativi). Tuttavia, dopo ripetuti fallimenti sui flussi transazionali, inizia a mostrare segni di \textbf{affaticamento decisionale} -- troppi passaggi, troppi termini sconosciuti, troppe scelte senza contesto.
    }
\end{tcolorbox}

\vspace{0.3cm}

% ================================================================
% C&A DIAGRAM (Competences and Abilities)
% ================================================================

\begin{tcolorbox}[colback=backgroundColor,colframe=primaryColor,arc=2mm]
    \textbf{Grafico delle competenze e abilit\`a (C\&A Diagram)} \\[4pt]
    {\small
    Ogni dimensione \`e valutata su una scala da 1 (molto bassa) a 5 (molto alta). \\
    I punteggi di David riflettono il profilo emerso dai test utente e dalle osservazioni sul campo.
    }

    \vspace{0.3cm}

    \begin{center}
    \includegraphics[width=0.70\textwidth,keepaspectratio]{images/david_chart.png}
    \end{center}

    \vspace{0.2cm}

    \begin{center}
    \renewcommand{\arraystretch}{1.3}
    \begin{tabular}{|l|c|c|}
        \hline
        \textbf{Dimensione} & \textbf{Punteggio (1--5)} & \textbf{Indicatore} \\
        \hline
        Competenza tecnica &
        \textbf{4} &
        \cellcolor{gray!10} \makebox[50pt]{\rule{10pt}{10pt}\rule{10pt}{10pt}\rule{10pt}{10pt}\rule{10pt}{10pt}\color{lightgray}\rule{10pt}{10pt}} \\
        \hline
        Competenza di dominio &
        \textbf{2} &
        \cellcolor{gray!10} \makebox[50pt]{\rule{10pt}{10pt}\rule{10pt}{10pt}\color{lightgray}\rule{10pt}{10pt}\rule{10pt}{10pt}\rule{10pt}{10pt}} \\
        \hline
        Competenza linguistica &
        \textbf{2} &
        \cellcolor{gray!10} \makebox[50pt]{\rule{10pt}{10pt}\rule{10pt}{10pt}\color{lightgray}\rule{10pt}{10pt}\rule{10pt}{10pt}\rule{10pt}{10pt}} \\
        \hline
        Abilit\`a fisica e visiva &
        \textbf{5} &
        \cellcolor{gray!10} \makebox[50pt]{\rule{10pt}{10pt}\rule{10pt}{10pt}\rule{10pt}{10pt}\rule{10pt}{10pt}\rule{10pt}{10pt}} \\
        \hline
        Motivazione &
        \textbf{4} &
        \cellcolor{gray!10} \makebox[50pt]{\rule{10pt}{10pt}\rule{10pt}{10pt}\rule{10pt}{10pt}\rule{10pt}{10pt}\color{lightgray}\rule{10pt}{10pt}} \\
        \hline
        Concentrazione &
        \textbf{3} &
        \cellcolor{gray!10} \makebox[50pt]{\rule{10pt}{10pt}\rule{10pt}{10pt}\rule{10pt}{10pt}\color{lightgray}\rule{10pt}{10pt}\rule{10pt}{10pt}} \\
        \hline
    \end{tabular}
    \end{center}

    \vspace{0.2cm}
    {\small \textbf{Legenda:} \rule{10pt}{10pt} = valore raggiunto \hfill \color{lightgray}\rule{10pt}{10pt}\color{black} = valore non raggiunto}
\end{tcolorbox}

\vspace{0.3cm}

% ================================================================
% DETTAGLI INTERESSANTI ADDIZIONALI
% ================================================================

\begin{tcolorbox}[colback=backgroundColor!50,colframe=borderColor,arc=2mm]
    \textbf{Dettagli interessanti addizionali} \\[4pt]
    {\small
    \begin{itemize}[leftmargin=*]
        \item Ha un'app sul telefono con tutti i mezzi pubblici delle citt\`a in cui ha vissuto (Londra, Manchester, Edimburgo). Ogni volta che arriva in una nuova citt\`a, cerca automaticamente l'app ufficiale del TPL. Per Bologna ha scaricato l'app TPER ma l'ha disinstallata dopo una settimana perch\'e ``non riusciva a capire come funzionasse la registrazione''.
        \item I colleghi italiani lo aiutano con le pratiche burocratiche, ma si sente in colpa a chiedere continuamente aiuto. Preferirebbe un sistema che pu\`o usare da solo.
        \item Ha provato a tradurre la homepage di TPER con Google Lens dal cellulare, ma la traduzione automatica ha reso le etichette ancora pi\`u incomprensibili.
        \item La domenica va a giocare a calcio in un campetto della zona universitaria con altri expat. Molti condividono la sua stessa frustrazione col sistema digitale italiano e si scambiano ``tips'' su come sopravvivere burocraticamente.
    \end{itemize}
    }
\end{tcolorbox}

\newpage

% ================================================================
% USE CASE #1
% ================================================================

\begin{tcolorbox}[colback=primaryColor!10,colframe=primaryColor,arc=2mm]
    \textbf{Descrizione del Use Case \#1: Acquisto del primo abbonamento mensile TPER come lavoratore straniero}
\end{tcolorbox}



% --- Narrative ---
\begin{tcolorbox}[colback=backgroundColor,colframe=borderColor,arc=2mm]
    \textbf{Narrative del Use Case} \\[4pt]
    {\small
    \textbf{Scenario:} David ha appena ricevuto il primo stipendio e decide di comprare un abbonamento mensile TPER -- ha calcolato che con gli spostamenti casa-lavoro-centro gli conviene rispetto ai biglietti singoli. Un collega italiano gli ha detto ``vai sul sito tper.it e cerca abbonamento mensile''. David apre il laptop aziendale, digita ``tper.it'' nella barra degli indirizzi (operazione familiare, nessun problema) e arriva alla homepage. \\[4pt]

    \textbf{Obiettivo:} Acquistare un abbonamento mensile impersonale TPER online, senza registrazione SPID, in autonomia completa. \\[4pt]

    \textbf{Sequenza delle azioni:}
    \begin{enumerate}
        \item Aprire il browser e navigare su tper.it.
        \item Identificare il percorso verso l'acquisto di un abbonamento nella homepage (fa fatica a trovare la sezione ``Abbonamenti'' tra news aziendali, avvisi e menu a tendina).
        \item Arrivare alla pagina ``Abbonamenti'' e leggere la lista di tipologie (Mensile Impersonale, Mensile Nominativo, Annuale, OVER 65...). L'etichetta ``Mensile Impersonale'' non \`e chiara per un non-madrelingua.
        \item Selezionare ``Mensile Impersonale'' e arrivare alla schermata di login. Non ha SPID. Cerca un pulsante ``Continua senza registrazione'' o ``Guest checkout'' -- non lo trova immediatamente perch\'e \`e poco visibile.
        \item Inserire i dati di pagamento (carta di credito internazionale). Il form chiede ``Codice Fiscale'' -- David non sa cosa sia, non ha un codice fiscale italiano ben visibile, e il campo non ha un'etichetta esplicativa in inglese.
        \item Completare l'acquisto dopo aver chiesto il codice fiscale a un collega italiano via WhatsApp. La conferma arriva in italiano con dettagli sulla validit\`a e le zone.
        \item Scaricare la ricevuta. L'abbonamento \`e digitale -- David non sa se deve mostrare la ricevuta o se c'\`e un QR code da attivare.
    \end{enumerate}

    \textbf{Barriere previste:}
    \begin{itemize}
        \item \textbf{Barriera linguistica (critica):} tutte le etichette sono in italiano, senza supporto multilingua. ``Mensile Impersonale'' non \`e autoesplicativo per un non-madrelingua.
        \item \textbf{Bloom cognitivo ``Codice Fiscale'':} campo obbligatorio senza spiegazione. Per uno straniero, un identificatore fiscale a 16 caratteri alfanumerici \`e un concetto sconosciuto.
        \item \textbf{SPID/autenticazione:} richiedere SPID per acquistare un abbonamento blocca chiunque non abbia cittadinanza italiana o residenza lunga.
        \item \textbf{Assenza di feedback multilingua:} i messaggi di errore e di conferma sono solo in italiano.
        \item \textbf{Validazione del titolo:} dopo l'acquisto, David non sa se l'abbonamento \`e gi\`a attivo o se deve fare qualcosa (attivare, timbrare, convalidare).
    \end{itemize}

    \textbf{Success criteria:}
    \begin{itemize}
        \item David completa l'acquisto in meno di 25 minuti senza assistenza umana.
        \item David non ha bisogno di contattare colleghi italiani per capire termini o passaggi.
        \item David riceve una conferma in inglese (o bilingue) con istruzioni chiare sull'uso del titolo.
        \item David \`e in grado di rinnovare l'abbonamento il mese successivo da solo.
        \item David riferisce di sentirsi ``in controllo'' del processo (non ``confuso'' o ``frustrato'').
    \end{itemize}

    \textbf{Stima dei tempi:} 30--40 minuti (stato attuale), 15--20 minuti (con supporto multilingua e guest path).
    }
\end{tcolorbox}

\newpage

% ================================================================
% USE CASE #2 (opzionale — secondo scenario)
% ================================================================

\begin{tcolorbox}[colback=primaryColor!10,colframe=primaryColor,arc=2mm]
    \textbf{Descrizione del Use Case \#2: Consultazione e verifica agevolazioni tariffarie per residente temporaneo}
\end{tcolorbox}

\vspace{0.3cm}

\begin{tcolorbox}[colback=backgroundColor,colframe=borderColor,arc=2mm]
    \textbf{Narrative del Use Case} \\[4pt]
    {\small
    \textbf{Scenario:} David ha sentito da un collega expat che esistono abbonamenti agevolati per redditi bassi o per categorie specifiche. Ha un reddito medio, ma come residente temporaneo si chiede se pu\`o ottenere uno sconto. Apre il sito TPER e cerca ``agevolazioni'' o ``sconti''. \\[4pt]

    \textbf{Obiettivo:} Trovare le informazioni sulle agevolazioni tariffarie disponibili e capire se come lavoratore straniero temporaneo pu\`o accedervi. \\[4pt]

    \textbf{Sequenza:}
    \begin{enumerate}
        \item David digita ``agevolazioni'' nella barra di ricerca interna del sito.
        \item I risultati non sono chiari -- titoli come ``Carta Agevolata'', ``ISEE'', ``DSU'' non significano nulla per lui.
        \item Prova a navigare nel menu ``Abbonamenti'' $\rightarrow$ ``Agevolazioni''. Arriva a una pagina densa di testo italiano con riferimenti a leggi regionali, requisiti ISEE e documentazione.
        \item Cerca un pulsante ``English'' o una bandierina per cambiare lingua -- non c'\`e.
        \item Abbandona la pagina e chiede direttamente a un collega italiano. Il collega gli spiega che le agevolazioni sono solo per residenti con ISEE basso e che lui probabilmente non ne ha diritto.
    \end{enumerate}

    \textbf{Success criteria:}
    \begin{itemize}
        \item David trova le informazioni sulle agevolazioni in meno di 10 minuti.
        \item David capisce se ha diritto o meno senza dover chiedere a nessuno.
        \item La pagina spiega i requisiti in linguaggio semplice, con termini tradotti o spiegati.
        \item Se non ha diritto, il sito gli mostra comunque le opzioni standard disponibili (non un ``vicolo cieco'' informativo).
    \end{itemize}

    \textbf{Stima dei tempi:} 15 minuti (stato attuale), 5 minuti (con contenuti semplificati e supporto multilingua).
    }
\end{tcolorbox}

\newpage

% ================================================================
% IMPATTO DEL PERSONA SUL DESIGN
% ================================================================

\begin{tcolorbox}[colback=primaryColor!10,colframe=primaryColor,arc=2mm]
    \textbf{Considerazioni conclusive --- Impatto del persona sul design}
\end{tcolorbox}

\vspace{0.3cm}

\begin{tcolorbox}[colback=backgroundColor,colframe=borderColor,arc=2mm]
    {\small
    \textbf{Impatto del persona David Carter sul design:} \\[4pt]
    Il profilo di David --- competenza tecnica 4/5, competenza di dominio 2/5, competenza linguistica 2/5 --- ha guidato le seguenti decisioni di design:
    \begin{itemize}[leftmargin=*]
        \item \textbf{Supporto multilingua (R12):} Il sito deve offrire almeno un'interfaccia in inglese per i flussi transazionali principali. Non si tratta di tradurre tutto il sito, ma di garantire che acquisto, rinnovo e agevolazioni siano fruibili in inglese con terminologia semplificata.
        \item \textbf{Glossario contestuale (R13):} Termini come ``Codice Fiscale'', ``SPID'', ``ISEE'', ``Mensile Impersonale'', ``Mi Muovo'' devono avere un tooltip esplicativo al primo utilizzo o un'icona informativa accanto. Per l'utente straniero, questi termini non sono autoesplicativi.
        \item \textbf{Percorso guest senza SPID (R7):} David non ha SPID e non pu\`o ottenerlo. Il guest path non \`e un'opzione ma un requisito indispensabile per l'utente straniero temporaneo.
        \item \textbf{Codice Fiscale spiegato (R14):} Il campo ``Codice Fiscale'' deve includere un messaggio esplicativo bilingue che spiega cos'\`e, come si trova (sul tesserino plastificato, sulla Tessera Sanitaria) e perch\'e serve. Per uno straniero senza tessera sanitaria italiana, deve essere prevista un'alternativa (es. numero passaporto + data di nascita).
        \item \textbf{Conferma multilingue e istruzioni post-acquisto (R12):} La schermata di conferma deve essere bilingue (italiano + inglese) e includere istruzioni chiare su come usare il titolo di viaggio acquistato (``L'abbonamento \`e gi\`a attivo. Mostra questa schermata all'autista o al controllore'').
        \item \textbf{Architettura informativa orientata al compito (R1/R2):} David \`e abituato a siti dove l'azione principale \`e la prima cosa che vedi. La homepage deve ridurre il rumore visivo e presentare il percorso di acquisto in modo prominente, come standard UX internazionale.
    \end{itemize}
    }
\end{tcolorbox}

\vspace{0.3cm}

\begin{tcolorbox}[colback=backgroundColor,colframe=borderColor,arc=2mm]
    {\small
    \textbf{Impatto del Use Case \#1 (Acquisto primo abbonamento) sul design:} \\[4pt]
    \begin{itemize}[leftmargin=*]
        \item Il flusso di acquisto deve funzionare senza SPID e senza registrazione obbligatoria. Il guest checkout deve essere l'opzione predefinita, non nascosta.
        \item Il form di pagamento deve accettare carte internazionali e gestire il caso ``straniero senza codice fiscale italiano'' con un fallback chiaro (es. ``Se non hai un Codice Fiscale italiano, usa il tuo numero di passaporto'').
        \item La conferma finale deve includere un messaggio celebrativo bilingue e l'indicazione chiara della prossima scadenza con formato data internazionale (GG/MM/AAAA).
    \end{itemize}
    \vspace{0.2cm}
    \textbf{Impatto del Use Case \#2 (Agevolazioni) sul design:} \\[4pt]
    \begin{itemize}[leftmargin=*]
        \item La sezione agevolazioni deve essere accessibile dal menu principale e deve offrire un riassunto in linguaggio semplice con domande guida (``Hai un ISEE basso? Sei residente in Emilia-Romagna da almeno 2 anni?''). Per David, anche una versione inglese della checklist.
        \item Se l'utente non ha diritto alle agevolazioni, il sistema non deve lasciarlo in un ``vicolo cieco'' ma deve mostrare le opzioni standard disponibili con un link diretto all'acquisto.
        \item La ricerca interna deve restituire risultati ordinati per pertinenza anche per query in inglese parziale (es. ``discount'', ``student'', ``low income'').
    \end{itemize}
    }
\end{tcolorbox}

\vspace{0.5cm}

\begin{center}
    {\color{primaryColor}\Large\bfseries Persona and Use Case Card} \\[4pt]
    {\color{secondaryColor}\normalsize Phase B --- Feasibility Study} \\[4pt]
    {\color{secondaryColor}\small Questo \textit{persona} \textbf{non} \`e il protagonista (cio\`e non \`e colui per cui stiamo progettando il nuovo strumento; \`e un membro del cast di supporto che per\`o va comunque tenuto in considerazione per garantire un design inclusivo).}
\end{center}

\vspace{0.3cm}
{\color{borderColor}\hrule}
\vspace{0.4cm}

\begin{minipage}[t]{0.25\textwidth}
    \centering
    \includegraphics[width=\textwidth,height=5cm,keepaspectratio]
    {images/roberto.png}
    \vspace{0.2cm}
\end{minipage}
\hfill
\begin{minipage}[t]{0.70\textwidth}
    \begin{tcolorbox}[colback=backgroundColor,colframe=primaryColor,arc=2mm]
        \textbf{\large Roberto Rinaldi}\label{persona:roberto_rinaldi} \\[4pt]
        {\color{secondaryColor}\small\textbf{Citazione significativa:}
        \textit{``Sinceramente, mettermi al computer a leggere pagine e pagine di istruzioni mi fa passare la voglia. Preferisco andare allo sportello così fa tutto l'operatore, anche se un po' mi imbarazza dover fare mille domande se c'è molta fila.''}} \\[4pt]
        {\color{secondaryColor}\small\textbf{Segmento corrispondente:}
        [S3] Il cittadino allo sportello (Italiano, 36--64 anni, competenza digitale bassa, relazione istituzionale e occasionale con TPER, forte bisogno di assistenza fisica).}
    \end{tcolorbox}
\end{minipage}

\vspace{0.3cm}

% ================================================================
% DESCRIZIONE DEL PERSONA (Basic Facts)
% ================================================================

\begin{tcolorbox}[colback=backgroundColor,colframe=borderColor,arc=2mm]
    \textbf{Descrizione del persona} \\[4pt]
    {\small
    \textbf{Età:} 54 anni \hfill
    \textbf{Genere:} Maschio \hfill
    \textbf{Nazionalità:} Italiana
    \\[2pt]
    \textbf{Residenza:} Bologna, quartiere periferico (Borgo Panigale) \hfill
    \textbf{Stato civile:} Separato \\[2pt]
    \textbf{Professione e Reddito:} Collaboratore scolastico (personale ATA) con contratto a tempo determinato. Ha un ISEE basso dovuto al recente periodo di disoccupazione e all'affitto che paga per una piccola stanza in condivisione. \\[2pt]
    \textbf{Competenza linguistica:} Madrelingua italiano. Comprende perfettamente la lingua, ma si annoia subito di fronte al linguaggio burocratico e ai muri di testo densi tipici dei siti istituzionali, preferendo evitare lo sforzo di leggerli e interpretarli. \\[2pt]
    \textbf{Problemi fisici rilevanti:} Nessun problema motorio grave; soffre di una leggera presbiopia (usa occhiali da lettura per vedere da vicino) e di alcuni acciacchi alla schiena causati dai lavori manuali svolti in passato. \\[2pt]
    \textbf{Profilo digitale:} Possiede uno smartphone di fascia bassa che usa per telefonare, mandare messaggi su WhatsApp ai figli e guardare qualche video o notizia. Non soffre di fobie tecnologiche o ansie legate alla sicurezza informatica, ma è estremamente pigro nell'apprendimento digitale: se una procedura richiede più di due passaggi o l'applicazione mentale su moduli strutturati, si stufa e abbandona. Non ha attivato lo SPID semplicemente perché ha continuato a rimandare la trafila, trovandola una noiosa perdita di tempo. \\[2pt]
    \textbf{Utilizzo TPL:} Fino a tre mesi fa usava solo l'automobile. Arrivato a Bologna, si è trovato senza macchina e deve affidarsi completamente agli autobus per andare a scuola. Al momento paga biglietti singoli o giornalieri, sprecando molto denaro, e non possiede ancora un abbonamento. Rimanda costantemente l'acquisto per pura inerzia, non conoscendo ancora bene le linee o le regole tariffarie del servizio.
    }
\end{tcolorbox}

\vspace{0.3cm}

% ================================================================
% BACKGROUND E CARATTERISTICHE PSICOLOGICHE
% ================================================================

\begin{tcolorbox}[colback=backgroundColor,colframe=borderColor,arc=2mm]
    \textbf{Background e caratteristiche psicologiche} \\[4pt]
    {\small
    Roberto ha lavorato per trent'anni come operaio in una piccola azienda artigiana sull'Appennino emiliano, in un paese dove tutti si conoscevano e ci si spostava solo in auto. A causa della chiusura dell'azienda, ha attraversato un periodo di disoccupazione. Tre mesi fa ha finalmente ottenuto un incarico annuale come collaboratore scolastico e si è trasferito a Bologna da solo. I suoi figli lavorano in altre regioni e lui tende a non disturbarli per questioni burocratiche quotidiane che reputa noiose. \\[4pt]

    \textbf{Scarsa Conoscenza del Dominio e Inerzia:} Bologna per Roberto è una realtà ancora nuova e dispersiva. Non è abituato alle mille linee di autobus e al concetto di ``zone tariffarie''. Avendo una scarsa conoscenza del dominio dei trasporti, preferisce non sforzarsi a studiarlo online, manifestando una forte apatia e un atteggiamento passivo verso le regole di TPER. \\[4pt]

    \textbf{Pigrizia e Mancanza di Motivazione/Concentrazione:} La sua ridotta attenzione non nasce dalla paura di sbagliare o da timori legati alla sicurezza dei dati, ma da una profonda pigrizia mentale nei confronti della burocrazia digitale. Di fronte a uno schermo informativo o a un modulo online, Roberto perde subito la concentrazione; legge a salti, si distrae con facilità e decide che è più comodo far sbrigare la pratica a qualcun altro. La sua strategia per risparmiare tempo ed energie cognitive è delegare l'intero sforzo direttamente all'operatore dello sportello fisico. \\[4pt]

    \textbf{Timidezza in Situazioni Scomode:} Nonostante la preferenza per il canale fisico, Roberto è un uomo timido e introverso che teme il giudizio altrui. Quando si presenta allo sportello, l'idea di dover fare troppe domande o di rallentare le persone in coda lo mette a disagio. In queste situazioni scomode si rivela un po' impacciato: estrae fogli spiegazzati dalle tasche in modo confuso, si scusa spesso con l'impiegato e, se sente la pressione della fila alle sue spalle, tende a sbrigarsi accettando passivamente le indicazioni senza approfondire i suoi reali dubbi.
    }
\end{tcolorbox}

\vspace{0.3cm}

% ================================================================
% C&A DIAGRAM
% ================================================================

\begin{tcolorbox}[colback=backgroundColor,colframe=primaryColor,arc=2mm]
    \textbf{Grafico delle competenze e abilità (C\&A Diagram)} \\[4pt]
    {\small
    }

     \begin{center}
    \includegraphics[width=0.70\textwidth,keepaspectratio]
    {images/robertochart.png}
    \end{center}

    \vspace{0.3cm}

    \begin{center}
    \renewcommand{\arraystretch}{1.3}
    \begin{tabular}{|l|c|c|}
        \hline
        \textbf{Dimensione} & \textbf{Punteggio (1--5)} & \textbf{Indicatore} \\
        \hline
        Competenza tecnica &
        \textbf{2} &
        \cellcolor{gray!10} \makebox[50pt]{\rule{10pt}{10pt}\rule{10pt}{10pt}\color{lightgray}\rule{10pt}{10pt}\rule{10pt}{10pt}\rule{10pt}{10pt}} \\
        \hline
        Competenza di dominio &
        \textbf{1} &
        \cellcolor{gray!10} \makebox[50pt]{\rule{10pt}{10pt}\color{lightgray}\rule{10pt}{10pt}\rule{10pt}{10pt}\rule{10pt}{10pt}\rule{10pt}{10pt}} \\
        \hline
        Competenza linguistica &
        \textbf{5} &
        \cellcolor{gray!10} \makebox[50pt]{\rule{10pt}{10pt}\rule{10pt}{10pt}\rule{10pt}{10pt}\rule{10pt}{10pt}\rule{10pt}{10pt}} \\
        \hline
        Abilità fisica e visiva &
        \textbf{4} &
        \cellcolor{gray!10} \makebox[50pt]{\rule{10pt}{10pt}\rule{10pt}{10pt}\rule{10pt}{10pt}\rule{10pt}{10pt}\color{lightgray}\rule{10pt}{10pt}} \\
        \hline
        Motivazione &
        \textbf{1} &
        \cellcolor{gray!10} \makebox[50pt]{\rule{10pt}{10pt}\color{lightgray}\rule{10pt}{10pt}\rule{10pt}{10pt}\rule{10pt}{10pt}\rule{10pt}{10pt}} \\
        \hline
        Concentrazione &
        \textbf{2} &
        \cellcolor{gray!10} \makebox[50pt]{\rule{10pt}{10pt}\rule{10pt}{10pt}\color{lightgray}\rule{10pt}{10pt}\rule{10pt}{10pt}\rule{10pt}{10pt}} \\
        \hline
    \end{tabular}
    \end{center}

    \vspace{0.2cm}
    {\small \textbf{Legenda:} \rule{10pt}{10pt} = valore raggiunto \hfill \color{lightgray}\rule{10pt}{10pt}\color{black} = valore non raggiunto}

\end{tcolorbox}

\vspace{0.3cm}

% ================================================================
% DETTAGLI INTERESSANTI ADDIZIONALI
% ================================================================

\begin{tcolorbox}[colback=backgroundColor!50,colframe=borderColor,arc=2mm]
    \textbf{Dettagli interessanti addizionali} \\[4pt]
    {\small
    \begin{itemize}[leftmargin=*]
        \item Quando prova ad aprire il portale web per inerzia si ferma alla homepage: l'assenza di un percorso ultra-rapido o di indicazioni immediate lo spinge a chiudere la scheda dopo pochissimi secondi, liquidando la lettura con un pigro \textit{``ci penserò un altro giorno''}.
        \item Sapendo di avere un ISEE basso ma non avendo alcuna voglia di decifrare i requisiti sul web, si fida ciecamente del passaparola di un collega che gli ha detto di andare di persona a chiedere \textit{``l'abbonamento con sconti ISEE''}. Si è segnato questa frase esatta su un pezzetto di carta per evitare di confondersi o bloccarsi davanti all'impiegato per via della sua timidezza.
        \item Rappresenta il classico esempio di utente che rischia di presentarsi fisicamente allo sportello sprovvisto dei documenti corretti (come la DSU o l'attestazione ISEE ufficiale valida, portando invece scartoffie inutili), non per un problema di digital divide o di ansia, ma perché ha attuato un controllo preventivo superficiale e svogliato sul sito pubblico.
    \end{itemize}
    }
\end{tcolorbox}

% ================================================================
% USE CASE #1: ROBERTO RINALDI
% ================================================================
\begin{tcolorbox}[colback=primaryColor!10,colframe=primaryColor,arc=2mm]
    \textbf{Descrizione del Use Case \#1: Tentativo di richiesta abbonamento agevolato tra sportello fisico e canale online}
\end{tcolorbox}

\vspace{0.3cm}

% --- Narrative Use Case 1 ---
\begin{tcolorbox}[colback=backgroundColor,colframe=borderColor,arc=2mm]
    {\small
    \textbf{Scenario:} Dopo tre mesi di vita a Bologna, Roberto si rende conto che comprare biglietti singoli e giornalieri per andare a scuola gli sta prosciugando il portafoglio. Un collega gli ha detto che, avendo un ISEE basso, ha diritto a un abbonamento annuale molto scontato. Armato di buona volontà ma di pessima organizzazione, si reca di persona alla biglietteria centrale TPER nel suo pomeriggio libero. Trova una fila lunghissima. Quando finalmente arriva il suo turno, l'operatore gli chiede l'attestazione ISEE in corso di validità e una fototessera. Roberto entra in confusione: ha portato con sé solo il codice fiscale di plastica e una vecchia busta paga. Sentendo i sospiri spazientiti delle persone in fila dietro di lui, la sua timidezza prende il sopravvento. Diventa rosso, balbetta delle scuse frettolose all'impiegato ed esce letteralmente in fuga dallo sportello, troppo imbarazzato per chiedere ulteriori spiegazioni su come procurarsi i documenti o se esista un'alternativa su internet. Tornato nella sua stanza in affitto, sentendosi in colpa per il tempo perso, decide di aprire il sito web di TPER dal suo smartphone per cercare di capire come risolvere la questione senza dover tornare ad affrontare la fila. \\[4pt]

    \textbf{Obiettivo:} Scoprire quali documenti esatti servono per l'abbonamento agevolato e capire se sia possibile avviare la pratica online per evitare di tornare fisicamente allo sportello. \\[4pt]

    \textbf{Sequenza delle azioni:}
    \begin{enumerate}
        \item Roberto affronta il canale fisico: si reca in biglietteria, fa 40 minuti di fila e fallisce l'operazione perché si presenta sprovvisto della documentazione corretta.
        \item A causa dell'imbarazzo e della pressione sociale della coda alle sue spalle, abbandona l'interazione fisica senza farsi rilasciare un promemoria o istruzioni dall'operatore.
        \item Rientrato a casa, apre il browser sul proprio smartphone e cerca ``abbonamento sconti TPER''.
        \item Atterra sulla pagina dedicata alle agevolazioni tariffarie. Trattandosi di un madrelingua italiano, non ha alcun problema a leggere le parole, ma si scontra con una pagina chilometrica ricca di spiegazioni tecniche, riferimenti a delibere comunali, scaglioni ISEE e moduli scaricabili in PDF.
        \item Roberto inizia a leggere le prime due righe, ma la sua bassissima concentrazione e la totale pigrizia mentale verso la burocrazia prendono subito il sopravvento. Trova faticoso dover fare lo sforzo di tradurre mentalmente il linguaggio amministrativo per capire a quale categoria appartenga.
        \item Scorre velocemente la pagina col dito verso il basso sperando di trovare un pulsante magico o una soluzione facile, ma vede solo altro testo. 
        \item Con uno sbuffo di noia e apatia, decide che ``è una cosa troppo lunga e complicata da fare adesso''. Chiude la scheda del browser, poggia il telefono sul tavolo e accende la televisione, ripromettendosi di pensarci in un secondo momento.
        \item \textbf{Conseguenza a lungo termine:} La sua inerzia vince sul bisogno di risparmiare. Continuerà a procrastinare l'operazione per settimane, pagando le corse a prezzo pieno e aggravando la sua situazione economica per pura pigrizia burocratica.
    \end{enumerate}

    \textbf{Barriere previste:}
    \begin{itemize}
        \item \textbf{Mancanza di continuità omnicanale (Fisico $\rightarrow$ Digitale):} L'assenza di un supporto proattivo allo sportello (es. consegna di un volantino con i passaggi chiave o un QR code in caso di pratica respinta) lascia l'utente solo e disorientato quando prova a passare al canale web.
        \item \textbf{Muro burocratico vs. Inerzia dell'utente:} La presenza di un'architettura informativa basata su testi densi e PDF normativi è totalmente incompatibile con un utente che si distrae subito e non ha alcuna motivazione allo sforzo cognitivo.
    \end{itemize}

    \textbf{Success criteria (non raggiunti):}
    \begin{itemize}
        \item L'utente fallisce la pratica allo sportello fisico a causa della mancanza di una checklist preventiva chiara prima di uscire di casa.
        \item L'utente abbandona la ricerca informativa sul sito web per noia e sfinimento cognitivo.
        \item L'azienda fallisce nel convertire un utente occasionale in abbonato, perdendo l'opportunità di fidelizzarlo.
    \end{itemize}

    \textbf{Stima dei tempi:} 45 minuti di tempo perso e frustrazione nel canale fisico; meno di 2 minuti di navigazione passiva e svogliata sul canale digitale prima dell'abbandono definitivo.
    }
\end{tcolorbox}

% ================================================================
% CONSIDERAZIONI CONCLUSIVE --- IMPATTO SUL DESIGN
% ================================================================
% ================================================================
% CONSIDERAZIONI CONCLUSIVE --- IMPATTO SUL DESIGN (FOCUS ISTITUZIONALE)
% ================================================================
\begin{tcolorbox}[colback=primaryColor!10,colframe=primaryColor,arc=2mm]
    \textbf{Considerazioni conclusive --- Impatto del persona sul design nel contesto istituzionale}
\end{tcolorbox}

\vspace{0.3cm}

\begin{tcolorbox}[colback=backgroundColor,colframe=borderColor,arc=2mm]
    {\small
    \textbf{Impatto del persona Roberto Rinaldi sul design:} \\[4pt]
    Il profilo di Roberto --- caratterizzato da motivazione 1/5, concentrazione 2/5 e una profonda pigrizia verso i processi amministrativi --- richiede un radicale ripensamento del portale, che deve smettere di comportarsi come un sito aziendale e iniziare ad agire come un facilitatore dell'ecosistema istituzionale pubblico:
    \begin{itemize}[leftmargin=*]
        \item \textbf{De-frammentazione dell'Ecosistema Pubblico nella Home Page (R1):} Per un utente con scarsa conoscenza del dominio, TPER, CAF, INPS e Comune fanno parte di un unico, faticoso ``apparato istituzionale''. La Home Page deve abbattere questi silos informativi: anziché mostrare l'organigramma dei servizi o tariffe commerciali, deve presentare in primo piano macro-aree orientate ai diritti del cittadino.
        \item \textbf{Flessibilità di Autenticazione Istituzionale (R7):} Sebbene si tratti di un servizio pubblico che solitamente richiede SPID o CIE, l'obbligo di login immediato rappresenta una barriera insormontabile per la pigrizia di Roberto. Il design deve consentire un \textit{Guest Checkout} istituzionale avanzato, permettendo di istruire la pratica inserendo solo i dati anagrafici e rimandando la validazione formale dell'identità a un momento successivo o delegandola alla verifica dell'operatore.
        \item \textbf{Tono Istituzionale Accogliente e De-burocratizzato (R9):} Pur essendo un madrelingua italiano (5/5), Roberto rifiuta il ``burocratese'' e le intimidazioni normative (riferimenti a delibere, sanzioni, decreti regionali). Il micro-copy deve adottare un tono di voce che passa dall'essere *autoritario* (imposizione di regole) a *collaborativo* (accompagnamento del cittadino), riducendo la fatica e la noia della lettura.
        \item \textbf{Checklist Istituzionale Preventiva (R11):} Per evitare l'annoso problema dei cittadini che fanno la fila a vuoto intasando gli sportelli a causa di controlli superficiali, ogni sezione informativa deve aprirsi con una checklist visiva dei documenti richiesti dagli altri enti, agendo come un filtro preventivo chiaro ed estremamente sintetico.
    \end{itemize}
    }
\end{tcolorbox}

\vspace{0.3cm}

\begin{tcolorbox}[colback=backgroundColor,colframe=borderColor,arc=2mm]
    {\small
    \textbf{Impatto del Use Case \#1 (Richiesta abbonamento agevolato) sul design:} \\[4pt]
    Il fallimento parallelo di Roberto --- che scappa dallo sportello fisico per l'imbarazzo della fila e abbandona il sito web per noia cognitiva --- guida le seguenti scelte di redesign volte a unificare il canale offline e online dell'istituzione:
    \begin{itemize}[leftmargin=*]
        \item \textbf{Attivazione dello Sportellista in ``Trainer Mode'' (R6):} Lo sportello fisico non deve essere un punto di rifiuto della pratica, ma un momento di onboarding digitale assistito. Se Roberto si presenta sprovvisto di documenti, l'operatore attiva a sistema la *Trainer Mode*: il terminale memorizza i dati anagrafici già forniti e stampa uno ``Scontrino Guida'' personalizzato. Questo scontrino contiene un QR-Code unico e una checklist fisica dei documenti mancanti, riducendo l'imbarazzo di Roberto e parlando chiaramente al posto suo.
        \item \textbf{Ponte Omnicanale e Pratica Pre-compilata da QR-Code (R6):} Inquadrando il QR-Code dello scontrino da casa, Roberto non atterra sulla homepage generica, ma viene catapultato all'interno di un flusso web protetto dove la sua richiesta di abbonamento agevolato è già parzialmente istruita dall'operatore dello sportello.
        \item \textbf{Sostituzione dei bandi pubblici con flussi interattivi a imbuto (R2 + R8):} È vietato obbligare l'utente a scaricare e interpretare PDF di bandi o delibere regionali per comprendere i requisiti ISEE. Il portale istituzionale deve sostituire la saggistica normativa con un wizard interattivo (questionario a risposte chiuse e passaggi lineari) che esegue il calcolo dei requisiti al posto dell'utente, minimizzando il carico cognitivo (Concentrazione 2/5).
    \end{itemize}
    }
\end{tcolorbox}


\end{document}